\chapter{The seven faces of the collision residue}
\label{ch:holographic-datum-master}

\begin{abstract}
A single genus-zero binary collision residue
$r_\cA(z)=\mathrm{Res}^{\mathrm{coll}}_{0,2}(\Theta_\cA)$ controls the
comparison surfaces on which seven constructions are simultaneously
defined. It is the universal twisting morphism of the bar-cobar
adjunction (Theorem~\ref{thm:bar-cobar-adjunction}); the dg-shifted
Yangian spectral $r$-matrix of the line-operator package of
Dimofte--Niu--Py~\cite{DNP25}; the leading coefficient of the
classical $\lambda$-bracket of the Khan--Zeng Poisson vertex
algebra~\cite{KhanZeng25}; the symbol of the commuting differential
operators on the sphere of Gaiotto--Zeng~\cite{GZ26}; Drinfeld's
rational spectral kernel~\cite{Drinfeld85}; the spectral kernel of
the Sklyanin Poisson bracket of Semenov-Tian-Shansky~\cite{STS83};
and the kernel of the Gaudin Hamiltonians of
Feigin--Frenkel--Reshetikhin~\cite{FFR94}. Faces \textup{(F1)},
\textup{(F3)}, and \textup{(F4)} are available on the
modular-Koszul surface named below; the DNP face requires a
line-operator realization; the Drinfeld, Sklyanin, and
Gaudin faces are the affine Kac--Moody specialization with the
trace-form and KZ normalizations related by an explicit Casimir
rescaling. The seven-face theorem identifies the common convolution
dg~Lie algebra in which these available faces agree. A disagreement
between two available faces is a normalization error, a failed
hypothesis, or a wrong computation.
\end{abstract}

\section{Thesis: six projections of a single boundary-bulk algebra}
\label{sec:hdm-thesis}

On the comparison surfaces specified in this chapter, seven independent
mathematical frameworks produce the same binary object. The collision
residue~$r_\cA(z)$ appears as a
twisting morphism, a spectral $R$-matrix, a $\lambda$-bracket
coefficient, a differential-operator symbol, a Yangian structure
constant, a Sklyanin bracket kernel, and a Gaudin Hamiltonian
kernel. Their coincidence is not a fortunate accident of matching
normalizations; it is the binary shadow of the boundary-bulk
algebra~$U_\cA$ in the ambient where the relevant comparison has been
constructed. The seven faces are the images of that algebra under
seven projection functors. The seven-face master theorem
(Theorem~\ref{thm:hdm-seven-way-master}) therefore identifies the
conditional surface on which the seven constructions descend from the
same boundary-bulk algebra.

The theorem sits inside a
larger claim. Fix a standard chiral algebra~$\cA$ in the standard
landscape \textup{(}Heisenberg, affine Kac--Moody, $\beta\gamma/bc$,
Virasoro, or principal $\mathcal{W}_N$\textup{)}, or more generally a
modular Koszul chiral algebra together with the extra hypotheses named
when the line-side comparison is invoked. The six-slot datum
\[
\mathcal{H}(T) \;=\;
\bigl(\cA,\; \cA^!,\; \cC,\; r_\cA(z),\; \Theta_\cA,\; \nabla^{\mathrm{hol}}\bigr),
\]
formally introduced in
Definition~\ref{def:holographic-modular-koszul-datum}, is not a
six-tuple of independent invariants. It is the six-slot projection of
a single boundary-bulk algebra, namely the Drinfeld double
$U_\cA := \cA \bowtie \cA^{!}_\infty$ assembled from the boundary
chiral algebra~$\cA$ and its Verdier-dual factorization companion.
Each slot of $\mathcal{H}(T)$ captures one aspect of the
boundary-bulk reconstruction, and the six aspects are interlocking
projections of the one algebra~$U_\cA$:
the boundary object is~$\cA$, the bulk object is the pinned derived
centre~$\cC$, the comparison map is the restriction
$\rho_{\partial}\colon \cC\to\cA$ induced by
$\mathrm{Ass}^{\mathrm{ch}}\hookrightarrow\mathrm{SC}^{\mathrm{ch,top}}$,
and the witnesses are the same triple
$(r_\cA(z),\Theta_\cA,\nabla^{\mathrm{hol}})$. Off the
modular-Koszul / boundary-linear / principal DS surfaces named below,
the strict edge $\cA^{!}_\infty\to\cA^!$ and the global-triangle
comparison to the line model are the active obstructions.
\begin{enumerate}[label=\textup{(\arabic*)}]
\item \emph{Boundary slot.} $\cA$ is the boundary chiral algebra,
 the primary factor of~$U_\cA$.
\item \emph{Verdier/Koszul slot.} $\cA^!$ means the strict chiral
 Koszul dual obtained from the Verdier-dual factorization companion
~$\cA^{!}_\infty$ when the rectification map
 $\cA^{!}_\infty \to \cA^!$ is proved. We never identify
 $\cA^!$ with $\Omega(\barB(\cA))$: $\Omega(\barB(\cA))\simeq \cA$ is
 bar--cobar inversion, not Koszul duality, and
 $(H^*(\barBch(\cA)))^\vee$ is only the linear shadow of the Verdier
 construction.
\item \emph{Bulk/line slot.} We pin
 \[
 \cC \;:=\; \cZ^{\mathrm{der}}_{\mathrm{ch}}(\cA)
 \;:=\; \mathcal{C}^{\bullet}_{\mathrm{ch}}(\cA,\cA)
 \]
 on the formal disk, and for the global variant we use the same symbol
 for the corresponding object on $X_{\mathrm{dR}}$. Every later
 appearance of~$\cC$ in this chapter refers to this pin. The line-side
 presentation $\cC_{\mathrm{line}} \simeq
 \cA^!_{\mathrm{line}}\text{-}\mathsf{mod}$ is a second model for the
 same slot only on the boundary-linear / Drinfeld--Sokolov lanes where
 the compact-generator comparison is proved.
\item \emph{$R$-matrix slot.} We pin the fourth slot to the
 family-normalized binary collision residue
 \[
 r_\cA(z) \;:=\; \operatorname{Res}^{\mathrm{coll}}_{0,2}(\Theta_\cA).
 \]
 In the standard families this specializes to
 \[
 r^{\mathrm{Heis}}_k(z)=\frac{k}{z}\,J\otimes J,\qquad
 r^{KM,\fg}_{\mathrm{coll},k}(z)=\frac{k\,\Omega_{\mathrm{tr}}}{z},
 \qquad
 r^{KM,\fg}_{\mathrm{KZ},k}(z)=\frac{\Omega_{\mathrm{KZ}}}{(k+h^\vee)z},
 \]
 and
 \[
 r^{\mathrm{Vir}}_c(z)=\frac{c/2}{z^3}+\frac{2T}{z}.
 \]
 An unspecialized bare $r$-symbol will therefore not be used below.
\item \emph{Maurer--Cartan slot.} In the chain-level lane we pin
 $\Theta_\cA$ to the degree-$1$ Maurer--Cartan element of the ordered
 convolution dGLA
 \[
 \mathfrak{g}^{E_1}_\cA
 \;:=\;
 \operatorname{Conv}\bigl(B^{\mathrm{ord}}(\cA),\cA\bigr),
 \qquad
 D\cdot\Theta_\cA + \tfrac{1}{2}[\Theta_\cA,\Theta_\cA]=0.
 \]
 Its averaged image in the modular convolution dGLA is denoted by the
 same symbol. The gauge class is presentation-independent on the
 generic cyclic-rigid locus, so the slot is universal up to gauge.
\item \emph{Holomorphic-connection slot.} We pin
 \[
 \nabla^{\mathrm{hol}}_{g,n}
 \;:=\;
 d-\operatorname{Sh}_{g,n}(\Theta_\cA)
 \]
 as the flat holomorphic connection on the ordered bar bundle over
 $\mathrm{Conf}_n(X)$ induced by the universal MC element. On the
 Heisenberg and affine Kac--Moody lanes this specializes to the KZ
 connections built from $r^{\mathrm{Heis}}_k(z)$ and
 $r^{KM,\fg}_{\mathrm{KZ},k}(z)$; on the Virasoro lane it specializes to the
 $\mathrm{SL}_2$-type conformal connection built from
 $r^{\mathrm{Vir}}_c(z)$. Every later appearance of
 $\nabla^{\mathrm{hol}}$ refers to this pin.
\end{enumerate}

In this chapter, the words \emph{holographic}, \emph{boundary-bulk},
and \emph{open/closed} always refer to the chain-level comparison
between the boundary object~$\cA$ and the chiral derived centre
$\cC = \cZ^{\mathrm{der}}_{\mathrm{ch}}(\cA)$ (T18(i), the
chain-level \emph{bulk object} below) along the
restriction map $\rho_{\partial}\colon \cC \to \cA$ and the operadic
inclusion $\mathrm{Ass}^{\mathrm{ch}} \hookrightarrow
\mathrm{SC}^{\mathrm{ch,top}}$, witnessed by the same trio
$(r_\cA(z),\Theta_\cA,\nabla^{\mathrm{hol}})$. In this
chapter \emph{bulk} always means T18(i); the $3$d HT bulk
(T18(iii)) appears only as the Brown--Henneaux partition-function
partner recalled below. The proved scope is the
modular-Koszul, boundary-linear, and principal DS comparison surfaces
named below; off those surfaces the obstruction is explicit: either the
strict edge $\cA^{!}_\infty \to \cA^!$ fails, or the global triangle
relating the bulk slot to the line model does not close.

Under this reading, the three structural constants of the thesis have
the following content. The chiral derived centre
$\cZ^{\mathrm{der}}_{\mathrm{ch}}(\cA)$ is the bulk algebra, computed
from~$\cA$ by chiral Hochschild cochains. The modular
characteristic~$\kappa(\cA)$ is the degree-$2$ scalar component of
$\Theta_\cA$ and the Brown--Henneaux coefficient of the bulk partition
function
(Theorem~\ref{thm:genus-universality}). The shadow obstruction tower
$\{S_r(\cA)\}_{r \geq 2}$ is the correction hierarchy, the
sequence of finite-order projections of~$\Theta_\cA$ whose termination
controls the depth of boundary-bulk reconstruction; and
$U_\cA = \cA \bowtie \cA^{!}_\infty$ is the universal boundary-bulk
algebra organizing the proved comparison maps from the boundary VOA to
the bulk slot $\cC=\cZ^{\mathrm{der}}_{\mathrm{ch}}(\cA)$ along
$\rho_{\partial}\colon \cC\to\cA$ and
$\mathrm{Ass}^{\mathrm{ch}}\hookrightarrow\mathrm{SC}^{\mathrm{ch,top}}$,
witnessed by $(r_\cA(z),\Theta_\cA,\nabla^{\mathrm{hol}})$ on the
named modular-Koszul / boundary-linear / principal DS surfaces
(introduction, Remark~\ref{rem:five-facets-boundary-bulk}).

The collision residue~$r_\cA(z)$, which is the subject of this chapter,
is slot~(4) of this six-tuple: the spectral $R$-matrix of the
quasitriangular structure of~$U_\cA$. The seven-face master theorem
is the verification checklist for slot~(4): the available independent
computations of~$r_\cA(z)$ across the seven frameworks must agree on
their common comparison domain, and when they do, the $R$-matrix slot
of~$U_\cA$ is verified on that domain.
The chapter therefore has two readings. Narrowly, it proves the
unification of the available computations of one spectral $r$-matrix.
Broadly, it verifies one of the six slots of the boundary-bulk
reconstructor~$U_\cA$, and in doing so instantiates the thesis that
$\mathcal{H}(T)$ is the image of a single algebra under six
projections. In that broad reading the comparison still runs between
boundary object~$\cA$ and bulk slot~$\cC$ via
$\rho_{\partial}$ and the Swiss-cheese inclusion, with witness triple
$(r_\cA(z),\Theta_\cA,\nabla^{\mathrm{hol}})$; off the named surfaces,
the strict edge and global-triangle obstructions keep the remaining
five slots from assembling into a fully strict theorem. The remaining
five slots admit analogous seven-face
verifications (not performed in this chapter); the structural
statement that they all derive from~$U_\cA$ is
Proposition~\ref{prop:hdm-u-a-recoverability} below.

\section{The collision residue as universal object}
\label{sec:hdm-collision-residue}

The ordered chain-level starting point is the degree-$1$
Maurer--Cartan element of the ordered convolution dGLA
\[
\Theta_\cA \in \mathfrak{g}^{E_1,1}_\cA
= \operatorname{Conv}\bigl(B^{\mathrm{ord}}(\cA),\cA\bigr)^1,
\qquad
D\cdot\Theta_\cA + \tfrac{1}{2}[\Theta_\cA, \Theta_\cA] = 0.
\]
Its averaged modular image is the bar-intrinsic element
\[
\Theta_\cA = D_\cA - d_0 \in
\mathrm{MC}(\Defcyc(\cA) \,\widehat{\otimes}\, \mathcal{G}_{\mathrm{mod}})$
,
\]
which is the central object of this monograph
\textup{(}Theorem~\textup{\ref{thm:mc2-bar-intrinsic}}\textup{)}. It is
the unique element of
the modular cyclic deformation complex satisfying
$D_\cA \Theta_\cA + \tfrac{1}{2}[\Theta_\cA, \Theta_\cA] = 0$
that lifts the chiral bracket to all degrees and all genera. The
\emph{collision residue} is the projection of $\Theta_\cA$ onto the
codimension-one boundary stratum of the genus-zero, two-point
moduli space:
\begin{equation}\label{eq:hdm-collision-residue-defn}
r_\cA(z)
\;:=\;
\operatorname{Res}^{\mathrm{coll}}_{0,2}(\Theta_\cA)
\;=\;
\Theta_\cA\bigm|_{\partial\overline{\mathcal{M}}_{0,2}^{\mathrm{coll}}}.
\end{equation}
The notation $r_\cA(z)$ records the residual spectral parameter that
survives the collision: when two marked points~$z_1, z_2$ approach
each other, the bar propagator $\eta(z_1,z_2) = d\log(z_1 - z_2)$
acquires a residue along the diagonal, leaving a one-parameter
family $r_\cA(z) \in \mathrm{End}_\cA(2)[\![z^{-1}]\!]$ indexed by
the relative coordinate $z = z_1 - z_2$.

The existence of $r_\cA(z)$ follows from the bar-intrinsic
construction. $\Theta_\cA$ exists by
Theorem~\ref{thm:mc2-bar-intrinsic}: the universal MC element is
$D_\cA - d_0$, and its restriction to any boundary stratum is
well-defined because $D_\cA$ is a graded morphism on the modular
cyclic complex. Uniqueness up to gauge equivalence is the content
of generic cyclic rigidity
(Theorem~\ref{thm:cyclic-rigidity-generic}): the cyclic deformation
complex has one-dimensional second cohomology in degree~$2$ on the
generic locus. On that locus the gauge class is therefore
presentation-independent. The collision residue is
therefore a canonical invariant of~$\cA$.

\begin{theorem}[Seven-face master theorem; \ClaimStatusConditional]
\label{thm:seven-faces-master}
\index{collision residue!seven-face master theorem|textbf}
\index{seven faces!collision residue|textbf}
Let $\cA$ be a modular Koszul chiral algebra on a smooth projective
curve~$X$. The collision residue
$r_\cA(z) = \mathrm{Res}^{\mathrm{coll}}_{0,2}(\Theta_\cA)$ admits
the following realizations on their stated comparison surfaces:
\begin{enumerate}[label=\textup{(F\arabic*)}]
\item the universal twisting morphism
	 $\pi_\cA \in \mathrm{Tw}(\barBch(\cA), \cA)$ of the bar-cobar
 adjunction;
\item the spectral $r$-matrix of the dg-shifted Yangian arising as
 the line-operator algebra of Dimofte--Niu--Py~\cite{DNP25};
\item the leading $\lambda$-bracket coefficient of the classical
 Poisson vertex algebra of Khan--Zeng~\cite{KhanZeng25} (genus~$0$ only);
\item the symbol of the commuting sphere differential operators of
 Gaiotto--Zeng~\cite{GZ26};
\item the spectral $r$-matrix of the Yangian quasitriangular
 structure of Drinfeld~\cite{Drinfeld85};
\item the classical $r$-matrix defining the Sklyanin Poisson bracket
 of Semenov-Tian-Shansky~\cite{STS83};
\item the kernel of the Gaudin Hamiltonians of
	 Feigin--Frenkel--Reshetikhin~\cite{FFR94}.
\end{enumerate}
Faces \textup{(F1)}, \textup{(F3)}, and \textup{(F4)} are available
for every modular Koszul~$\cA$ in the genus-zero comparison
ambient. Face \textup{(F2)} is available when~$\cA$ is realised as
the closed colour of the Dimofte--Niu--Py line-operator package.
Faces \textup{(F5)}--\textup{(F7)} are the affine
Kac--Moody specialization. On the common domain of any two available
faces, their realizations are equal as elements of
$\mathrm{End}_\cA(2)[\![z^{-1}]\!]$, after the explicit
normalization conventions of \S\ref{sec:hdm-master}. On the affine
lane the trace-form Drinfeld kernel
$k\,\Omega_{\mathrm{tr}}/z$ and the KZ collision residue
$\Omega/((k+h^\vee)z)$ are not literal rational-function identities;
they are representatives of the same binary comparison after the
finite Casimir rescaling
$k\,\Omega_{\mathrm{tr}}\mapsto\Omega/(k+h^\vee)$.
\end{theorem}

\noindent The proof of Theorem~\ref{thm:seven-faces-master} is
distributed across \S\ref{sec:hdm-face-1}--\S\ref{sec:hdm-face-7}:
each section establishes one face on its stated comparison surface and
identifies it with the collision residue there, citing the formal
cross-volume theorem that proves the identification.
\S\ref{sec:hdm-master} assembles the available identifications into a
common-domain equation.

\section{Face 1: the bar-cobar twisting morphism}
\label{sec:hdm-face-1}

The first incarnation of $r_\cA(z)$ is the one closest to the
algebraic engine of this monograph. The bar-cobar adjunction
$\Omega^{\mathrm{ch}} \dashv B^{\mathrm{ch}}$ (cobar left, bar right) between chiral
algebras and conilpotent chiral coalgebras (Theorem~\ref{thm:bar-cobar-adjunction}) has unit and
counit that give a canonical \emph{universal twisting morphism}
\begin{equation}\label{eq:hdm-twisting-morphism}
\pi_\cA \;\colon\; \barBch(\cA) \;\longrightarrow\; \cA,
\qquad
\pi_\cA \in \mathrm{Tw}(\barBch(\cA), \cA),
\end{equation}
characterized by the property that $\pi_\cA$ corresponds, under the
three bijections of Corollary~\ref{cor:three-bijections}, to the
identity element of $\Hom_{\mathrm{cdgca}}(\barBch(\cA), \barBch(\cA))$
and the identity element of
$\Hom_{\mathrm{dga}}(\Omega(\barBch(\cA)), \cA)$. Concretely,
$\pi_\cA$ extracts the chiral bracket against the logarithmic
propagator: for $[s^{-1}\bar{a}]$ a generating cycle of $\barBch(\cA)$,
\begin{equation}\label{eq:hdm-twisting-propagator}
\pi_\cA([s^{-1}\bar{a}])(z)
\;=\;
\operatorname{Res}_{w = z}
\bar{a}(w) \cdot d\log(z - w),
\end{equation}
which is precisely the propagator-extraction formula of
Proposition~\ref{prop:twisting-morphism-propagator}.

\begin{theorem}[Face~1: bar-cobar twisting; \ClaimStatusConditional]
\label{thm:hdm-face-1}
The collision residue equals the universal twisting morphism:
\begin{equation}\label{eq:hdm-face-1}
r_\cA(z)
\;=\;
\operatorname{Res}^{\mathrm{coll}}_{0,2}(\Theta_\cA)
\;=\;
\pi_\cA
\;\in\;
\mathrm{Tw}(\barBch(\cA), \cA).
\end{equation}
After passage to Koszul coalgebra coordinates
$\cA^{\mathrm{i}}:=H^*(\barBch(\cA))$, and then to strict Verdier-dual
algebra coordinates when Upgrade~\textup{(U2)} holds,
$\cA^!=(\cA^{\mathrm{i}})^\vee$, the morphism~$\pi_\cA$ becomes an
element of $\cA^! \otimes \cA^![\![z^{-1}]\!]$ that satisfies the
classical Yang--Baxter equation.
\end{theorem}

\begin{proof}
The identification $r_\cA(z) = \pi_\cA$ is
Theorem~\ref{thm:collision-residue-twisting}: the bar-intrinsic
MC element restricted to the genus-zero, two-point collision stratum
recovers the universal twisting morphism, with the geometric
realization~\eqref{eq:hdm-twisting-propagator}. The Koszul-dual
passage and the Yang--Baxter equation are
Theorem~\ref{thm:collision-residue-twisting}(ii) and
Theorem~\ref{thm:collision-depth-2-ybe}, respectively.
\end{proof}

This is the structurally cleanest face: it makes no reference to
spectral parameters, Yangians, or sigma models. The collision
residue is the morphism through which the bar coalgebra acts on the
algebra; this is the content of bar-cobar duality at the level of
twisting morphisms.

\section{Face 2: the line-operator $r$-matrix of Dimofte--Niu--Py}
\label{sec:hdm-face-2}

The 2025 paper of Dimofte, Niu, and Py~\cite{DNP25} studies line
operators in three-dimensional holomorphic-topological gauge theory
and constructs an associative meromorphic tensor category of line
operators whose endomorphism algebra is a dg-shifted Yangian. The
$A_\infty$-Yang--Baxter element of the DNP framework is a spectral
$r$-matrix
\begin{equation}\label{eq:hdm-dnp-rmatrix}
r^{\mathrm{DNP}}(z)
\;\in\;
\cA^! \otimes \cA^! [\![z^{-1}]\!],
\end{equation}
controlling the meromorphic braiding of line operators in the
$z$-plane transverse to the holomorphic direction. DNP prove that
$r^{\mathrm{DNP}}(z)$ satisfies a one-loop exact $A_\infty$-YBE, and
that the meromorphic tensor category is non-renormalized beyond
one loop. The framework is geometric: $r^{\mathrm{DNP}}(z)$ is
extracted from the BRST-closed two-point function of line operators
in the 3d HT theory whose closed colour is~$\cA$.

\begin{theorem}[Face~2: DNP line operators; \ClaimStatusConditional\ (identification with collision residue); \ClaimStatusProvedElsewhere\ (DNP line-operator package: Dimofte--Niu--Py 2025; Vol.~II cross-volume matching)]
\label{thm:hdm-face-2}
For $\cA$ a chirally Koszul logarithmic $\mathrm{SC}^{\mathrm{ch,top}}$-algebra
realised as the closed colour of the DNP 3d HT theory, the DNP
spectral $r$-matrix equals the collision residue:
\begin{equation}\label{eq:hdm-face-2}
r^{\mathrm{DNP}}(z) \;=\; r_\cA(z).
\end{equation}
The non-renormalization theorem of~\textup{\cite{DNP25}} corresponds
to the $E_2$-collapse of the bar spectral sequence for chirally
Koszul~$\cA$, and the $A_\infty$-YBE of $r^{\mathrm{DNP}}(z)$ is the
infinitesimal braid relation derived from collision depth~$2$ in the
modular MC equation.
\end{theorem}

\begin{proof}
The cross-volume identification is established in
Vol.~II, where the DNP line-operator package is matched against the
bar-cobar twisting package face by face. Part~(i) of that
identification proves the monoidal equivalence between the DNP
meromorphic tensor product and the $R$-twisted tensor product
internal to $\barBch(\cA)$-comodules; part~(ii) proves
$r^{\mathrm{DNP}}(z) = \pi_\cA = r_\cA(z)$ by matching the
Feynman-graph extraction of the DNP $r$-matrix with the bar-complex
extraction of the collision residue; part~(iii) proves that the
one-loop exactness and non-renormalization claims of~\cite{DNP25}
correspond to the $E_2$-collapse of the bar spectral sequence,
which is the chiral Koszulness condition (item~(ii) of
Theorem~\ref{thm:koszul-equivalences-meta}). The matching of
spectral parameters is the content of the dg-shifted Yangian
recognition theorem~\cite[Thm.~5.5]{DNP25} together with
Face~1. The result is documented in Vol.~II, line-operators
chapter.
\end{proof}

The chirally Koszul qualifier is sharp: off that locus the bar
spectral sequence need not collapse at~$E_2$, the Verdier-dual
companion need not rectify to a strict~$\cA^!$, and the line-operator
kernel need not admit a proved identification with the pinned bulk
slot~$\cC$.

Face~2 is the open/closed face of the collision residue. The boundary
object is~$\cA$, the bulk object is the DNP holomorphic-topological
line-operator package, and the comparison map is the chain-level
identification of the restriction map
$\rho_{\partial}\colon \cZ^{\mathrm{der}}_{\mathrm{ch}}(\cA)\to \cA$
with the operadic inclusion
$\mathrm{Ass}^{\mathrm{ch}}\hookrightarrow\mathrm{SC}^{\mathrm{ch,top}}$
transported through the Vol.~II equivalence between line operators and
$\barBch(\cA)$-comodules. The witnesses are the same three pieces of
homotopy data as in the chapter-top pin:
$r^{\mathrm{DNP}}(z)=r_\cA(z)$, the universal MC element~$\Theta_\cA$,
and the genus-zero shadow connection~$\nabla^{\mathrm{hol}}_{0,2}$.
Its exact scope is the chirally Koszul DNP-realized closed-colour
lane of Theorem~\ref{thm:hdm-face-2}; outside that lane the comparison
with the pinned bulk slot~$\cC$ remains conjectural because the bar
spectral sequence need not collapse, the strict edge
$\cA^{!}_\infty\to\cA^!$ need not rectify, and the line-operator kernel
need not identify with~$\cC$.

\section{Face 3: the classical $\lambda$-bracket of Khan--Zeng \textup{(}genus~$0$ only\textup{)}}
\label{sec:hdm-face-3}

The 2025 paper of Khan and Zeng~\cite{KhanZeng25} formulates a
3d Poisson sigma model whose target is a Poisson vertex algebra
$(\cP, \{\cdot_\lambda\cdot\})$ in the sense of Barakat--De Sole--Kac,
and proves that the gauge invariance of the sigma-model action is
equivalent to the $\lambda$-Jacobi identity for the Poisson vertex
bracket. The classical $\lambda$-bracket is a Laurent series in the
formal variable~$\lambda$:
\begin{equation}\label{eq:hdm-kz-lambda-bracket}
\{a_\lambda b\}
\;=\;
\sum_{n \geq 0} \frac{\lambda^n}{n!}\, a_{(n)}b
\;\in\;
\cP[\lambda],
\end{equation}
where $a_{(n)}b$ are the OPE-mode coefficients of the underlying
chiral algebra (with the divided-power conversion ). The
\emph{leading coefficient} $a_{(0)}b$ is the binary $\lambda^0$-bracket;
the higher coefficients $a_{(n)}b$ for $n \geq 1$ are encoded in the
$\lambda$-dependence.

\begin{theorem}[Face~3: Khan--Zeng PVA \textup{(}genus~$0$ only\textup{)}; \ClaimStatusProvedHere]
\label{thm:hdm-face-3}
Let $\cP = \mathrm{gr}^{\mathrm{Li}}(\cA)$ be the classical Poisson
vertex algebra obtained from the Li filtration of~$\cA$. Then the
collision residue equals the generating function of the
$\lambda$-bracket:
\begin{equation}\label{eq:hdm-face-3}
r_\cA(z)
\;=\;
\sum_{n \geq 0} \frac{1}{z^{n+1}}\, \mathrm{cl}(a_{(n)}b)
\;\in\;
\cP \otimes \cP [\![z^{-1}]\!],
\end{equation}
where $\mathrm{cl}\colon \cA \to \cP$ is the classical limit and
the right-hand side is the formal Laplace transform of the
$\lambda$-bracket~\eqref{eq:hdm-kz-lambda-bracket}. The
$\lambda$-Jacobi identity for $\{\cdot_\lambda\cdot\}$ corresponds
to the classical Yang--Baxter equation for $r_\cA(z)$.
\end{theorem}

\begin{proof}
The classical limit of the chiral OPE is the Poisson vertex algebra
structure on the associated graded $\mathrm{gr}^{\mathrm{Li}}(\cA)$
under the Li filtration, by the standard chiral-to-classical
construction. The collision residue
$r_\cA(z) = \pi_\cA(z)$ extracts the OPE pole expansion against the
logarithmic propagator $d\log(z-w)$, which is the Laplace transform
of the $\lambda$-bracket up to the divided-power normalization. The
identification of the classical $\lambda$-bracket leading coefficient
$a_{(0)}b$ with $\mathrm{Res}_{z=0}(z\, r_\cA(z))$ is direct. The
$\lambda$-Jacobi identity becomes the classical YBE for $r_\cA(z)$
under the Laplace transform; the gauge invariance of the
Khan--Zeng sigma model is therefore equivalent to the IBR of
Theorem~\ref{thm:collision-depth-2-ybe} after passage through
Face~1. This proves the genus-zero formal-disk comparison asserted
here. The higher-genus deformation-quantization comparison with the
bar-intrinsic MC element is a separate problem and is not used in
this face.
\end{proof}

Face~3 \textup{(}genus~$0$ only\textup{)} is the classical-Poisson incarnation of the collision
residue. Where Face~1 produces a morphism in the dg category of
twisting morphisms and Face~2 produces a quantum spectral $r$-matrix,
Face~3 produces the classical limit: the binary leading coefficient
of the classical $\lambda$-bracket. The Khan--Zeng paper provides
the geometric realization of this classical limit as a 3d Poisson
sigma model. The higher-genus extension is not merely unwritten: the
missing input is a proved comparison between higher-genus chiral
Hochschild classes and the Khan--Zeng sigma-model moduli problem, so
the face is presently pinned only on the genus-zero formal-disk/sphere
sector.

\section{Face 4: the commuting sphere Hamiltonians of Gaiotto--Zeng}
\label{sec:hdm-face-4}

The 2026 paper of Gaiotto and Zeng~\cite{GZ26} studies interface
minimal-model holography and constructs, for each chiral algebra~$\cA$
in their framework, a family of commuting differential operators
$\{H_i^{\mathrm{GZ}}\}_{i=1}^n$ on the configuration space
$\mathrm{Conf}_n(\mathbb{P}^1)$ of $n$ marked points on the sphere.
The operators are derived from the chiral OPE data via a contour
manipulation in the holomorphic-topological bulk theory, pulled back to
the boundary sphere by the restriction map
$\rho_{\partial}\colon \cZ^{\mathrm{der}}_{\mathrm{ch}}(\cA)\to \cA$ and
the binary part of the shadow connection~$\nabla^{\mathrm{hol}}_{0,n}$.
Their commutativity
$[H_i^{\mathrm{GZ}}, H_j^{\mathrm{GZ}}] = 0$ is a non-trivial
consequence of the OPE structure. The order of $H_i^{\mathrm{GZ}}$
as a differential operator on $\mathrm{Conf}_n(\mathbb{P}^1)$
depends on the collision depth $k_{\max}(\cA)$ of $\cA$ via the
operator-order trichotomy of Chapter~\ref{ch:three-invariants}.

\begin{theorem}[Face~4: GZ26 commuting Hamiltonians; \ClaimStatusConditional]
\label{thm:hdm-face-4}
The Gaiotto--Zeng commuting Hamiltonians on the sphere are
generated by the collision residue: each $H_i^{\mathrm{GZ}}$ is the
contraction of $r_\cA(z_i - z_j)$ against a chosen test state at
the $j$-th marked point, summed over $j \neq i$:
\begin{equation}\label{eq:hdm-face-4}
H_i^{\mathrm{GZ}}
\;=\;
\sum_{j \neq i}
\bigl[ r_\cA(z_i - z_j) \bigr]_{i,j}
\;\in\;
\mathrm{Diff}^{k_{\max}(\cA)-1}\bigl(\mathrm{Conf}_n(\mathbb{P}^1)\bigr).
\end{equation}
The commutativity $[H_i^{\mathrm{GZ}}, H_j^{\mathrm{GZ}}] = 0$ is
the flatness of the modular shadow connection $\nabla^{\mathrm{hol}}_{0,n}$
on the sphere. In particular, the GZ26 sphere reconstruction
theorem is a consequence of the sphere flatness theorem
\textup{(}Theorem~\textup{\ref{thm:sphere-reconstruction})}.
\end{theorem}

The sphere qualifier is essential: beyond genus~$0$ one would have to
identify the higher-genus shadow connection
$\nabla^{\mathrm{hol}}_{g,n}$ with the GZ contour construction on
$\mathrm{Conf}_n(X)$, and no such physical comparison is presently
proved.

\begin{proof}
The GZ26 commuting differential operators are defined by a contour
integral that, after the bar-complex translation of
Theorem~\ref{thm:collision-residue-twisting}, becomes the chiral
bracket against the logarithmic propagator $d\log(z_i - z_j)$.
This is the binary collision residue of $\Theta_\cA$ acting between
the $i$-th and $j$-th marked points: see~\eqref{eq:hdm-twisting-propagator}.
The summation over $j \neq i$ produces the Hamiltonian on the
$n$-point sphere because the modular shadow connection
$\nabla^{\mathrm{hol}}_{0,n} = d - \mathrm{Sh}_{0,n}(\Theta_\cA)$
is built precisely from these binary collision residues
(Theorem~\ref{thm:sphere-reconstruction}(i)). The commutativity
$[H_i^{\mathrm{GZ}}, H_j^{\mathrm{GZ}}] = 0$ is the flatness
$(\nabla^{\mathrm{hol}}_{0,n})^2 = 0$, which follows from the
infinitesimal braid relation
(Theorem~\ref{thm:collision-depth-2-ybe}) and the Arnold relation
on $\overline{\cM}_{0,n}$. The order $k_{\max}(\cA) - 1$ of
$H_i^{\mathrm{GZ}}$ as a differential operator is the operator-order
trichotomy of Chapter~\ref{ch:three-invariants}, controlled by the
collision depth.
\end{proof}

Face~4 is the second-order differential incarnation of the
collision residue: where Faces~1--3 produce algebraic objects
(morphisms, $r$-matrices, brackets), Face~4 produces commuting
differential operators on a moduli space. The order of these
operators is governed by the collision depth, which equals
$p_{\max}(\cA) - 1$: trivial for $\beta\gamma$ ($k_{\max} = 0$),
multiplication operators for affine Kac--Moody and Heisenberg
($k_{\max} = 1$), and genuine differential operators of order
$\geq 2$ for Virasoro and $\Walg_N$ ($k_{\max} \geq 3$).

\section{Face 5: the Yangian $r$-matrix of Drinfeld}
\label{sec:hdm-face-5}

The original spectral $r$-matrix is Drinfeld's~\cite{Drinfeld85},
introduced in connection with the quantum Yang--Baxter equation and
the Yangian deformation $Y_\hbar(\fg)$ of the universal enveloping
algebra of a simple Lie algebra~$\fg$. In the trace-form
normalization used in this volume, the classical $r$-matrix is
\begin{equation}\label{eq:hdm-drinfeld-rmatrix}
r^{\mathrm{Dr}}(z)
\;=\;
\frac{k\,\Omega_{\mathrm{tr}}}{z}
\;\in\;
\fg \otimes \fg [\![z^{-1}]\!],
\qquad
\Omega_{\mathrm{tr}} \;=\; \sum_a J^a \otimes J^a,
\end{equation}
where $\{J^a\}$ is an orthonormal basis of $\fg$ with respect to the
Killing form, satisfies the classical
Yang--Baxter equation
$[r^{\mathrm{Dr}}_{12}(z), r^{\mathrm{Dr}}_{13}(z+w)]
+ [r^{\mathrm{Dr}}_{12}(z), r^{\mathrm{Dr}}_{23}(w)]
+ [r^{\mathrm{Dr}}_{13}(z+w), r^{\mathrm{Dr}}_{23}(w)] = 0$. This
is the prototype of all spectral $r$-matrices and the historical
origin of the Yangian as the quantization of the Lie bialgebra
$(\fg[\![z^{-1}]\!], r^{\mathrm{Dr}})$.
The historical literature often absorbs the level into the chosen
invariant form. In the trace-form normalization used in Volume~I we
keep the level explicit as the prefactor~$k$, so
$r^{\mathrm{Dr}}(z)|_{k=0} = 0$.

\begin{theorem}[Face~5: Drinfeld $r$-matrix; \ClaimStatusConditional\ (identification with collision residue); \ClaimStatusProvedElsewhere\ (classical $r$-matrix theory: Drinfeld 1985)]
\label{thm:hdm-face-5}
For $\cA = \widehat{\fg}_k$ the affine Kac--Moody algebra at level
$k \neq -h^\vee$, the KZ-normalized collision residue is
\begin{equation}\label{eq:hdm-face-5}
r^{KM,\fg}_{\mathrm{KZ},k}(z)
\;=\;
\frac{\Omega}{(k + h^\vee)\, z}
\;\in\;
\fg \otimes \fg [\![z^{-1}]\!].
\end{equation}
On the generic trace-form locus this is matched with
Drinfeld's level-prefixed kernel
$r^{\mathrm{Dr}}(z)=k\,\Omega_{\mathrm{tr}}/z$ by the finite Casimir
rescaling $k\,\Omega_{\mathrm{tr}}\mapsto\Omega/(k+h^\vee)$.
The classical Yang--Baxter equation for $r^{\mathrm{Dr}}(z)$
follows from the infinitesimal braid relation
$[\Omega_{ij}, \Omega_{ik} + \Omega_{jk}] = 0$, which is the
collision depth~$2$ projection of the modular MC equation
\textup{(}Theorem~\textup{\ref{thm:collision-depth-2-ybe}})}.
\end{theorem}

\begin{proof}
The Kac--Moody OPE
$J^a(z)\, J^b(w) \sim k\,\delta^{ab}/(z-w)^2 + f^{ab}_{\;\;c}\,
J^c(w)/(z-w)$ has poles of order~$2$ and~$1$. The bar propagator
$d\log(z-w)$ absorbs one power, so the collision residue
extracts the simple-pole coefficient: the zeroth product
$a_{(0)}b = f^{ab}_{\;\;c} J^c$ (the Lie bracket). The double-pole
term $k\,\delta^{ab}$ contributes to the curvature~$m_0$, not to
the collision residue
(Computation~\ref{comp:sl2-collision-residue-kz}, Step~2).
Dualizing via the level-shifted Killing form $(k+h^\vee)\kappa$
(the Sugawara normalization) gives the Casimir $r$-matrix
$\Omega/((k+h^\vee)z)$, which is~\eqref{eq:hdm-face-5}.
At the critical level $k = -h^\vee$ the denominator vanishes,
tracking the Sugawara singularity.
At $k = 0$ the collision residue is $\Omega/(h^\vee z) \neq 0$:
the Lie bracket of~$\fg$ persists at vanishing level, so the
abelian vanishing $r = 0$ that characterizes the Heisenberg
$r$-matrix does not hold for non-abelian Kac--Moody.
The Yang--Baxter equation is then the IBR of
Theorem~\ref{thm:collision-depth-2-ybe} after evaluation against
the Casimir, which is the pure algebraic content of the Arnold
relation tensored with the structure constants of~$\fg$. The full
derivation is the shadow/KZ comparison theorem
(Theorem~\ref{thm:shadow-connection-kz}).
\end{proof}

Face~5 is the historical face: the original spectral $r$-matrix that
launched the theory of quantum groups and Yangians. The
identification with the collision residue of $\Theta_\cA$ realizes
Drinfeld's $r$-matrix as a specific projection of a much larger
universal object: $r_\cA(z)$ exists for every modular Koszul chiral
algebra, not just affine Kac--Moody, and Drinfeld's case is the
$\widehat{\fg}_k$ specialization.

\begin{remark}[Face 5 on 4d class-$\mathcal S$ punctured spheres: Beem--Rastelli as $\widehat{\fg}_k$ specialisation]
\label{rem:holo-master-class-S}
\index{class-$\mathcal S$!Beem--Rastelli chiral algebra at slot~(4)}
\index{VOA/$4$d $\mathcal N=2$ correspondence!Drinfeld $r$-matrix face}
The $\widehat{\fg}_k$ specialisation of Face~5 sits at a precise
physical checkpoint in the landscape surveyed in
Part~\ref{part:standard-landscape}: the 4d $\mathcal N = 2$ theory
$T_g[\Sigma_{0,n}]$ of class~$\mathcal S$ type $A_1$ on an $n$-punctured
sphere, whose Schur-index chiral algebra by the
Beem--Lemos--Liendo--Peelaers--Rastelli--van Rees
correspondence~\cite{BLLPRvR15} is
$\chi[T[A_1,\Sigma_{0,n}]] = \widehat{\mathfrak{sl}}_2$ at the
4d-2d-matched level $k_{2d} = -2$ with central charge
$c_{2d} = -12\,c_{4d}$. On the class-$\mathcal S$ sphere
$\Sigma_{0,24}$ with twenty-four regular punctures, the
representative 4d trinion census
(Gaiotto~\cite{Gaiotto12}, Chacaltana--Distler~\cite{ChacaltanaDistler10})
gives $c_{4d} = 107/6$ (Shapere--Tachikawa 2008 \S 3, with the
canonical normalization recorded in Vol~I
Section~\ref{sec:thqg-open-closed-realization}) and
hence $c_{2d} = -214$, placing the
associated $\widehat{\mathfrak{sl}}_2$ at the VOA-side central charge
$c(\widehat{\mathfrak{sl}}_2, k = -2) = -6$ rescaled by the
class-$\mathcal S$ Schur normalisation; the Drinfeld $r$-matrix
$r^{\mathrm{Dr}}(z) = k\,\Omega_{\mathrm{tr}}/z$ in
\eqref{eq:hdm-drinfeld-rmatrix} specialises at $k = -2$ to
$-2\,\Omega_{\mathrm{tr}}/z$, which in KZ normalisation
\eqref{eq:hdm-face-5} becomes
$r^{KM,\mathfrak{sl}_2}_{-2}(z) = \Omega / ((k+h^\vee) z) =
\infty \cdot \Omega / z$ at the critical level $k + h^\vee = 0$.
This critical-level divergence of the KZ-normalised collision residue
\emph{is} the class-$\mathcal S$ statement that the Schur-index chiral
algebra lives at the Sugawara denominator zero locus of
Face~5; the Drinfeld $r$-matrix in trace-form normalisation remains
finite ($k = -2$, hence
$-2\,\Omega_{\mathrm{tr}}/z \in \mathfrak{sl}_2 \otimes \mathfrak{sl}_2
[\![z^{-1}]\!]$), so the 4d parent class of
$\Sigma_{0,24}$ is read off cleanly from Face~5 in the trace-form
gauge. The Beem--Rastelli central-charge identity
$c_{2d} = -12\,c_{4d}$ is the numerical incarnation of the
4d-to-2d specialisation of slot~(5) at the sphere vacuum.
\end{remark}

\section{Face 6: the Sklyanin Poisson bracket of Semenov-Tian-Shansky}
\label{sec:hdm-face-6}

The 1983 paper of Semenov-Tian-Shansky~\cite{STS83} formulated the
notion of a classical $r$-matrix from the perspective of Poisson
geometry and showed that any classical $r$-matrix on a Lie algebra
$\fg$ defines a quadratic Poisson bracket on the dual $\fg^*$, the
\emph{Sklyanin Poisson bracket}, by
\begin{equation}\label{eq:hdm-sklyanin-bracket}
\{f, g\}_{\mathrm{STS}}(\xi)
\;=\;
\bigl\langle \xi,\, [df,\, dg]_{r^{\mathrm{cl}}}\bigr\rangle,
\qquad
[X,\, Y]_{r^{\mathrm{cl}}}
\;:=\;
[r^{\mathrm{cl}}(X),\, Y] + [X,\, r^{\mathrm{cl}}(Y)],
\end{equation}
where $f, g \in C^\infty(\fg^*)$, $\xi \in \fg^*$, and
$r^{\mathrm{cl}}\colon \fg \to \fg$ is a linear endomorphism whose
antisymmetric part defines a $2$-tensor in $\fg \otimes \fg$. The
classical Yang--Baxter equation for $r^{\mathrm{cl}}$ is equivalent
to the Jacobi identity for $\{\cdot,\cdot\}_{\mathrm{STS}}$. This is
the Poisson-geometric face of the spectral $r$-matrix: where Faces~1
and~5 produce algebraic morphisms or operator-valued functions, Face~6
produces an actual Poisson bracket on a phase space.

\begin{theorem}[Face~6: Sklyanin bracket; \ClaimStatusProvedHere\ (identification with classical limit of collision residue); \ClaimStatusProvedElsewhere\ (Sklyanin bracket: Semenov-Tian-Shansky 1983)]
\label{thm:hdm-face-6}
For $\cA = \widehat{\fg}_k$, the Sklyanin Poisson bracket on
$\fg^*$ defined by~\eqref{eq:hdm-sklyanin-bracket} with
level-$k$ Drinfeld rational $r$-matrix
$r^{\mathrm{cl}} = r^{\mathrm{Dr}} = k\Omega_{\mathrm{tr}}/z$
\textup{(}vanishing at $k=0$ and related at generic $k$ to the
KZ-normalized collision residue by
$k\,\Omega_{\mathrm{tr}}\mapsto
\Omega/(k+h^\vee)$ through finite Casimir rescaling\textup{)} is the classical
limit of the collision residue $r_\cA(z)$ in the sense that the
Lie--Poisson bracket on $\fg^* \cong \mathrm{Spec}\,\mathrm{gr}^{\mathrm{Li}}(\cA)$
agrees with the leading-order term of $r_\cA(z)$ under the classical
limit functor of Face~3 (genus~$0$ only). Equivalently, the Sklyanin bracket is the
$\hbar \to 0$ limit of the Yangian quantum bracket
$[\Delta_\hbar(x), y \otimes 1]/\hbar$, where $\Delta_\hbar$ is the
Yangian coproduct of Drinfeld~\cite{Drinfeld85}.
\end{theorem}

\begin{proof}
The Sklyanin bracket is constructed from a spectral classical
$r$-kernel. For
$r^{\mathrm{cl}} = r^{\mathrm{Dr}} = k\Omega_{\mathrm{tr}}/z$
\textup{(}level $k$, vanishing at $k=0$\textup{)}, the tensor
$\Omega_{\mathrm{tr}}$ is symmetric, but the full spectral kernel is
skew in the standard sense
$r^{\mathrm{Dr}}_{21}(-z)=-r^{\mathrm{Dr}}_{12}(z)$. Thus the
Sklyanin bivector is obtained from the spectral skew kernel, not by
antisymmetrising the Casimir tensor alone. Its residue at $z=0$ is
the invariant Casimir controlling the Lie--Poisson bracket on
$\fg^*$. Under the classical limit
$\mathrm{cl}\colon \cA \to \mathrm{gr}^{\mathrm{Li}}(\cA)$ of
Face~3, the chiral OPE pole structure of $\widehat{\fg}_k$ becomes
the Lie--Poisson bracket on $\fg^*$, and the binary collision
residue $r^{KM,\fg}_{\mathrm{KZ},k}(z) = \Omega/((k+h^\vee)z)$
is the same generic family as
$r^{KM,\fg}_{\mathrm{coll},k}(z)=k\Omega_{\mathrm{tr}}/z$ only after
the finite Casimir-rescaling gauge
$k\,\Omega_{\mathrm{tr}}\mapsto\Omega/(k+h^\vee)$, so it becomes the
Lie--Poisson bracket twisted by $\Omega$. This is the Sklyanin
bracket up to the level normalization. The $\hbar \to 0$ statement
is Drinfeld's classical-limit theorem: the Yangian
$Y_\hbar(\fg) = \mathcal{U}(\fg[\![z^{-1}]\!])\otimes\mathbb{C}[\![\hbar]\!]$
quantizes the Sklyanin bracket, and the leading term in $\hbar$ of
the coproduct commutator recovers the Sklyanin bracket as a
classical Poisson structure on the dual. This is part~(i) of the
three-parameter $\hbar$ identification
(\S\ref{sec:hdm-hbar-identification}).
\end{proof}

Face~6 records the Poisson-geometric origin of the spectral
$r$-matrix. The Semenov-Tian-Shansky construction shows that the
$r$-matrix is the structure constant of a Poisson bracket; the
collision residue identification shows that this Poisson bracket is
the classical limit of the bar-cobar twisting morphism.

\section{Face 7: the Gaudin Hamiltonians of Feigin--Frenkel--Reshetikhin}
\label{sec:hdm-face-7}

The 1994 paper of Feigin, Frenkel, and Reshetikhin~\cite{FFR94}
constructed an integrable system on the tensor product of $n$ finite
representations of a simple Lie algebra~$\fg$, called the
\emph{Gaudin model}, by means of the commuting Hamiltonians
\begin{equation}\label{eq:hdm-gaudin-hamiltonians}
H_i^{\mathrm{Gaudin}}
\;=\;
\sum_{j \neq i} \frac{\Omega_{ij}}{z_i - z_j}
\;\in\;
\mathcal{U}(\fg)^{\otimes n} \otimes \mathbb{C}[z_1, \ldots, z_n,
\Delta^{-1}],
\qquad
i = 1, \ldots, n,
\end{equation}
where $z_1, \ldots, z_n$ are distinct points on the line and
$\Omega_{ij}$ is the Casimir acting on the $i$-th and $j$-th tensor
factors. The commutativity $[H_i^{\mathrm{Gaudin}}, H_j^{\mathrm{Gaudin}}] = 0$
follows from the infinitesimal braid relation and is the
algebraic-Bethe-ansatz starting point for solving the Gaudin model.
FFR proved that the Gaudin Hamiltonians arise as the critical-level
$k = -h^\vee$ limit of the affine Sugawara construction, and that
the Gaudin eigenfunctions are computed by the Bethe ansatz.

\begin{theorem}[Face~7: Gaudin Hamiltonians; \ClaimStatusConditional]
\label{thm:hdm-face-7}
For $\cA = \widehat{\fg}_k$, the Gaudin Hamiltonians of
Feigin--Frenkel--Reshetikhin are the kernel of the modular shadow
connection on the sphere:
\begin{equation}\label{eq:hdm-face-7}
H_i^{\mathrm{Gaudin}}
\;=\;
(k+h^\vee)\,
\sum_{j \neq i} \bigl[r^{KM,\fg}_{\mathrm{KZ},k}(z_i - z_j)\bigr]_{i,j}
\;=\;
\sum_{j \neq i} \frac{\Omega_{ij}}{z_i - z_j}.
\end{equation}
The level dependence cancels: the factor $(k+h^\vee)$ in the
prefactor absorbs the denominator $(k+h^\vee)$ of the collision
residue~\eqref{eq:hdm-face-5}, leaving a level-independent
expression.
Equivalently, $H_i^{\mathrm{Gaudin}} = (k+h^\vee)\, H_i^{\mathrm{GZ}}$,
where $H_i^{\mathrm{GZ}}$ is the Gaiotto--Zeng Hamiltonian of
Face~4. The classical Yang--Baxter equation for $r^{\mathrm{Dr}}(z)$
is the commutativity of the Gaudin Hamiltonians.
\end{theorem}

\begin{proof}
The collision residue for affine Kac--Moody is
$r^{KM,\fg}_{\mathrm{KZ},k}(z) = \Omega/((k+h^\vee)z)$
(Theorem~\ref{thm:hdm-face-5}; the Sugawara denominator $k+h^\vee$
arises from dualizing the Lie bracket via the level-shifted Killing
form). Substituting into the Gaiotto--Zeng
formula~\eqref{eq:hdm-face-4} of Face~4 and rescaling by
$(k+h^\vee)$ produces~\eqref{eq:hdm-face-7}, which is the
standard Gaudin Hamiltonian formula. The commutativity
$[H_i^{\mathrm{Gaudin}}, H_j^{\mathrm{Gaudin}}] = 0$ is equivalent
to $[\Omega_{ij}, \Omega_{ik} + \Omega_{jk}] = 0$, which is the
infinitesimal braid relation
(Theorem~\ref{thm:collision-depth-2-ybe}). The critical-level
limit $k \to -h^\vee$ is the Feigin--Frenkel--Reshetikhin recovery
of the Gaudin model from the affine Sugawara construction; in the
present framework, the critical level is the locus where the
Sugawara normalization $1/(k+h^\vee)$ diverges and the Gaudin Hamiltonians
become the singular limit of the modular shadow connection.
\end{proof}

Face~7 is the integrable-system face: the Gaudin Hamiltonians are
the historically first commuting integrable Hamiltonians built from
a classical $r$-matrix, and they are recovered here as a special
case of the modular shadow connection on the sphere. The
identification with the Gaiotto--Zeng face is exact, up to the
level rescaling.

\section{The seven-way master theorem}
\label{sec:hdm-master}

We assemble the seven faces on their common comparison domains. Let
$\cA$ be a modular Koszul chiral algebra, and let $r_\cA(z)$ be the
collision residue of $\Theta_\cA$. The seven-face master theorem
asserts that the following comparison diagram commutes whenever the
displayed faces are simultaneously defined:
\begin{equation}\label{eq:hdm-master-diagram}
\xymatrix{
\pi_\cA \in \mathrm{Tw}(\barBch(\cA), \cA)
 \ar@{<->}[r]^-{\textrm{(F1)}}
 \ar@{<->}[d]_{\textrm{(F3)}}
& r^{\mathrm{DNP}}(z) \in \cA^! \otimes \cA^! [\![z^{-1}]\!]
 \ar@{<->}[d]^{\textrm{(F2)}}\\
\{\cdot_\lambda \cdot\}\;\textrm{(KZ25)}
 \ar@{<->}[r]
 \ar@{<->}[d]
& \{H_i^{\mathrm{GZ}}\}\;\textrm{(GZ26)}
 \ar@{<->}[d]^{\textrm{(F4)}}\\
\{\cdot,\cdot\}_{\mathrm{STS}}\;\textrm{(STS83)}
 \ar@{<->}[r]^-{\textrm{(F6)}}
& \{H_i^{\mathrm{Gaudin}}\}\;\textrm{(FFR94)}\\
& k\,\Omega_{\mathrm{tr}}/z\;\textrm{(Drinfeld85)}
 \ar@{<->}[u]^{\textrm{(F5--F7)}}
}
\end{equation}

\begin{theorem}[Seven-way master theorem; \ClaimStatusConditional]
\label{thm:hdm-seven-way-master}
\index{seven faces!master theorem|textbf}
\index{collision residue!seven-way master|textbf}
Let $I(\cA)\subset\{1,\ldots,7\}$ be the set of faces whose comparison
hypotheses hold for~$\cA$. Then the faces
$\{F_i\}_{i\in I(\cA)}$ of $r_\cA(z)$ are equal as elements of
$\mathrm{End}_\cA(2)[\![z^{-1}]\!]$ after the standard normalization.
For $\cA=\widehat{\fg}_k$ at non-critical level in the affine KZ lane,
the concrete master equation is
\begin{equation}\label{eq:hdm-master-equation}
\boxed{\;
r^{KM,\fg}_{\mathrm{KZ},k}(z)
\;=\;
\pi_\cA
\;=\;
r^{\mathrm{DNP}}(z)
\;=\;
\sum_{n \geq 0} \frac{\mathrm{cl}(a_{(n)}b)}{z^{n+1}}
\;=\;
H_i^{\mathrm{GZ}}\bigr|_{j\textrm{-coeff}}
\;=\;
\frac{\Omega}{(k+h^\vee)z}
\;=\;
\bigl\{\cdot,\cdot\bigr\}_{\mathrm{STS}}
\;=\;
\frac{H_i^{\mathrm{Gaudin}}}{k + h^\vee}\;}
\end{equation}
in the genus-zero, two-point sector of the modular cyclic deformation
complex. For a general modular Koszul~$\cA$, the same equation is read
with the available face-specific expressions replacing the affine
Casimir formula.
\end{theorem}

\begin{proof}
Face~1 is Theorem~\ref{thm:hdm-face-1}. Face~2 is
Theorem~\ref{thm:hdm-face-2}. Face~3 is Theorem~\ref{thm:hdm-face-3} (genus~$0$ only).
Face~4 is Theorem~\ref{thm:hdm-face-4}. Face~5 is
Theorem~\ref{thm:hdm-face-5}. Face~6 is Theorem~\ref{thm:hdm-face-6}.
Face~7 is Theorem~\ref{thm:hdm-face-7}. The transitive closure of
the available equalities yields the master
equation~\eqref{eq:hdm-master-equation}.
The normalization conventions match on the common domain. In the
affine lane the factor $1/(k+h^\vee)$ appears in Faces~5 and~7
because the KZ collision residue is dualized through the
level-shifted Killing form; Faces~3 and~6 see the same symmetric
Casimir after the finite trace-form/KZ rescaling.
\end{proof}

\begin{remark}[Scope: which faces hold for which~$\cA$]
\label{rem:hdm-master-scope}
Faces~1, 3, 4 hold for every modular Koszul chiral algebra $\cA$:
they are universal in the algebraic engine. Face~2 holds for
$\cA$ realized as the closed colour of a 3d HT theory in the sense
of~\cite{DNP25}; this includes only those standard-landscape algebras
for which that closed-colour realization and the Vol.~II comparison
are proved. Faces~5--7 are specialized to
$\cA = \widehat{\fg}_k$, where the historical $r$-matrix and Gaudin
formulas live. For a general modular Koszul~$\cA$ these affine
Casimir formulas are not asserted; the corresponding slot is the
available face-specific expression of the universal collision residue
$r_\cA(z)$.
\end{remark}

\subsection{Averaging-map compatibility on the seven faces}
\label{subsec:hdm-averaging-compatibility}

The seven-way master theorem is stated at the level of the bar
hub~$F_1$, where the universal collision residue
$r_\cA(z) = \Res^{\mathrm{coll}}_{0,2}(\Theta_\cA)$ lives in the
$E_\infty$-modular dGLA $\mathfrak{g}_\cA^{\mathrm{mod}}$. The
seven faces $F_2, \dots, F_7$ are constructed in different ambient
homotopy frames: some symmetric ($F_3, F_5, F_6, F_7$ in their
classical-Sugawara presentations), some intrinsically ordered
($F_1$ as a bar twisting morphism, $F_2$ as a colour-ordered
DNP $R$-twist, $F_4$ as a Gaiotto--Zeng commuting-differential
generator on a coloured boundary). The averaging map
\[
 \av\colon \mathfrak{g}_\cA^{E_1} \;\twoheadrightarrow\;
 \mathfrak{g}_\cA^{\mathrm{mod}}
\]
of Definition~\ref{def:e1-modular-convolution} therefore acts
non-trivially on the ordered faces and trivially on the symmetric
faces; the master theorem demands that the two records be
consistent. The following proposition records the compatibility
class-by-class and face-by-face. It is the compatibility interface between
the bar-intrinsic master theorem
(Theorem~\ref{thm:hdm-seven-way-master}) and the
$\Sigma_n$-symmetric descent
(Proposition~\ref{prop:symmetric-descent} of
Chapter~\ref{chap:e1-modular-koszul}).

\begin{proposition}[Averaging compatibility of the seven faces;
 \ClaimStatusConditional]
\label{prop:hdm-averaging-compatibility}
\index{averaging map!seven faces compatibility|textbf}
\index{seven faces!averaging compatibility|textbf}
\index{collision residue!ordered vs symmetric image}
Let $\cA$ be a chirally Koszul chiral algebra in the standard
landscape, and let $r_\cA^{(F_i)}(z)$ for $i \in \{1,\dots,7\}$
denote the seven faces of the collision residue in their
canonical chain-level presentations
(Definition~\ref{V2-def:face-space} of Vol~II,
Chapter~\ref{V2-ch:grt-seven-faces}). Write
$r_\cA^{(F_i),E_1}(z)$ for the lift of face $F_i$ to the
$E_1$-modular convolution dGLA $\mathfrak{g}_\cA^{E_1}$ obtained
by either (a) the bar-intrinsic universal MC element
$\Theta_\cA^{E_1}$ projected to bar-degree $2$ in the canonical
chain model of $F_i$, when this lift is canonical, or (b) the
trivial lift via a chosen section
$\iota^{\mathrm{rib}}_{0,2}\colon \cM_\Com(0,2) \to
\cM_\Ass(0,2)$ when $F_i$ is intrinsically symmetric. Then for
each shadow class $\mathbf{X} \in \{\mathbf{G}, \mathbf{L},
\mathbf{C}, \mathbf{M}\}$:
\begin{enumerate}[label=\textup{(\roman*)}]
\item \emph{Class~$\mathbf{G}$ (Gaussian; archetype $\cH_k$).}
 For every face $F_i$,
 \[
 \av\bigl(r_{\cH_k}^{(F_i),E_1}(z)\bigr) \;=\;
 r_{\cH_k}^{(F_i)}(z) \;=\; \frac{k}{z}\, J\otimes J,
 \]
 with no Sugawara shift. Equivalently the kernel
 $\ker(\av)|_{r}$ vanishes at every face: the Heisenberg algebra
 is $E_\infty$-chiral
 (Proposition~\ref{prop:symmetric-descent}(i)) and the
 averaging is the identity at degree~$2$.
\item \emph{Class~$\mathbf{L}$ (Lie/tree; archetype
 $\widehat{\fg}_k$).} For every face $F_i$, the binary
 collision residue is the symmetric Killing--Casimir tensor in
 one of the two normalisations recorded in the landscape census
 (lines 380--381, 644 of \texttt{landscape\_census.tex}):
 \begin{align*}
 r_{\widehat{\fg}_k}^{(F_i)}(z) \;=\;
 k\,\Omega_{\mathrm{tr}}/z
 \quad &\text{(trace-form, ``Drinfeld 1985 / DNP'' lane),} \\
 r_{\widehat{\fg}_k}^{(F_i)}(z) \;=\;
 \Omega/((k+h^\vee)z)
 \quad &\text{(KZ lane, ``FFR / Sugawara / KZ-PVA'' lane).}
 \end{align*}
 The two normalisations are matched by the Casimir rescaling
 $k\,\Omega_{\mathrm{tr}}\mapsto
 \Omega/(k+h^\vee)$, which is a finite inner gauge
 $\Phi^{\mathrm{lane}} \in \GRT^{\mathrm{fin}}$ in the sense
 of Vol~II, \S\ref{V2-sec:grtfin}. For every face $F_i$ the
 averaging map is the identity on the binary slice in either
 lane:
 \[
 \av\bigl(r_{\widehat{\fg}_k}^{(F_i),E_1}(z)\bigr) \;=\;
 r_{\widehat{\fg}_k}^{(F_i)}(z),
 \]
 and the kernel $\ker(\av)|_r$ vanishes lane-by-lane.
\item \emph{Class~$\mathbf{C}$ (contact; archetype $\beta\gamma$).}
 For every face $F_i$,
 \[
 \av\bigl(r_{\beta\gamma}^{(F_i),E_1}(z)\bigr) \;=\;
 r_{\beta\gamma}^{(F_i)}(z) \;=\; 0,
 \]
 the residue is zero on every face. The non-trivial datum
 is carried by the higher shadow $r_{\max}(\beta\gamma) = 4$
 component of the obstruction tower
 (Theorem~\ref{thm:genus-universality}, class~$\mathbf{C}$
 entry); $\av$ is compatible with the higher-shadow projection
 face by face by the same coinvariant argument.
\item \emph{Class~$\mathbf{M}$ (mixed/infinite; archetype
 $\mathrm{Vir}_c$).} For every face $F_i$ at which $r$ admits
 a closed-form chain-level presentation,
 \[
 \av\bigl(r_{\mathrm{Vir}_c}^{(F_i),E_1}(z)\bigr) \;=\;
 r_{\mathrm{Vir}_c}^{(F_i)}(z) \;=\;
 \frac{c/2}{z^3} + \frac{2T}{z},
 \]
 with the cubic central pole and the simple field-dependent pole
 (the class~$\mathbf{M}$ landscape census convention). The kernel of
 $\av|_r$ at degree~$2$ vanishes; the
 GRT-data that the $\Sigma_2$-coinvariant projection annihilates
 in Proposition~\ref{prop:symmetric-descent}(ii) lives at degree
 $\geq 3$, where it is detected by the higher shadow tower
 $S_3, S_4, S_5, \dots$ and not by the binary residue~$r$.
\end{enumerate}
\end{proposition}

\begin{proof}
At binary degree $n = 2$ the averaging map of
Definition~\ref{def:e1-modular-convolution} is
\[
 \av(\phi)(g, 2) \;=\; \tfrac{1}{2}\bigl(\phi(g,2) +
 \tau \cdot \phi(g,2)\bigr),
\]
where $\tau \in \Sigma_2$ is the transposition acting on the
two external legs. The proposition reduces to the question
\emph{is $r_\cA^{(F_i),E_1}(z)$ already $\Sigma_2$-invariant
in its chosen chain model?}; if yes, $\av$ is the identity
at degree~$2$ on $F_i$; if no, $\av$ extracts the symmetric part
and the antisymmetric part lies in $\ker\av$.

\smallskip
\emph{(i) Class~$\mathbf{G}$.} The Heisenberg single-generator
double-pole OPE $J(z)J(w) \sim k/(z-w)^2$ is symmetric in
$(J, J)$ by the abelian commutation relation. The $d\log$
absorption preserves this symmetry, so the chain-level lift
of $r^{\mathrm{Heis}}_k(z) = (k/z)\,J\otimes J$ at every face $F_i$ is
$\Sigma_2$-invariant: $\tau \cdot (J\otimes J) = J\otimes J$.
Therefore $\av(r^{(F_i),E_1}_{\cH_k}(z)) = r^{(F_i),E_1}_{\cH_k}(z)
= r_{\cH_k}^{(F_i)}(z)$ pointwise. The Sugawara correction
vanishes because $\cH_k$ is abelian.

\smallskip
\emph{(ii) Class~$\mathbf{L}$.} The diagonal off-diagonal
OPE $J^a(z)J^b(w) \sim k\,\kappa^{ab}/(z-w)^2 + f^{ab}{}_c
J^c(w)/(z-w)$ has symmetric double-pole part
$k\,\kappa^{ab}$ (because $\kappa^{ab} = \kappa^{ba}$ for the
Killing form) and antisymmetric simple-pole part
$f^{ab}{}_c J^c$ (because $f^{ab}{}_c = -f^{ba}{}_c$). The
$d\log$ absorption strips the simple pole, leaving only the
symmetric Casimir $\Omega^{ab} = \kappa^{ab} I_a \otimes I_b$.
Hence at the bar hub $F_1$ the chain-level lift of the
family-normalized Kac--Moody kernel
$r^{KM,\fg}_{\mathrm{coll},k}(z)$ or
$r^{KM,\fg}_{\mathrm{KZ},k}(z)$ is
$\Sigma_2$-invariant: $\av$ is the identity. The transport
to any other face $F_i$ ($i = 2, \ldots, 7$) is via a finite
inner gauge $\Phi_{1i} \in \GRT^{\mathrm{fin}}$
(Theorem~\ref{V2-thm:seven-faces-as-GRT-torsor} of Vol~II), which
is a finite Casimir rescaling and (in the KZ lane) a simple-pole
$k \mapsto k+h^\vee$ shift. The rescaled tensor remains
$\Sigma_2$-invariant because $\Omega/(k+h^\vee)$ is a scalar
multiple of the same Casimir~$\Omega^{ab}$, and scalar
rescalings preserve the $\Sigma_2$-action on
$\Omega^{ab} I_a \otimes I_b$. The two normalisation lanes belong
to the same gauge class (the chained equality
$k\,\Omega_{\mathrm{tr}} = \Omega/(k+h^\vee)$ is shorthand, not a
rational-function identity, for a Casimir rescaling
in $\GRT^{\mathrm{fin}}$, per
Section~\ref{V2-sec:grtfin} of Vol~II); the per-face assignment of
lane (trace-form vs KZ) is a convention determined by the
historical literature attached to each $F_i$ (cf.\ the
$\Phi_{12}, \dots, \Phi_{17}$ enumeration in Vol~II,
\S\ref{V2-sec:grt-torsor-thm}).

\smallskip
\emph{(iii) Class~$\mathbf{C}$.} Both the $\beta\gamma$ pair
and the $bc$ ghost system have OPEs of the form
$\beta(z)\gamma(w) \sim 1/(z-w)$, $\gamma(z)\beta(w) \sim
\pm 1/(z-w)$ (sign by parity), with $p_{\max} = 1$ and
$k_{\max} = 0$. The $d\log$ absorption strips the only pole,
leaving zero. Hence $r_{\beta\gamma}^{(F_i)}(z) = 0$ for
every face by direct computation (landscape census,
class~$\mathbf{C}$ entry in
\texttt{landscape\_census.tex}). The averaging of zero is
zero, trivially. The non-trivial information of class~$\mathbf{C}$
is the shadow obstruction at $r = 4$
(Theorem~\ref{thm:genus-universality}, class~$\mathbf{C}$
witness, with the class-$\mathbf{C}$ shadow tower
$r_{\max}(\beta\gamma) = 4$
recorded in the landscape census). The averaging-compatibility
statement at the binary level is therefore vacuous: both sides
are zero.

\smallskip
\emph{(iv) Class~$\mathbf{M}$.} The Virasoro $TT$ OPE
$T(z)T(w) \sim (c/2)/(z-w)^4 + 2T(w)/(z-w)^2 + \partial T(w)/(z-w)$
has all three pole terms symmetric in $(T, T)$ as functions
of the legs after pull-back along the $\tau$-swap: under
the transposition $\tau\colon (z, w) \mapsto (w, z)$, the
quartic central pole maps to itself (scalar), the double-pole
$2T(w)/(z-w)^2$ maps to $2T(z)/(z-w)^2$ which on the
deconcatenation chain model is $\tau$-equivalent up to a
$d\log$-derivative, and the simple-pole $\partial T(w)/(z-w)$
contributes the antisymmetric component that $d\log$
absorption removes. The chain-level lift of
$r^{\mathrm{Vir}}_c(z) = (c/2)/z^3 + 2T/z$ at every face is therefore
$\Sigma_2$-invariant once the
$d\log$-absorption is applied; the cubic-plus-simple pole pattern,
not a quartic pole, is the post-absorption form. Hence $\av$ is the
identity at degree~$2$ on every
face. The kernel
$\ker(\av) \cap \mathfrak{g}_{\mathrm{Vir}_c}^{E_1}|_{(g,n) = (0,2)}$
vanishes; the GRT-data that the $\Sigma_n$-coinvariant
projection annihilates in
Proposition~\ref{prop:symmetric-descent}(ii) lives at
degree $n \geq 3$, where it is the higher Virasoro shadow
$S_3, S_4, S_5, \dots$
(Theorem~\ref{thm:virasoro-shadow-generating-function}, Vol~I,
Chapter~\texttt{w\_algebras})
and not the binary $r$.
\end{proof}

\begin{remark}[The four-class compatibility table]
\label{rem:hdm-averaging-compatibility-table}
The four-class statement of
Proposition~\ref{prop:hdm-averaging-compatibility} can be
summarised by the dichotomy:
\begin{itemize}
\item \emph{Identity at degree~$2$.} For class~$\mathbf{G}$,
 $\mathbf{L}$, $\mathbf{M}$, the binary residue $r_\cA(z)$ is
 already $\Sigma_2$-invariant on every face's canonical chain
 model, so $\av(r) = r$ and the seven-face master equation
 lifts unchanged from $\mathfrak{g}_\cA^{\mathrm{mod}}$ to
 $\mathfrak{g}_\cA^{E_1}$.
\item \emph{Vacuous at degree~$2$.} For class~$\mathbf{C}$,
 $r = 0$ on every face, so $\av(r) = 0$ trivially. The
 non-trivial averaging compatibility moves to the higher shadow
 tower at $r_{\max} = 4$.
\end{itemize}
The structural reading is that the GRT-torsor structure of
$\Face(\cA)$ (Vol~II, Theorem~\ref{V2-thm:seven-faces-as-GRT-torsor})
acts on the binary slice through finite inner gauges only on
class~$\mathbf{L}$ (the trace-form $\leftrightarrow$ KZ
rescaling), and acts trivially on classes $\mathbf{G},
\mathbf{C}, \mathbf{M}$ at the binary level. The non-trivial
GRT data first appears at higher degree $n \geq 3$, and it is
this data that the averaging map projects out via
$\Sigma_n$-coinvariants for $n \geq 3$ in the genuine $E_1$-chiral
case (Proposition~\ref{prop:symmetric-descent}(ii)).
\end{remark}

\begin{remark}[Cross-volume bridge]
\label{rem:hdm-averaging-compatibility-xvol}
\index{cross-volume bridge!averaging compatibility}
Proposition~\ref{prop:hdm-averaging-compatibility} is the
binary-degree shadow of the GRT-torsor structure of Vol~II
(\S\ref{V2-sec:grt-torsor-thm}, Theorem~\ref{V2-thm:seven-faces-as-GRT-torsor}).
On the Vol~III side, the analogous proposition for
$r_{CY}(z) = \Res^{\mathrm{coll}}_{0,2}(\Theta_{A_\cC})$
becomes a statement about the $E_1$-chiral lift of the seven CY
faces (Vol~III,
Theorem~\ref{V3-thm:seven-faces-of-r-cy}); the four-class
compatibility table extends to the CY setting via the
dimension-stratified correspondence $\{\Phi_d\}_{d\ge 1}$
(Vol~III, \S\ref{sec:cy-holographic-anchors}). The
proposition is therefore a single cross-volume bridge: the
seven faces of $r_\cA(z)$ in Vol~I, the GRT-orbit representatives
in Vol~II, and the seven CY faces of $r_{CY}(z)$ in Vol~III
all share the same averaging-compatibility table
class-by-class. The class-$\mathbf{L}$ instance at
$\widehat{\mathfrak{sl}}_{2,k}$ on a genus-$1$ curve is
verified by direct chain-level computation in
\texttt{compute/lib/averaging\_kernel\_sl2\_genus1\_engine.py}
(test in
\texttt{compute/tests/test\_averaging\_kernel\_sl2\_genus1\_engine.py}),
which witnesses the kernel statement
$\ker(\av)|_r = 0$ at degree~$2$ on the bar hub for
$\widehat{\mathfrak{sl}}_2$ at non-critical level.
\end{remark}

\subsection{Vol~I face-cocycle record: degree-$2$ collapse of the
\texorpdfstring{$\GRT^{\mathrm{fin}}$}{GRT-fin} torsor}
\label{subsec:hdm-grt-cocycle-record}

\providecommand{\Face}{\operatorname{Face}}
\providecommand{\fgrt}{\mathfrak{grt}}
\providecommand{\GRTone}{\operatorname{GRT}_{1}}

The seven-way master theorem
(Theorem~\ref{thm:hdm-seven-way-master}) asserts the equality of the
available chain-level presentations of the binary collision residue
$r_\cA(z)$ on their common comparison domain.
Vol~II's master theorem
(Theorem~\ref{V2-thm:seven-faces-as-GRT-torsor}) refines that equality to
the statement that the face space $\Face(\cA)$ is a principal
$\GRT(\mathbb{Q})$-torsor and that the seven historical faces are
Q-rational orbit representatives. The proposition below is the Vol~I
shadow of that torsor on the binary slice: it tabulates the seven
explicit cocycles $\Phi_{1i}\in\GRT^{\mathrm{fin}}$ that take the bar
hub $F_1$ to each face $F_i$, records the first bar-degree at which
each cocycle is non-trivial, and proves that the seven-way master
equality is the binary-degree collapse
$\GRT(\mathbb{Q})\twoheadrightarrow\GRT^{(2)}=\{1\}$ of the
$\fgrt$-action.

\begin{theorem}[Vol~I face-cocycle record; degree-$2$ collapse;
\ClaimStatusConditional]
\label{thm:hdm-grt-cocycle-record}
\index{seven faces!cocycle record|textbf}
\index{GRT torsor!degree-2 collapse|textbf}
\index{collision residue!MZV invisibility on binary slice|textbf}
Let $\cA$ be a chirally Koszul chiral algebra in the standard
landscape, and let $\Theta_\cA^{E_1}$ be its bar-intrinsic universal
Maurer--Cartan element in the ordered convolution dGLA
$\mathfrak{g}_\cA^{E_1}$ \textup{(}Definition~\textup{\ref{def:e1-modular-convolution}}\textup{)}.
Write $\GRT^{(\geq 3)}\subset\GRT(\mathbb{Q})$ for the pro-unipotent
radical of associators trivial at bar-degree $2$, and
$\GRT^{(2)}:=\GRT(\mathbb{Q})/\GRT^{(\geq 3)}$ for the binary-degree
quotient. Then:
\begin{enumerate}[label=\textup{(\roman*)}]
\item \emph{Cocycle record.} The seven cocycles
$\{\Phi_{1i}\}_{i=2}^{7}\subset\GRT(\mathbb{Q})$ realising
$F_i = \Phi_{1i}\cdot F_1$ \textup{(}Vol~II,
Theorem~\textup{\ref{V2-thm:seven-faces-as-GRT-torsor}}\textup{)}
all lie in the finite inner gauge subgroup $\GRT^{\mathrm{fin}}$.
Each $\Phi_{1i}$ is supported at a definite first non-trivial
bar-degree $d^\star(\Phi_{1i})\in\{2,3,\geq 3\}$ recorded in
Table~\textup{\ref{tab:hdm-cocycle-record}}.
\item \emph{Degree-$2$ collapse.} The natural projection
\[
 \mu_2\colon\;\GRT(\mathbb{Q})\;\longrightarrow\;\GRT^{(2)},
 \qquad
 \mu_2(\GRT^{\mathrm{fin}})\;=\;\{1\},
\]
sends every $\Phi_{1i}$ \textup{(}for $i=2,\dots,7$\textup{)} to the
identity. Consequently, on the bar-degree-$2$ slice the seven
faces are pointwise equal:
\[
 r_\cA^{(F_1)}(z) \;=\;
 r_\cA^{(F_2)}(z) \;=\;\cdots\;=\;
 r_\cA^{(F_7)}(z)
 \;\in\;
 \cA^{!}\otimes\cA^{!}\,[\![z^{-1}]\!],
\]
which is precisely the seven-way master equation
\eqref{eq:hdm-master-equation}.
\item \emph{MZV invisibility on the binary slice.} The Drinfeld
associator $\Phi_{\KZ}$ and the rational associator $\Phi_{\rat}$
satisfy
\[
 \Phi_{\KZ}\cdot r_\cA(z)
 \;=\;
 \Phi_{\rat}\cdot r_\cA(z),
\]
and the difference $\Phi_{\KZ}-\Phi_{\rat}$ is a sum of
$\fgrt$-elements of weight $w\geq 3$, hence supported at bar-degree
$\geq 3$ and invisible to the binary residue.
\item \emph{Class-by-class refinement.} Combined with
Proposition~\textup{\ref{prop:hdm-averaging-compatibility}} and
Remark~\textup{\ref{rem:hdm-averaging-compatibility-table}}, the
$\GRT$-cocycle action on the binary slice is non-trivial only on
class~$\mathbf{L}$, where $\Phi_{13}$ realises the
trace-form/KZ Casimir rescaling
$k\,\Omega_{\mathrm{tr}}\leftrightarrow\Omega/(k+h^\vee)$ as a
finite inner gauge \textup{(}Vol~II,
Proposition~\textup{\ref{V2-prop:grtfin-casimir}}\textup{)}; on classes
$\mathbf{G}$, $\mathbf{C}$, $\mathbf{M}$, the $\GRT^{\mathrm{fin}}$
action on the binary residue is the identity.
\end{enumerate}
\end{theorem}

\begin{table}[htbp]
\centering
\renewcommand{\arraystretch}{1.2}
\begin{tabular}{|c|l|c|c|l|}
\hline
\textbf{Face} & \textbf{Cocycle $\Phi_{1i}$} & $d^\star$
 & \textbf{In $\GRT^{\mathrm{fin}}$?} & \textbf{Non-trivial on classes} \\
\hline\hline
$F_1$ & $\Phi_{11}=\mathrm{id}$ (bar hub origin) &; & yes &; (origin) \\
\hline
$F_2$ & $R$-twist coproduct deformation
 $\Delta\to R\Delta R^{-1}$ at order $\hbar$
 & $2$ & yes & none (identity at degree~$2$) \\
\hline
$F_3$ & Casimir rescaling $\Omega_{\mathrm{tr}}\to\Omega/(2h^\vee)$
 $+$ level shift $k\to k+h^\vee$
 & $2$ & yes & class $\mathbf{L}$ \\
\hline
$F_4$ & Laplace transform of bar kernel
 & $2$ & yes & none (identity at degree~$2$) \\
\hline
$F_5$ & Yangian classical-limit of $R(u)=1+\hbar P/u+O(\hbar^2)$
 & $2$ & yes & class $\mathbf{L}$ \\
\hline
$F_6$ & Sugawara denominator $1/(k+h^\vee)$ inversion
 & $2$ & yes & class $\mathbf{L}$ \\
\hline
$F_7$ & Simple-pole truncation of bar-degree-$3$ component
 & $\geq 3$ & yes (G/L);
 acquires MZV at $w\geq 3$ on class $\mathbf{M}$ ($F_7'$)
 & class $\mathbf{M}$ at $d\geq 3$ \\
\hline
\end{tabular}
\caption{The seven Vol~I cocycles $\Phi_{1i}$ realising $F_i$ from
the bar hub $F_1$. The column $d^\star$ records the first bar-degree
at which $\Phi_{1i}$ is non-trivial; the right column records the
shadow class on which the cocycle action is non-trivial. On the
binary slice $d=2$, every cocycle in $\GRT^{\mathrm{fin}}$ projects
to the identity, recovering the seven-way master equation
\eqref{eq:hdm-master-equation}.}
\label{tab:hdm-cocycle-record}
\end{table}

\begin{proof}
\emph{(i) Cocycle record.} For each pair $(F_1,F_i)$ with
$i=2,\dots,7$, Vol~II's
Theorem~\ref{V2-thm:seven-faces-as-GRT-torsor} constructs the explicit
Q-rational associator $\Phi_{1i}\in\GRT(\mathbb{Q})$. The proof
there records that each $\Phi_{1i}$ either is a finite inner gauge
(Casimir rescaling, level shift, $R$-twist, Laplace transform,
simple-pole truncation) or differs from such a gauge by a
bar-degree-$\geq 3$ correction with no binary projection. The
finite-gauge subgroup $\GRT^{\mathrm{fin}}$ is by definition the
kernel of the projection
$\GRT(\mathbb{Q})\to\GRT(\mathbb{Q})/\GRT^{\mathrm{fin}}$ that
quotients out Casimir rescalings and simple-pole shifts; the seven
explicit cocycles all factor through it. The first non-trivial
bar-degree $d^\star(\Phi_{1i})$ recorded in
Table~\ref{tab:hdm-cocycle-record} is read off the $\fgrt$-weight
expansion of $\log\Phi_{1i}$:
\begin{itemize}
\item $\Phi_{12}$ is the $R$-twist of the deconcatenation coproduct,
$\Delta^R = R\Delta R^{-1}$ with $R = 1+\hbar r_\cA(z)+O(\hbar^2)$. At
binary degree, conjugation acts by the inner adjoint, which is the
identity on the symmetric Casimir; the order-$\hbar$ correction is
$[r_\cA(z),\Delta]$ and lies in the kernel of the
$\Sigma_2$-coinvariant projection because $r_\cA(z)$ is itself
symmetric. Hence $\Phi_{12}$ acts as the identity on the symmetric
binary residue, and $d^\star(\Phi_{12})=2$ records the ordered slot
of first-non-triviality before symmetrisation.
\item $\Phi_{13}$ is the composite of the Casimir rescaling
$\Omega_{\mathrm{tr}}\mapsto\Omega/(2h^\vee)$
(Vol~II, Proposition~\ref{V2-prop:grtfin-casimir}) and the
Sugawara level shift $k\mapsto k+h^\vee$. Both belong to
$\GRT^{\mathrm{fin}}$; both act non-trivially at bar-degree~$2$ on
class~$\mathbf{L}$ but trivially on classes $\mathbf{G},\mathbf{C},\mathbf{M}$
(where $h^\vee=0$ structurally, the residue is $0$, or the algebra
is non-Lie respectively).
\item $\Phi_{14}$ is the Laplace transform that converts the
$\lambda$-bracket presentation into a generating function in $1/z$.
The Laplace transform is an isomorphism of formal vector spaces
that fixes the binary symmetric tensor; hence $\Phi_{14}$ acts as
the identity on the bar-degree-$2$ slice.
\item $\Phi_{15}$ is the classical limit of the Yangian
$R$-matrix $R_{\mathrm{Yang}}(u)-1=\hbar P/u+O(\hbar^2)$. The
leading coefficient $P/u$ is the rational $r$-matrix; the
$O(\hbar^2)$ tail carries no MZV at weight~$2$ (because the
weight-$2$ generator of $\fgrt$ vanishes by the pentagon axiom).
Hence $\Phi_{15}$ is identity at bar-degree~$2$ modulo
$\GRT^{\mathrm{fin}}$.
\item $\Phi_{16}$ inverts the Sugawara denominator $1/(k+h^\vee)$
to convert the KZ residue $\Omega/((k+h^\vee)z)$ into the Gaudin
generator $\Omega_{ij}/(z_i-z_j)$. This is a multiplicative
rescaling by the level-shift factor; on class~$\mathbf{L}$ it acts
non-trivially (inverse of $\Phi_{13}$ on the level-shift component);
on classes $\mathbf{G},\mathbf{C},\mathbf{M}$ it acts trivially.
\item $\Phi_{17}$ on classes $\mathbf{G},\mathbf{L}$ is the
simple-pole truncation of the bar-degree-$3$ component, which
restricts to the identity on the binary slice $(d=2)$. On
class~$\mathbf{M}$ the correct entry is $F_7'$ (top $A_\infty$
operation $m_3$), and the cocycle $\Phi_{17}'$ acquires MZV-valued
corrections at every weight $w\geq 3$
(Vol~II, Proposition~\ref{V2-prop:f7-split}); these corrections are
supported entirely at bar-degree~$\geq 3$ and invisible to the
binary residue.
\end{itemize}

\smallskip\noindent
\emph{(ii) Degree-$2$ collapse.} The pro-unipotent radical
$\GRT^{(\geq 3)}$ is by definition the closed normal subgroup of
$\GRT(\mathbb{Q})$ generated by associators whose
$\fgrt$-logarithm has weight $\geq 3$. The quotient
$\GRT^{(2)}=\GRT(\mathbb{Q})/\GRT^{(\geq 3)}$ is the abelianisation
of the bar-degree-$2$ action on the symmetric Casimir, which is the
trivial group: at bar-degree~$2$, the only $\fgrt$-element is the
weight-$2$ component, and that component vanishes identically by the
pentagon axiom (Drinfeld 1991; Bar-Natan 1998). The finite inner
gauge $\GRT^{\mathrm{fin}}$ acts on $r_\cA(z)$ via Casimir
rescalings and level shifts, both of which fix the underlying
$\fg$-equivariant tensor up to scalar; the scalar agrees lane-by-lane
because the seven faces are stated in their canonical chain models
with matched normalisations
(Proposition~\ref{prop:hdm-averaging-compatibility}, lane-by-lane
identity). Hence $\mu_2(\Phi_{1i})=1$ for every $i$, and the seven
faces are pointwise equal on the binary slice.

\smallskip\noindent
\emph{(iii) MZV invisibility.} This is exactly Vol~II,
Theorem~\ref{V2-thm:f1-f4-injection} (F1$\leftrightarrow$F4 injection),
transported to the Vol~I bar-intrinsic perspective: the binary-degree
projection of the $\fgrt$-action factors through $\GRT^{(2)}$, and
$\Phi_{\KZ}-\Phi_{\rat}$ has zero weight-$2$ component by the
pentagon axiom. The Gerstenhaber-bracket computation
$\Phi\cdot r_\cA(z)
=
r_\cA(z)+\sum_{w\geq 2}\hbar^w[\phi_w,r_\cA(z)]_{\mathrm{Ger}}$
shows that the MZV tail at weight $w\geq 3$ acts only at bar-degree
$\geq w+1\geq 4$, hence vanishes on the binary $r$.

\smallskip\noindent
\emph{(iv) Class-by-class refinement.} The class-$\mathbf{G}$ residue
$r^{\mathrm{Heis}}_k(z)=k\,J\otimes J/z$ is a single tensor with no Lie-bracket
structure; Casimir rescalings act by scalar multiplication, which is
absorbed into the level normalisation. The class-$\mathbf{C}$ residue
vanishes identically ($r_{\beta\gamma}=0$); any cocycle action is
vacuous. The class-$\mathbf{M}$ residue
$r^{\mathrm{Vir}}_c(z)=(c/2)/z^3+2T/z$ is fixed by every Casimir
rescaling (no $\fg$-Casimir is involved) and by every Sugawara level
shift (no level parameter); the only non-trivial cocycle action is
$\Phi_{17}'$ at bar-degree~$\geq 3$, which projects to zero on the
binary slice. The class-$\mathbf{L}$ residue is the only one on which
$\Phi_{13}$ (Casimir rescaling $+$ level shift) and $\Phi_{16}$
(Sugawara-denominator inversion) act non-trivially at the binary
level. The composite traversal $\Phi_{16}\circ\Phi_{13}$ along the
path $F_1\to F_3\to F_6$ (or via $F_7$) yields the identity on the
binary symmetric tensor, because the inverse level factor
$(k+h^\vee)^{-1}$ introduced by $\Phi_{13}$ is annulled by the
Sugawara prefactor $(k+h^\vee)$ supplied by $\Phi_{16}$ along the
Gaudin chain
\eqref{eq:hdm-face-7}. The lane-to-lane chained equality
$k\,\Omega_{\mathrm{tr}}/z\leftrightarrow\Omega/((k+h^\vee)z)$ is
therefore realised as a $\GRT^{\mathrm{fin}}$ gauge with no residual
binary content (Vol~II,
Proposition~\ref{V2-prop:grtfin-casimir}).
\end{proof}

\begin{remark}[$\GRT$ as a torsor structure: pro-algebraic
vs $\infty$-groupoid]
\label{rem:hdm-grt-torsor-ambient}
\index{GRT torsor!ambient (pro-algebraic vs $\infty$-groupoid)}
The $\GRT(\mathbb{Q})$-action on $\Face(\cA)$ in
Vol~II's Proposition~\ref{V2-prop:grt-action-face} is a free and
transitive action of the \emph{pro-algebraic group} $\GRT_1$ over
$\mathbb{Q}$ (equivalently, $\GRT(\mathbb{Q})$ as a pro-unipotent
group with $\fgrt$ as its pro-nilpotent Lie algebra). At the
$(\infty,1)$-categorical level the action lifts to a
free-and-transitive $\infty$-action of the $\infty$-group
$B\GRT_1$ on the $\infty$-groupoid of presentations, but the
ambient framework throughout this monograph is the strict
pro-algebraic one: the cocycles $\Phi_{1i}$ are honest elements of
$\GRT(\mathbb{Q})$, not weak equivalences in an
$\infty$-groupoid. The chain-level vs $\infty$-categorical
disambiguation is taken up in
Section~\ref{sec:e1-koszul-duality-categorical-mode-key}; the
Vol~I face-cocycle record of
Theorem~\ref{thm:hdm-grt-cocycle-record} is stated at the
chain-level (strict $\MC$-gauge equivalence, per
Definition~\ref{V2-def:face-space} of Vol~II), and is
$\infty$-categorically lifted by the formality theorems of
Tamarkin~\cite{Tamarkin00} and
Kontsevich~\cite{Kon03}.
\end{remark}

\begin{remark}[Genus-$1$ activation: where the binary collapse fails]
\label{rem:hdm-grt-genus1-activation}
\index{GRT torsor!genus-1 activation|textbf}
\index{seven faces!genus-1 splitting from $\GRT^{(\geq 3)}$|textbf}
The degree-$2$ collapse of
Theorem~\ref{thm:hdm-grt-cocycle-record}(ii) is a genus-$0$
phenomenon. At genus~$1$ the bar propagator
$\eta_{12}^{E_\tau}=\zeta_\tau(z_1-z_2)\,d(z_1-z_2)$
\textup{(}Definition~\ref{def:g1sf-bar-propagator}\textup{)}
acquires the quasi-period datum
$(\eta_1,\eta_\tau)$, and the seven faces of
Chapter~\ref{ch:genus1-seven-faces} split into seven genuinely
distinct objects. From the cocycle perspective, the genus-$1$
splitting is the activation of $\GRT^{(\geq 3)}$-data: the $B$-cycle
monodromy of the Weierstrass zeta function couples the binary
residue to higher-bar-degree corrections at $w=3,5,7,\dots$, where
the MZV tail enters as $\Phi_{\KZ}-\Phi_{\rat}\not=0$ on the
elliptic completion.

The structural reading is:
\begin{itemize}
\item At genus~$0$, $\GRT^{(\geq 3)}$ acts trivially on the
family-normalized residue $r_\cA(z)$, the
seven faces collapse, and the master equation
\eqref{eq:hdm-master-equation} holds pointwise.
\item At genus~$1$, $\GRT^{(\geq 3)}$ acts non-trivially through
the $B$-cycle monodromy, the seven elliptic faces $F_1^{E_\tau},
\dots,F_7^{E_\tau}$ separate, and the genus-$1$ master theorem
\textup{(}Chapter~\ref{ch:genus1-seven-faces},
Theorem~\textup{\ref{thm:g1sf-master}}, when proved\textup{)}
records seven distinct elliptic specialisations linked by
quasi-periodic gauge.
\item In the degeneration $\tau\to i\infty$, the $B$-cycle
monodromy collapses and the genus-$1$ splitting reduces uniformly
to the genus-$0$ collapse \textup{(}Chapter~\ref{ch:genus1-seven-faces},
\S\textup{\ref{sec:g1sf-degeneration}}\textup{)}.
\end{itemize}
This explains the geometric content of the seven-face splitting:
it is the binary-vs-higher-bar-degree dichotomy of the $\fgrt$-action,
made visible by the curve-genus $g$ through the monodromy of the
bar propagator on $\overline{\mathcal{M}}_{g,n}$.
\end{remark}

\begin{remark}[Cross-volume bridge: cocycle record and Vol~III
$\Phi_d$ correspondence]
\label{rem:hdm-grt-cocycle-xvol}
\index{cross-volume bridge!cocycle record}
On the Vol~III side, the analogous cocycle record lives on the
seven CY faces of $r_{CY}(z) = \Res^{\mathrm{coll}}_{0,2}
(\Theta_{A_\cC})$ \textup{(}Vol~III,
Theorem~\textup{\ref{V3-thm:seven-faces-of-r-cy}}\textup{)}; the
class-by-class refinement of
Theorem~\ref{thm:hdm-grt-cocycle-record}(iv) extends to the CY
setting via the dimension-stratified correspondence
$\{\Phi_d\}_{d\geq 1}$ on the Vol~III holographic CY datum
\textup{(}Vol~III, Chapter~``\texttt{cy\_holographic\_datum\_master}''\textup{)},
under which the class-$\mathbf{L}$ Casimir-rescaling cocycle
$\Phi_{13}$ corresponds to the toric-CY3 cocycle controlling the
COHA-to-Yangian comparison, and the class-$\mathbf{M}$ MZV-tail
cocycle $\Phi_{17}'$ corresponds to the $K3\times E$ Borcherds
cocycle controlling the primitive BKM denominator
\textup{(}Vol~III, $\Delta_5$; the weight-$10$ square is
$\Phi_{10}^{\mathrm{un}}=\Delta_5^2$\textup{)}. The
binary-degree collapse of the $\GRT^{\mathrm{fin}}$ action is
therefore a single cross-volume statement: the seven faces of
$r_\cA(z)$ (Vol~I), the GRT-orbit representatives of $\Face(\cA)$
(Vol~II), and the seven CY faces of $r_{CY}(z)$ (Vol~III) all
collapse on the binary slice through the same finite inner gauge
group, which is the $\fgrt$-weight-$2$-trivial quotient
$\GRT^{(2)}=\{1\}$.
\end{remark}

\begin{remark}[Igusa $\Delta_5$ as the gauge-class invariant of the Hall--Drinfeld double at $K3 \times E$]
\label{rem:holo-master-igusa}
\index{Igusa cusp form $\Delta_5$!gauge-class invariant of $U_\cA$}
\index{Hall--Drinfeld double!$H^2$ classification at $K3 \times E$}
The binary-degree collapse asserted above leaves a residual
cross-volume classification: the $\mathbb{Z}/2 \times K(1)$-invariants
of the degree-$2$ gauge cohomology of the Drinfeld double
$U_\cA$ at the $K3 \times E$ fibration, where $\mathbb{Z}/2$ is the
Fourier--Mukai involution acting on the K3 factor and $K(1)$ is the
elliptic weight-one modular action on the $E$ factor. The classical
Kac--Peterson--Lepowsky statement in the Borcherds--Gritsenko--Nikulin
setup is
\[
 H^2\bigl(\mathfrak{g}_{\Delta_5}\bigr)^{\mathbb{Z}/2,\, K(1)}
 \;=\;
 \mathbb{C}\cdot \Delta_5,
\]
with $\mathfrak{g}_{\Delta_5}$ the Borcherds--Kac--Moody Lie algebra
whose denominator is the Igusa cusp form $\Delta_5$ of weight~$5$ on
$\mathrm{Sp}_4(\mathbb{Z})$
(Gritsenko--Nikulin~\cite{GritsenkoNikulin18},
Borcherds~\cite{Borcherds95} on singular-theta pushforward). The
present reading of this one-dimensional $H^2$ is that \emph{the
Igusa form is the gauge-class invariant of the Hall--Drinfeld
double}: the class-$\mathbf{M}$ MZV-tail cocycle $\Phi_{17}'$
(Theorem~\ref{thm:hdm-grt-cocycle-record}(iv)) corresponds at
 $K3 \times E$ to the BKM denominator cocycle of Vol~III
(the primitive $\Delta_5$ denominator, with
$\Phi_{10}^{\mathrm{un}}=\Delta_5^2$ only after squaring), and the
one-dimensionality of the $(\mathbb{Z}/2, K(1))$-invariant subspace
forces the gauge class of the Drinfeld double $U_\cA$ at
$K3 \times E$ to be a scalar multiple of $\Delta_5$ on the
boundary-linear, chirally Koszul locus in
Part~\ref{part:standard-landscape}. This is the scalar rigidity
statement underlying the six-slot recovery in the BKM family: the
$K3 \times E$ specialisation of $U_\cA$ admits no further gauge
freedom once the weight is fixed, and the weight is fixed by the
trace identity of Remark~\ref{rem:hdm-anchors}(W2) evaluated
on the $K3 \times E$ BRST ghost sector.
\end{remark}

\subsection{Recovery of the six-slot datum from the Drinfeld double}
\label{subsec:hdm-u-a-recoverability}

The seven-way master theorem verifies one of the six slots of
$\mathcal{H}(T)$: the spectral $R$-matrix~$r_\cA(z)$. The remaining five
slots are not independent invariants; they are co-projections of the
same boundary-bulk algebra~$U_\cA = \cA \bowtie \cA^{!}_\infty$
introduced in Section~\ref{sec:hdm-thesis}. Here boundary/bulk still
means boundary object~$\cA$, bulk slot
$\cC=\cZ^{\mathrm{der}}_{\mathrm{ch}}(\cA)$, restriction map
$\rho_{\partial}\colon \cC\to\cA$ coming from
$\mathrm{Ass}^{\mathrm{ch}}\hookrightarrow\mathrm{SC}^{\mathrm{ch,top}}$,
and witnesses $(r_\cA(z),\Theta_\cA,\nabla^{\mathrm{hol}})$. Off the
modular-Koszul / boundary-linear / principal DS comparison surfaces,
the strict edge $\cA^{!}_\infty\to\cA^!$ and the nonlinear quartic
contact class prevent a strict global closure. The present subsection
separates three upgrades that must not be conflated: existence of the
universal MC element, strict Verdier/Koszul rectification, and the
bulk/centre/line identification. Only after those three upgrades have
been stated separately does the umbrella recoverability proposition
make sense.

\begin{proposition}[Upgrade~(U1): existence and universality of
 $\Theta_\cA$; \ClaimStatusConditional]
\label{prop:hdm-upgrade-theta}
Let $\cA$ be a modular Koszul chiral algebra on a smooth curve~$X$.
In the chain-level lane there is a degree-$1$ Maurer--Cartan element
\[
\Theta_\cA \in
\mathfrak{g}^{E_1,1}_\cA
:=
\operatorname{Conv}\bigl(B^{\mathrm{ord}}(\cA),\cA\bigr)^1
\]
satisfying
\[
D\cdot\Theta_\cA + \tfrac{1}{2}[\Theta_\cA,\Theta_\cA]=0.
\]
Its averaged modular image is the bar-intrinsic class
$D_\cA-d_0 \in
\mathrm{MC}(\Defcyc(\cA)\widehat{\otimes}\mathcal G_{\mathrm{mod}})$.
On the generic cyclic-rigid locus its gauge class depends only on
$\cA$ and not on a choice of generators or presentation. This upgrade
proves slot~\textup{(5)} of~$\mathcal H(T)$ and, by projection,
constructs the pinned shadow connection
$\nabla^{\mathrm{hol}}=d-\operatorname{Sh}(\Theta_\cA)$; it says
nothing yet about strict rectification or about the bulk slot~$\cC$.
\end{proposition}

\begin{proof}
The modular class is exactly Theorem~\ref{thm:mc2-bar-intrinsic}. The
ordered lift is the $E_1$-Maurer--Cartan element whose averaging gives
the modular class \textup{(}Definition~\textup{\ref{def:e1-modular-convolution}}
and Proposition~\textup{\ref{prop:symmetric-descent}}\textup{)}.
Gauge uniqueness on the generic cyclic-rigid locus is
Theorem~\ref{thm:cyclic-rigidity-generic}. This is the slot~\textup{(5)}
upgrade only.
\end{proof}

\begin{proposition}[Upgrade~(U2): strict Verdier/Koszul rectification;
 \ClaimStatusConditional]
\label{prop:hdm-upgrade-rectification}
Let $\cA$ be a modular Koszul chiral algebra on a smooth curve~$X$.
On the modular Koszul locus the Verdier-dual factorization companion
$\cA^{!}_\infty$ rectifies to the strict chiral Koszul dual
$\cA^!$. Off the Koszul locus one retains only the coderived,
weight-completed factorization companion together with the bar--cobar
correspondence; the strict edge $\cA^{!}_\infty \to \cA^!$ is
conjectural there. The boundary obstruction is explicit: for the
admissible affine algebra $L_{-1/2}(\mathfrak{sl}_2)$ the bar class
$\eta_4$ survives in bar cohomology and blocks the quasi-isomorphism
$\Omega^{\mathrm{ch}}(\bar B^{\mathrm{ch}}(\cA))\to\cA$, hence blocks
the strict rectification.
\end{proposition}

\begin{proof}
The strict statement is the Koszul-locus content of
Theorem~\ref{thm:bar-cobar-inversion-qi} together with the
Verdier/Koszul comparison of
Theorem~\ref{thm:frontier-protected-bulk-antiinvolution}. Off the
Koszul locus the surviving coderived statement is precisely what
Theorem~\ref{thm:bar-cobar-inversion-qi} leaves in place. The explicit
obstruction class $\eta_4$ is the content of
Example~\ref{ex:admissible-sl2-failure} and will reappear below in
Remark~\ref{rem:hdm-boundary-linear-vs-global-triangle}. This is the
slot~\textup{(2)} upgrade only.
\end{proof}

\begin{proposition}[Upgrade~(U3): the bulk slot is the derived chiral
 centre, with line-side comparison on named lanes;
 \ClaimStatusConditional]
\label{prop:hdm-upgrade-bulk-slot}
For every curved $A_\infty$-chiral algebra~$\cA$, the bulk slot pinned
at the chapter top is
\[
\cC
\;=\;
\cZ^{\mathrm{der}}_{\mathrm{ch}}(\cA)
\;=\;
\mathcal{C}^{\bullet}_{\mathrm{ch}}(\cA,\cA),
\]
and the boundary restriction map
$\rho_{\partial}\colon \cC \to \cA$ is induced by the Swiss-cheese
operadic inclusion
$\mathrm{Ass}^{\mathrm{ch}}\hookrightarrow\mathrm{SC}^{\mathrm{ch,top}}$
and witnessed by the same three pieces of homotopy data
$(r_\cA(z),\Theta_\cA,\nabla^{\mathrm{hol}})$. On chirally Koszul
boundary-linear families of classes $\mathsf{G},\mathsf{L},\mathsf{C}$
and on the principal Drinfeld--Sokolov lane at non-critical level,
Volume~II identifies the same slot with the line model
$\cC_{\mathrm{line}}\simeq \cA^!_{\mathrm{line}}\text{-}\mathsf{mod}$.
Outside those lanes the bulk object~$\cC$ remains defined, but the
global triangle comparing it to the line model is conjectural because
the strict edge $\cA^{!}_\infty\to\cA^!$ can fail and, on nonlinear
class~$\mathsf{M}$, the quartic contact class $10/[c(5c+22)]$ blocks
the boundary-linear reduction.
\end{proposition}

\begin{proof}
The open/closed statement
$\cZ^{\mathrm{der}}_{\mathrm{ch}}(\cA)
\simeq\mathcal{C}^{\bullet}_{\mathrm{ch}}(\cA,\cA)$ is
Theorem~\ref{thm:thqg-swiss-cheese}. The line-side comparison is the
precise content of the Vol.~II theorems quoted in
Proposition~\ref{prop:hdm-four-recovery-directions}. The obstruction
boundary is not rhetorical: away from the named lanes the strict edge
$\cA^{!}_\infty\to\cA^!$ fails to glue, and on nonlinear
class~$\mathsf{M}$ the quartic contact class
$10/[c(5c+22)]$ blocks the boundary-linear reduction. This is the
slot~\textup{(3)} upgrade only.
\end{proof}

\begin{proposition}[Recovery of $\mathcal{H}(T)$ from $U_\cA$;
 \ClaimStatusConjectured]
\label{prop:hdm-u-a-recoverability}
\index{Drinfeld double!boundary-bulk recoverability|textbf}
\index{holographic modular Koszul datum!recovery from $U_\cA$|textbf}
Let $\cA$ be a modular Koszul chiral algebra on a smooth projective
curve~$X$, and let $U_\cA := \cA \bowtie \cA^{!}_\infty$ be the
Drinfeld double of~$\cA$ and its Verdier-dual factorization companion.
The six components of the holographic modular Koszul datum
$\mathcal{H}(T) = (\cA, \cA^!, \cC, r_\cA(z), \Theta_\cA, \nabla^{\mathrm{hol}})$
(Definition~\ref{def:holographic-modular-koszul-datum}) are
projections of~$U_\cA$ as follows.
\begin{enumerate}[label=\textup{(R\arabic*)}]
\item \emph{Boundary factor.} The first tensor factor of~$U_\cA$ is
 the boundary chiral algebra~$\cA$.
\item \emph{Dual factor.} The second tensor factor of~$U_\cA$ is the
 Verdier-dual factorization companion~$\cA^{!}_\infty$; when
 Upgrade~\textup{(U2)} holds, its strict cohomological shadow is the
 chiral Koszul dual~$\cA^!$.
\item \emph{Bulk/line slot.} The same boundary-bulk package determines
 the pinned bulk algebra
 $\cC=\cZ^{\mathrm{der}}_{\mathrm{ch}}(\cA)
 \simeq \mathcal{C}^{\bullet}_{\mathrm{ch}}(\cA,\cA)$, and on the
 chirally Koszul line side it presents the line category by
 $\cC_{\mathrm{line}} \simeq
 \cA^!_{\mathrm{line}}\text{-}\mathsf{mod}$ once the compact-generator
 comparison is available. This is the third slot~$\cC$ of
 $\mathcal{H}(T)$; the base curve remains ambient data.
\item \emph{Quasitriangular $R$-matrix.} The spectral $R$-matrix of
 the quasitriangular structure on~$U_\cA$ equals the collision
 residue~$r_\cA(z) = \operatorname{Res}^{\mathrm{coll}}_{0,2}(\Theta_\cA)$
 by the seven-way master
 theorem (Theorem~\ref{thm:hdm-seven-way-master}).
\item \emph{Universal MC element.} The bar-intrinsic differential on
 the convolution dg~Lie algebra
 $\mathfrak{g}_\cA^{\mathrm{mod}}$ associated to~$U_\cA$ produces the
 universal Maurer--Cartan element~$\Theta_\cA$ by
 $\Theta_\cA = D_\cA - d_0$
 (Theorem~\ref{thm:mc2-bar-intrinsic}).
\item \emph{Shadow connection.} The modular shadow connection
 $\nabla^{\mathrm{hol}}_{g,n} = d - \operatorname{Sh}_{g,n}(\Theta_\cA)$
 is the flat connection produced by projecting~$\Theta_\cA$ to the
 $(g,n)$-sector, and therefore depends only on~$U_\cA$.
\end{enumerate}
Consequently, the heuristic claim that $\mathcal{H}(T)$ is recoverable
from~$U_\cA$ alone is the conjunction of Upgrades~\textup{(U1)}--%
\textup{(U3)} together with the seven-way master theorem. The six
slots carry no information beyond the Drinfeld double only where those
upgrades and comparison maps are simultaneously available.
\end{proposition}

\begin{remark}[Recovery scope]
\label{rem:hdm-u-a-recoverability-scope}
Items~(R1) and~(R4) are unconditional for modular Koszul~$\cA$:
the factor structure of~$U_\cA$ and the seven-way master theorem
supply them directly. Item~(R5) is exactly Upgrade~(U1), hence
 unconditional as a construction of~$\Theta_\cA$ and of
 $\nabla^{\mathrm{hol}}$; this proves only the MC/connection slot and
 does not yet rectify~$\cA^{!}_\infty$ or glue the bulk triangle.
 Item~(R2) is exactly Upgrade~(U2). The Koszul-locus qualifier is
 sharp: outside it the bar complex no longer resolves the strict dual,
 and the admissible class $\eta_4$ of
 Example~\ref{ex:admissible-sl2-failure} blocks the edge
 $\cA^{!}_\infty\to\cA^!$. Item~(R3) is exactly Upgrade~(U3). The
 open/closed bulk algebra
 $\cC=\cZ^{\mathrm{der}}_{\mathrm{ch}}(\cA)
 \simeq \mathcal{C}^{\bullet}_{\mathrm{ch}}(\cA,\cA)$ is proved for
 every~$\cA$, but the line-side presentation and the bulk/line
 identification are proved only on the boundary-linear / DS comparison
 lanes; off those lanes the global triangle fails first at the same
 strict edge and, on nonlinear class~$\mathsf{M}$, at the quartic
 contact class. The overall recoverability statement is therefore
 tagged \ClaimStatusHeuristic: the slots are pinned and partially
 linked, but the global assembly into a single invariant of~$U_\cA$
 remains conjectural beyond the named comparison surfaces.
\end{remark}

\subsection{Universal property of $U_\cA$}
\label{subsec:hdm-u-a-universal-property}

The six recovery clauses of
Proposition~\ref{prop:hdm-u-a-recoverability} name~$U_\cA$ and list
what it projects to. They do not yet answer what~$U_\cA$ \emph{is}.
The $\bowtie$ symbol is load-bearing: $\cA$ is an honest chiral
algebra, while $\cA^{!}_\infty$ is the coderived, weight-completed
factorization companion, strictly multiplicative only after
Upgrade~\textup{(U2)}. A strict matched-pair Hopf product is therefore
available only on the modular Koszul locus; off that locus the product
$\bowtie$ must be read as the twisted tensor product
$\cA \otimes^{\tau}_{\cA^{!}_\infty}$ whose twisting cocycle is the
universal Maurer--Cartan element~$\Theta_\cA$. The universal property
below states~$U_\cA$ as a representing object, independently of the
two possible readings of~$\bowtie$.

Let $\mathrm{ChirAlg}^{\mathrm{MK}}_X$ be the $(\infty,1)$-category of
modular Koszul chiral algebras on~$X$, and let
$\mathrm{Pair}^{\mathrm{Kos}}_\cA(T)$ denote, for $T \in
\mathrm{ChirAlg}^{\mathrm{MK}}_X$, the space of \emph{Koszul pairs
with~$\cA$ over~$T$}: pairs $(P, q)$ of a dg chiral coalgebra~$P$ over
the factorization base and a twisting morphism $q\colon P \to \cA$
\textup{(}Definition~\textup{\ref{def:twisting-morphism}}\textup{)},
together with a twisting cocycle $\alpha\colon P \to T$ making
$(\cA, T, P, q, \alpha)$ into a compatible pair in the sense that the
induced map
$P \to \cA \otimes^{\tau} T$ is a bar-cobar inverse to the universal
twisting. The assignment $T \mapsto \mathrm{Pair}^{\mathrm{Kos}}_\cA(T)$
extends to a 2-functor
$\mathrm{Pair}^{\mathrm{Kos}}_\cA\colon
\mathrm{ChirAlg}^{\mathrm{MK}}_X \to \mathrm{Cat}_{(\infty,1)}$
by pullback along algebra maps.

\begin{theorem}[Universal property of $U_\cA$;
\ClaimStatusConditional]
\label{thm:hdm-u-a-universal-property}
\index{Drinfeld double!universal property of $U_\cA$|textbf}
\index{universal property!boundary-bulk algebra $U_\cA$|textbf}
Let $\cA$ be a modular Koszul chiral algebra on a smooth projective
curve~$X$. The boundary-bulk algebra
$U_\cA := \cA \bowtie \cA^{!}_\infty$
represents the 2-functor of Koszul pairs with~$\cA$:
\begin{equation}\label{eq:hdm-u-a-representable}
\mathrm{Pair}^{\mathrm{Kos}}_\cA(T)
\;\simeq\;
\mathrm{Map}_{\mathrm{ChirAlg}^{\mathrm{MK}}_X}(U_\cA,\,T)
\qquad\text{for every } T \in \mathrm{ChirAlg}^{\mathrm{MK}}_X.
\end{equation}
Equivalently, $U_\cA$ is the initial object, in the
$(\infty,1)$-category of twisted tensor Koszul pairs under~$\cA$,
whose underlying twisting morphism is the universal one~$\pi_\cA$. On
the modular Koszul locus this initial object is computed as the
relative derived tensor product over the bar coalgebra,
\begin{equation}\label{eq:hdm-u-a-relative-tensor}
U_\cA
\;\simeq\;
\cA \,\otimes^{\mathbb L}_{\barBch(\cA)}\, \cA^{!}_\infty,
\end{equation}
and is quasitriangular, with universal spectral $R$-matrix equal to
the collision residue $r_\cA(z) = \mathrm{Res}^{\mathrm{coll}}_{0,2}
(\Theta_\cA)$. The unit $\eta\colon \underline{k} \to U_\cA$ is the
joint unit $1_\cA \otimes 1_{\cA^{!}_\infty}$, the counit
$\epsilon\colon U_\cA \to \underline{k}$ is the product of
augmentations, and the universal twisting morphism $\pi_\cA\colon
\barBch(\cA) \to \cA$ factors as the composition of the inclusion
$\barBch(\cA) \hookrightarrow U_\cA$ followed by the boundary
projection $U_\cA \twoheadrightarrow \cA$.
\end{theorem}

\begin{proof}
On the modular Koszul locus Upgrade~\textup{(U2)}
(Proposition~\ref{prop:hdm-upgrade-rectification}) strictifies the
second factor to~$\cA^!$, turning $\bowtie$ into a matched-pair Hopf
product. In that case the representability
statement~\eqref{eq:hdm-u-a-representable} is the chiral upgrade of
the Cartan--Eilenberg--Positselski universal coalgebra property: the
twisting morphism functor $T \mapsto \mathrm{Tw}(\barBch(\cA), T)$ is
corepresented on chiral coalgebras by the Koszul dual coalgebra
$\cA^{\mathrm{i}}:=H^\star\bigl(\barBch(\cA)\bigr)$. The strict
Koszul dual algebra is obtained only after Verdier duality,
$\cA^!=(\cA^{\mathrm{i}})^\vee$, when Upgrade~\textup{(U2)} holds.
The enrichment to Koszul pairs is then corepresented by the tensor
product of the boundary factor~$\cA$ with the Verdier-dual
factor~$\cA^!$ \textup{(}Theorem~\textup{\ref{thm:bar-cobar-adjunction}}
in the bar--cobar inversion lane and
Theorem~\textup{\ref{thm:frontier-protected-bulk-antiinvolution}}
in the Verdier/Koszul lane\textup{)}. The relative tensor-product
presentation~\eqref{eq:hdm-u-a-relative-tensor} is then the
$\infty$-categorical formula for the pushout in the 2-functor
$\mathrm{Pair}^{\mathrm{Kos}}_\cA$, computed using Lurie's
$\mathrm{HA}.5.5$ machinery for relative tensor products in stable
$\infty$-categories over the bar coalgebra~$\barBch(\cA)$.

Off the modular Koszul locus the strict edge
$\cA^{!}_\infty \to \cA^!$ is conjectural
\textup{(}Proposition~\textup{\ref{prop:hdm-upgrade-rectification}}\textup{)},
and the symbol~$\bowtie$ must then be read as the bar-twisted tensor
product. The initial-object statement nevertheless holds at the
coderived level: the universal twisting element~$\Theta_\cA$
(Upgrade~\textup{(U1)},
Proposition~\ref{prop:hdm-upgrade-theta}) provides a canonical twisting
cocycle for $\cA \otimes^\tau \cA^{!}_\infty$, and the universal
property of the derived tensor product under~$\cA$ factors every
Koszul pair $(P, q, \alpha)$ through~$U_\cA$ at the level of
coderived factorization $\infty$-categories. Quasitriangularity and
the identification of the universal $R$-matrix with the collision
residue is the seven-way master
theorem~\textup{(}Theorem~\textup{\ref{thm:hdm-seven-way-master}}\textup{)}
applied to slot~\textup{(4)} of~\eqref{eq:hdm-u-a-representable}.
\end{proof}

\begin{corollary}[Six projections as forgetful functors;
\ClaimStatusConditional]
\label{cor:hdm-u-a-six-projections}
\index{holographic modular Koszul datum!six projections from $U_\cA$|textbf}
The six slots of
$\mathcal H(T) = (\cA, \cA^!, \cC, r_\cA(z),
\Theta_\cA, \nabla^{\mathrm{hol}})$ are the images of the representing
object~$U_\cA$ under the following six forgetful functors, each of
which is a projection of the representability
isomorphism~\eqref{eq:hdm-u-a-representable}:
\begin{enumerate}[label=\textup{(P\arabic*)}]
\item \emph{Boundary projection} $U_\cA \mapsto U_\cA \otimes_{\cA^{!}_\infty}
 \underline{k} = \cA$. Forgets the dual factor. Recovers slot~\textup{(1)}.
\item \emph{Dual projection} $U_\cA \mapsto \underline{k} \otimes_\cA
 U_\cA = \cA^{!}_\infty$. Forgets the boundary factor. Strict
 edge to~$\cA^!$ on the modular Koszul locus is Upgrade~\textup{(U2)}.
 Recovers slot~\textup{(2)}.
\item \emph{Derived-centre projection} $U_\cA \mapsto
 \mathrm{End}^{\bullet}_{U_\cA\text{-}\mathrm{mod}}(\cA) =
 \cC^{\bullet}_{\mathrm{ch}}(\cA, \cA) = \cC$. The derived
 centre of~$U_\cA$ acting on its boundary representation is the chiral
 Hochschild cochain complex. Recovers slot~\textup{(3)} via
 Upgrade~\textup{(U3)}.
\item \emph{Universal $R$-matrix projection}
 $U_\cA \mapsto R_{U_\cA} \in
 U_\cA \otimes U_\cA[\![z^{-1}]\!]$, followed by restriction to the
 binary collision stratum, produces~$r_\cA(z)$. Recovers
 slot~\textup{(4)}; this is exactly Theorem~\ref{thm:hdm-seven-way-master}.
\item \emph{Bar-differential projection} $U_\cA \mapsto \Theta_\cA \in
 \mathrm{MC}(\mathfrak g^{\mathrm{mod}}_\cA)$, where~$\Theta_\cA$ is
 the universal twisting cocycle defining the twisted tensor product in
\eqref{eq:hdm-u-a-relative-tensor}. Recovers slot~\textup{(5)}.
\item \emph{Shadow-connection projection} $U_\cA \mapsto
 \nabla^{\mathrm{hol}} = d - \mathrm{Sh}(\Theta_\cA)$, obtained by
 pushing forward the bar-differential projection along the shadow map.
 Recovers slot~\textup{(6)}.
\end{enumerate}
In particular, the holographic datum $\mathcal H(T)$ carries no
information beyond~$U_\cA$ on the modular-Koszul / boundary-linear /
principal DS surfaces: all six slots are values of the representable
functor~\eqref{eq:hdm-u-a-representable} under the six forgetful
functors~\textup{(P1)}--\textup{(P6)}.
\end{corollary}

\begin{proof}
Projections~\textup{(P1)}, \textup{(P2)} are the two tensor-factor
projections of the twisted tensor product
$\cA \otimes^\tau \cA^{!}_\infty$, computed by tensoring against the
opposite factor's augmentation module. Projection~\textup{(P3)} is
the $\mathrm{HA}.5.5$ formula for the derived endomorphism algebra of
the boundary module, together with
Theorem~\ref{thm:thqg-swiss-cheese} identifying it with chiral
Hochschild cochains. Projection~\textup{(P4)} is the content of
Theorem~\ref{thm:hdm-seven-way-master}: the universal $R$-matrix of
the quasitriangular structure on~$U_\cA$ equals the collision
residue. Projection~\textup{(P5)} is the defining datum of
$U_\cA$ as a twisted tensor product: the twisting cocycle is the
universal Maurer--Cartan element~$\Theta_\cA$
(Theorem~\ref{thm:mc2-bar-intrinsic}). Projection~\textup{(P6)} is
the shadow-connection construction of
Proposition~\ref{prop:hdm-upgrade-theta}.
\end{proof}

\begin{example}[Flagship verifications]
\label{ex:hdm-u-a-flagships}
For the three boundary flagships the universal property of
Theorem~\ref{thm:hdm-u-a-universal-property} computes as follows.

\emph{Heisenberg $\mathcal H_k$.} The Verdier-dual companion is
self-dual, $\cA^{!}_\infty \simeq \mathrm{Sym}^{\mathrm{ch}}(V^*[-1])$
on the Heisenberg line $\kappa = k$. The representing object is the
free chiral Weyl algebra
$U_{\mathcal H_k} \simeq \mathrm{Sym}^{\mathrm{ch}}(V \oplus V^*[-1])$,
whose quasitriangular $R$-matrix is
$r^{\mathrm{Heis}}_k(z) = (k/z)\, J \otimes J$, recovering slot~\textup{(4)}.

\emph{Affine Kac--Moody $V_k(\mathfrak g)$.} The coderived dual is
$V_k(\mathfrak g)^{!}_\infty \simeq
V_{-k - 2 h^\vee}(\mathfrak g)^{\text{copy}}_\infty$ on the
trace-form family, and the representing object
$U_{V_k(\mathfrak g)} = V_k(\mathfrak g) \bowtie
V_k(\mathfrak g)^{!}_\infty$ has universal $R$-matrix
$r^{KM,\fg}_{\mathrm{coll},k}(z) = k\, \Omega_{\mathrm{tr}}/z$
in the trace-form collision lane, equivalently
$r^{KM,\fg}_{\mathrm{KZ},k}(z)=\Omega/((k+h^\vee)z)$ after the
Sugawara/KZ Casimir rescaling, recovering slot~\textup{(4)}.
Theorem~C (central-charge complementarity,
Theorem~\ref{thm:chiral-center-theorem}) reads here as
$\kappa(V_k) + \kappa(V_k^!) = \dim(\mathfrak g)/2$ on the trace-form
normalization.

\emph{Virasoro $\mathrm{Vir}_c$.} The coderived dual is
$\mathrm{Vir}_c^{!}_\infty \simeq \mathrm{Vir}_{c'}^{\text{copy}}_\infty$
with $c + c' = 26$ (the critical-central-charge complementarity of
Theorem~C on the Virasoro lane, $c/2 + c'/2 = 13$). The representing
object $U_{\mathrm{Vir}_c} = \mathrm{Vir}_c \bowtie
\mathrm{Vir}_{26 - c}^{!}_\infty$ has universal $R$-matrix
$r^{\mathrm{Vir}}_c(z) = (c/2)/z^3 + 2T/z$, recovering slot~\textup{(4)}.
On the nonlinear class~$\mathsf M$ locus the quartic contact class
$10/[c(5c+22)]$ (Zamolodchikov norm) is the obstruction to the
boundary-linear presentation of~$U_{\mathrm{Vir}_c}$, consistent with
Proposition~\ref{prop:hdm-upgrade-bulk-slot}.
\end{example}

\begin{proposition}[Four recovery directions with precise scope;
\ClaimStatusConditional]
\label{prop:hdm-four-recovery-directions}
\index{holographic modular Koszul datum!recovery directions and scope|textbf}
For the six-slot datum
$\mathcal{H}(T)=(\cA,\cA^!,\cC,r_\cA(z),\Theta_\cA,\nabla^{\mathrm{hol}})$
of Definition~\ref{def:holographic-modular-koszul-datum}, the four
recovery directions currently proved in the manuscript, with the
off-Koszul Verdier/Koszul direction retaining the conditional scope of
Theorem~\textup{\ref{thm:bar-cobar-inversion-qi}}, are the following.
\begin{enumerate}[label=\textup{(\roman*)}]
\item \emph{Boundary recovery from a realised HT theory.}
 On the source category
 $\mathrm{ChirAlg}^{\omega,\mathrm{BL}}_X$
 of boundary-linear or Drinfeld--Sokolov origin, Vol~II constructs a
 canonical HT theory~$\mathcal T_\cA$ with
 $\mathrm{Obs}^{\partial}(\mathcal T_\cA)\simeq \cA$
 \textup{(}the theorem ``Universal holography, class-$\mathsf{M}$
 chain-level'', part~\textup{(i)}\textup{)}.
 This is a proved recovery of the boundary from the \emph{constructed}
 theory~$\mathcal T_\cA$ on that lane. No theorem in Volume~I or
 Volume~II recovers~$\cA$ from an arbitrary bulk object~$\cC$ alone.
 The obstruction is structural: the current theorem begins with a
 chosen theory~$\mathcal T_\cA$, and no bulk-to-boundary restriction
 map from an arbitrary bulk object to~$\cA$ has been constructed.
\item \emph{Verdier/Koszul recovery.}
 On the modular Koszul locus there are canonical equivalences
 $\mathbb H_X(\cA)\simeq \cA^!$ and
 $\mathbb H_X^2(\cA)\simeq \cA$
 \textup{(}Theorem~\ref{thm:frontier-protected-bulk-antiinvolution}\textup{)}.
 Off the Koszul locus one retains only the factorization companion
 $\cA^{!}_\infty$ together with the coderived bar-cobar correspondence
 of Theorem~\ref{thm:bar-cobar-inversion-qi}; the strict recovery of
~$\cA^!$ is conjectural there. The obstruction is the failure of the
strict edge $\cA^{!}_\infty\to\cA^!$; the admissible
$L_{-1/2}(\mathfrak{sl}_2)$ class $\eta_4$ is the concrete witness.
\item \emph{Recovery of the bulk/line slot.}
 For every curved $A_\infty$-chiral algebra~$\cA$, the universal
 open/closed pair identifies the pinned bulk slot with
 $\cC=\mathcal{C}^{\bullet}_{\mathrm{ch}}(\cA,\cA)
 \simeq \cZ^{\mathrm{der}}_{\mathrm{ch}}(\cA)$
 \textup{(}Theorem~\ref{thm:thqg-swiss-cheese}\textup{)}.
 On the chirally Koszul line side one has
 $\cC_{\mathrm{line}} \simeq
 \cA^!_{\mathrm{line}}\text{-}\mathbf{mod}$
 \textup{(}Vol~II, theorem ``Lines as Modules for the Open-Colour
 Dual''\textup{)}.
 The global triangle identifies these two presentations for
 chirally Koszul boundary-linear families of classes
 $\mathsf{G}$, $\mathsf{L}$, $\mathsf{C}$
 \textup{(}Vol~II, theorem ``Global triangle for boundary-linear
 families''\textup{)} and on
 the principal DS lane at non-critical level
 \textup{(}Vol~II, the theorem ``Universal holography,
 class-$\mathsf{M}$ chain-level'' together with the corollary
 ``Class~$\mathsf{M}$ chain-level bulk reconstruction''\textup{)}.
 There is presently no theorem reconstructing this third slot from the
 smaller pair~$(\cA,r_\cA(z))$ alone. Outside the named lanes the
 global triangle fails first at the strict edge
 $\cA^{!}_\infty\to\cA^!$, and on nonlinear class~$\mathsf{M}$ the
 quartic contact class $10/[c(5c+22)]$ blocks the boundary-linear
 reduction even when a DS comparison theory exists. The binary datum
 $r_\cA(z)$ therefore does not determine~$\cC$ by itself.
\item \emph{Shadow connection recovery.}
 For modular Koszul~$\cA$, the Maurer--Cartan element determines the
 flat genus-zero shadow connection
 $\nabla^{\mathrm{hol}}_{0,n}=d-\operatorname{Sh}_{0,n}(\Theta_\cA)$,
 and on the established comparison surfaces its horizontal sections are
 the sphere correlators
 \textup{(}Theorem~\ref{thm:sphere-reconstruction}\textup{)}.
 Thus the construction
 $\Theta_\cA \mapsto \nabla^{\mathrm{hol}}$ is proved. What remains
 conjectural is a uniform comparison with arbitrary physical moduli
 problems outside those comparison surfaces. The obstruction is that no
 theorem presently identifies the horizontal sections of
 $\nabla^{\mathrm{hol}}_{g,n}$ with a physical moduli problem beyond
 the listed sphere/comparison lanes.
\end{enumerate}
\end{proposition}

\begin{proof}
Item~(i) is exactly the boundary-recovery clause of Vol~II's theorem
``Universal holography, class-$\mathsf{M}$ chain-level'' on the source category
$\mathrm{ChirAlg}^{\omega,\mathrm{BL}}_X$; the proof there is by
construction of~$\mathcal T_\cA$. No other theorem statement in
Volumes~I--II upgrades a general bulk object~$\cC$ to a boundary
reconstruction theorem, so the final sentence of~(i) is the exact
present scope.

Item~(ii) is Theorem~\ref{thm:frontier-protected-bulk-antiinvolution}
on the modular Koszul locus, while the off-locus limitation is
Theorem~\ref{thm:bar-cobar-inversion-qi}: strict inversion is replaced
there by a coderived statement. The explicit obstruction class is the
same $\eta_4$ exhibited in Example~\ref{ex:admissible-sl2-failure}.

For item~(iii), Theorem~\ref{thm:thqg-swiss-cheese} identifies the
universal bulk with chiral Hochschild cochains for every~$\cA$.
Vol~II's theorem ``Lines as Modules for the Open-Colour Dual''
identifies the line category with modules over the open-colour dual on
the chirally Koszul locus. Vol~II's theorem ``Global triangle for
boundary-linear families'' and the theorem ``Universal holography,
class-$\mathsf{M}$ chain-level'' together with the corollary
``Class~$\mathsf{M}$ chain-level bulk reconstruction'' are the precise
loci where the bulk and line presentations are glued. The search surface
of the manuscript contains no theorem whose hypotheses are merely the
pair~$(\cA,r_\cA(z))$, so the last sentence of~(iii) is the current exact
theorem surface; the obstruction boundary is the strict edge
$\cA^{!}_\infty\to\cA^!$ together with the nonlinear quartic contact
class already named in Remark~\ref{rem:hdm-boundary-linear-vs-global-triangle}.

Item~(iv) is the content of
Theorem~\ref{thm:sphere-reconstruction}: flatness follows from the MC
equation, while the correlator interpretation is proved on the listed
comparison surfaces only. Beyond them the missing input is an
independent physical moduli problem whose boundary restriction agrees
with the horizontal sections of~$\nabla^{\mathrm{hol}}_{g,n}$.
\end{proof}

\begin{remark}[Verification at Heisenberg]
\label{rem:hdm-heisenberg-verification}
\index{Heisenberg algebra!holographic datum verification}
Take $\cA = \cH_k$ with OPE $J(z)\,J(w) \sim k/(z-w)^2$ and
$k \neq 0$. Each slot of $\mathcal{H}(T)$ can be computed directly
and matched against
Proposition~\ref{prop:hdm-u-a-recoverability}:
\begin{enumerate}[label=\textup{(H\arabic*)}]
\item Boundary factor: $\cA = \cH_k$.
\item Dual factor: $\cA^! = \cH_{-k}$
 (Koszul self-duality of Heisenberg, with the sign on the level
 coming from the central-charge complementarity for free-field
 algebras, Theorem~\ref{thm:central-charge-complementarity}).
\item Bulk/line slot:
 $\cC = \mathcal{C}^{\bullet}_{\mathrm{ch}}(\cH_k,\cH_k)
 \simeq \cZ^{\mathrm{der}}_{\mathrm{ch}}(\cH_k)$ on the open/closed
 side, and in the free-field line presentation
 $\cC_{\mathrm{line}} \simeq \cH_{-k}\text{-}\mathsf{mod}$.
\item Quasitriangular $R$-matrix: the single-generator OPE has
 double pole $k$, and the $d\log$ absorption strips one
 pole, producing
 \[
 r^{\mathrm{Heis}}_k(z) \;=\; \frac{k}{z}\, J \otimes J,
 \]
 which is the level-$k$ Drinfeld $r$-matrix at $\fg = \C$; vanishing
 at $k = 0$ is the level sanity check.
\item Universal MC element: $\Theta_{\cH_k}$ has a single non-zero
 degree~$2$ component (class~$\mathbf{G}$, $r = 2$), and
 $\kappa(\cH_k) = k$ by the degree-$2$ projection
 (Heisenberg entry).
\item Shadow connection: the shadow tower terminates at $r = 2$
 (class~$\mathbf{G}$), so $\nabla^{\mathrm{hol}}$ is the trivial
 extension of the flat KZ connection beyond the binary level; all
 higher-order corrections vanish.
\end{enumerate}
The six slots therefore reduce to the rank-one free-field data:
$(\cH_k, \cH_{-k}, \cC, r^{\mathrm{Heis}}_k(z), \Theta_{\cH_k}, \nabla^{\mathrm{KZ}})$.
All six are projections of the Drinfeld double
$U_{\cH_k} = \cH_k \bowtie \cH_{-k}$, which is the free bosonic
bulk algebra corresponding to abelian Chern--Simons. This is the
atomic case of the boundary-bulk reconstruction thesis: the free
boundary algebra $\cH_k$ recovers the free abelian Chern--Simons
bulk via~$U_{\cH_k}$, with $\kappa = k$ as the degree-$2$ scalar
coefficient of~$\Theta_{\cH_k}$
and no higher corrections. Concretely, the bulk slot is
$\cC=\cZ^{\mathrm{der}}_{\mathrm{ch}}(\cH_k)$, the restriction map
$\rho_{\partial}\colon \cC\to\cH_k$ is the free Swiss-cheese
restriction, and the witness triple is
$(r^{\mathrm{Heis}}_k(z),\Theta_{\cH_k},\nabla^{\mathrm{KZ}})$. No
off-scope obstruction survives in this example because Heisenberg is
already boundary-linear and chirally Koszul.
\end{remark}

\subsection{$U_\cA$ as Drinfeld $\infty$-centre on the Koszul locus}
\label{subsec:hdm-u-a-drinfeld-center}

Theorem~\ref{thm:hdm-u-a-universal-property} names $U_\cA$ as the
representing object for Koszul pairs with~$\cA$. The chiral
Deligne--Tamarkin theorem
(Theorem~\ref{thm:chiral-deligne-tamarkin}) names
$\cZ^{\mathrm{der}}_{\mathrm{ch}}(\cA)
\simeq \mathcal{C}^{\bullet}_{\mathrm{ch}}(\cA,\cA)$ as the
terminal closed algebra acting on~$\cA$. The Ben-Zvi--Francis--Nadler
identification~\cite{BZFN10} names the Drinfeld $\infty$-centre of
$\mathrm{Rep}^{E_1}(\cA)$ as the Hochschild cochain object. One
question remains: do these three $E_2$-algebras coincide? The answer
is locus-dependent: on the modular Koszul locus the three agree on
the nose; off that locus the agreement holds only after
$\mathrm{Pro}$-completion, and the fibre obstruction is the same
strict edge $\cA^{!}_\infty\to\cA^!$ that obstructs Upgrade~(U2).

\begin{theorem}[$U_\cA$ as Drinfeld $\infty$-centre, on and off locus;
\ClaimStatusConditional]
\label{thm:hdm-u-a-drinfeld-center}
\index{Drinfeld $\infty$-centre!identification with $U_\cA$|textbf}
\index{Ben-Zvi--Francis--Nadler identification!chiral lift|textbf}
\index{Drinfeld double!$\infty$-centre interpretation|textbf}
Let $\cA$ be a curved $A_\infty$-chiral algebra on a smooth
projective curve~$X$, and let
$\mathcal Z_{(\infty,1)}\bigl(\mathrm{Rep}^{E_1}(\cA)\bigr)$ denote
the Drinfeld $\infty$-centre of its $E_1$-representation
$\infty$-category, namely the $\infty$-category of pairs
$(M,\sigma_{M,-})$ consisting of an $E_1$-module~$M$ together with a
half-braiding $\sigma_{M,N}\colon M\otimes N \simeq N\otimes M$
natural in $N\in \mathrm{Rep}^{E_1}(\cA)$, equipped with its canonical
$E_2$-algebra structure (Lurie~\cite[5.3]{HA}).
\begin{enumerate}[label=\textup{(\roman*)}]
\item \emph{On the modular Koszul locus.} If $\cA$ is modular Koszul,
there is an $E_2$-equivalence
\begin{equation}\label{eq:hdm-u-a-drinfeld-center-koszul}
U_\cA
\;\simeq\;
\mathcal Z_{(\infty,1)}\bigl(\mathrm{Rep}^{E_1}(\cA)\bigr)
\qquad \text{in } \mathrm{Alg}_{E_2}\bigl(\mathrm{Pr}^L_{\mathrm{st}}\bigr).
\end{equation}
The equivalence factors through the
Ben-Zvi--Francis--Nadler identification~\cite{BZFN10}:
$\mathcal Z_{(\infty,1)}\bigl(\mathrm{Rep}^{E_1}(\cA)\bigr)
\simeq \mathcal{C}^{\bullet}_{\mathrm{ch}}(\cA,\cA)
\simeq \cZ^{\mathrm{der}}_{\mathrm{ch}}(\cA)$, matched with
projection~\textup{(P3)} of
Corollary~\textup{\ref{cor:hdm-u-a-six-projections}}.
\item \emph{Off the Koszul locus.} In general there is an
$E_2$-equivalence only after $\mathrm{Pro}$-completion of both
sides:
\begin{equation}\label{eq:hdm-u-a-drinfeld-center-pro}
\mathrm{Pro}(U_\cA)
\;\simeq\;
\mathrm{Pro}\,\mathcal Z_{(\infty,1)}\bigl(\mathrm{Rep}^{E_1}(\cA)\bigr),
\end{equation}
with the natural map from~\eqref{eq:hdm-u-a-drinfeld-center-koszul}
to the $\mathrm{Pro}$ level inducing a strict $E_2$-equivalence if
and only if the edge $\cA^{!}_\infty\to\cA^!$ is a
quasi-isomorphism \textup{(}Upgrade~\textup{(U2)}\textup{)}.
\end{enumerate}
On the Koszul locus the boundary-bulk algebra is not merely a
representing object for Koszul pairs but simultaneously the Drinfeld
$\infty$-centre of its own module category; off the Koszul locus the
identification holds in the weight-completed limit only.
\end{theorem}

\begin{proof}
\emph{(i).} On the modular Koszul locus Upgrade~(U2) is available,
and the representing-object formula
$U_\cA \simeq \cA\otimes^{\mathbb L}_{\barBch(\cA)}\cA^{!}_\infty$
(equation~\eqref{eq:hdm-u-a-relative-tensor}) rectifies to the
matched-pair Hopf product $\cA\bowtie \cA^!$. Take derived
endomorphisms over the diagonal $\cA$-bimodule: projection~(P3) of
Corollary~\ref{cor:hdm-u-a-six-projections} identifies
$\mathrm{End}_{U_\cA\text{-}\mathrm{mod}}^\bullet(\cA)
\simeq \mathcal{C}^{\bullet}_{\mathrm{ch}}(\cA,\cA)$; the chiral
Deligne--Tamarkin theorem
(Theorem~\ref{thm:chiral-deligne-tamarkin}) supplies the $E_2$-brace
structure on the right. Ben-Zvi--Francis--Nadler~\cite{BZFN10}
identify Drinfeld $\infty$-centres with Hochschild cochains in the
derived-algebraic setting; the chiral refinement proceeds by replacing
ordinary tensor products with chiral tensor over the Ran space, which
is Lurie's~\cite{HA} factorisable $\infty$-operadic formalism applied
to the factorisation $\infty$-category $\mathrm{ChirAlg}^{\mathrm{MK}}_X$.
The composite
\[
U_\cA
\xrightarrow{\text{(P3)}}
\mathcal{C}^{\bullet}_{\mathrm{ch}}(\cA,\cA)
\xrightarrow[\cong]{\text{BZFN}}
\mathcal Z_{(\infty,1)}\bigl(\mathrm{Rep}^{E_1}(\cA)\bigr)
\]
is an $E_2$-equivalence: the first arrow is the universal property of
the relative tensor product evaluated on the boundary regular module
(Lurie~\cite[5.5]{HA}); the second is~\cite{BZFN10} Thm.~1.2 adapted
to the factorisation setting. The $E_2$-structure on
$\mathcal Z_{(\infty,1)}$ is the half-braiding monoidal structure; its
match with the Gerstenhaber bracket on Hochschild cochains is
Deligne's conjecture, proved in the chiral lane by
Theorem~\ref{thm:chiral-deligne-tamarkin}.

\emph{(ii).} Off the modular Koszul locus the strict edge
$\cA^{!}_\infty\to\cA^!$ is only a quasi-isomorphism of
pro-objects, not of strict chain complexes; the same obstruction
class $\eta_4$ that blocks Upgrade~(U2)
(Example~\ref{ex:admissible-sl2-failure}) blocks the strict
identification~\eqref{eq:hdm-u-a-drinfeld-center-koszul}. After
applying $\mathrm{Pro}$ to both sides, the weight-completed dual
$\mathrm{Pro}(\cA^{!}_\infty)$ is strictly inverse to~$\barBch(\cA)$
by Theorem~\ref{thm:bar-cobar-inversion-qi}, and the relative tensor
product formula~\eqref{eq:hdm-u-a-relative-tensor} becomes an
equivalence of $\mathrm{Pro}$-objects; BZFN~\cite{BZFN10} applied in
$\mathrm{Pro}$-completion then yields~\eqref{eq:hdm-u-a-drinfeld-center-pro}.
The converse implication is by inspection: a strict
$E_2$-equivalence $U_\cA \simeq \mathcal Z_{(\infty,1)}$ forces the
forgetful map $U_\cA \twoheadrightarrow \cA^{!}_\infty$ to factor
through~$\cA^!$, which is exactly Upgrade~(U2).
\end{proof}

\begin{conjecture}[Strict off-locus Drinfeld-centre identification]
\label{conj:hdm-u-a-drinfeld-center-strict-off-locus}
\index{Drinfeld $\infty$-centre!strict off-locus conjecture}
For every curved $A_\infty$-chiral algebra $\cA$ outside the modular
Koszul locus, the natural comparison
$U_\cA \to \mathcal Z_{(\infty,1)}\bigl(\mathrm{Rep}^{E_1}(\cA)\bigr)$
fails to be a strict $E_2$-equivalence before $\mathrm{Pro}$-completion,
with obstruction measured by
$\mathrm{coker}(\cA^{!}_\infty\to\cA^!)$ as in
equation~\eqref{eq:hdm-u-a-drinfeld-fibre}. Strict equivalence is
equivalent to Upgrade~\textup{(U2)} and is therefore as open as
Upgrade~(U2) itself.
\end{conjecture}

\begin{remark}[Drinfeld centre is not averaging]
\label{rem:hdm-drinfeld-not-averaging}
\index{averaging map!not the Drinfeld centre|textbf}
The identification of Theorem~\ref{thm:hdm-u-a-drinfeld-center}
is not a statement about the averaging map
$\mathrm{av}\colon \mathfrak g^{E_1}\to \mathfrak g^{\mathrm{mod}}$.
The averaging map is a strict projection of ordered onto symmetric
chiral homology and produces a smaller object: its image sits inside
the Drinfeld $\infty$-centre as the $S_n$-invariants of the
half-braided pair $(M,\sigma_{M,-})$ for $n$ the number of inputs.
The Drinfeld centre records all half-braidings, including the
non-symmetrisable ones that averaging annihilates. On the Virasoro
lane at generic~$c$ the averaging map kills the antisymmetric
component of the degree-$2$ bracket, whereas
$\mathcal Z_{(\infty,1)}(\mathrm{Rep}^{E_1}(\mathrm{Vir}_c))$ retains
it as the half-braiding $\sigma_{T,L_{-2}v_0}$ on the stress-tensor
module. The inscription $U_\cA \simeq \mathcal Z_{(\infty,1)}$ must
therefore be read in the non-averaged ordered lane; passing to the
symmetric chiral realm via averaging would produce a strictly
smaller centre: the Drinfeld centre is computed pre-averaging on the
ordered lane, not post-averaging.
\end{remark}

\begin{corollary}[Flagship verifications of the Drinfeld-centre
identification; \ClaimStatusConditional]
\label{cor:hdm-u-a-drinfeld-center-flagships}
\index{Drinfeld $\infty$-centre!Heisenberg verification}
\index{Drinfeld $\infty$-centre!affine Kac--Moody verification}
\index{Drinfeld $\infty$-centre!Kazhdan--Lusztig lift}
For the three boundary flagships
Theorem~\ref{thm:hdm-u-a-drinfeld-center}\textup{(i)} reads as
follows.
\begin{enumerate}[label=\textup{(F\arabic*)}]
\item \emph{Heisenberg $\cH_k$ at $k\neq 0$.} The
$E_1$-representation category is
$\mathrm{Rep}^{E_1}(\cH_k) \simeq \mathrm{Fock}_k$, the category of
Fock modules at level~$k$. Its Drinfeld $\infty$-centre is the
bosonic bulk algebra $\cH_k \otimes \cH_{-k}$, equal as an
$E_2$-algebra to $U_{\cH_k} = \cH_k \bowtie \cH_{-k}$; the
half-braiding is the $U(1)$ monodromy
$\sigma_{V_\alpha,V_\beta}= e^{i\pi \alpha\beta / k}$ on vertex operators.
\item \emph{Affine Kac--Moody $V_k(\fg)$ at non-critical $k$.}
The Kazhdan--Lusztig equivalence
$\mathrm{Rep}^{E_1}(V_k(\fg)) \simeq U_q(\fg)\text{-}\mathrm{mod}$
with $q = e^{i\pi/(k+h^\vee)}$ lifts to an $E_2$-equivalence
$U_{V_k(\fg)} \simeq
\mathcal Z_{(\infty,1)}(U_q(\fg)\text{-}\mathrm{mod})
\simeq D(U_q(\fg))\text{-}\mathrm{mod}$,
the module category of the Lusztig quantum double. The universal
$R$-matrix of~$U_{V_k(\fg)}$ is the trace-form collision residue
$r^{KM,\fg}_{\mathrm{coll},k}(z) = k\,\Omega_{\mathrm{tr}}/z$
and, after the KZ Casimir rescaling, matches the quasitriangular
structure on $D(U_q(\fg))$ under this equivalence.
\item \emph{Virasoro $\mathrm{Vir}_c$ on the Koszul locus.}
For~$c$ on the Koszul locus of Theorem~C,
$\mathrm{Rep}^{E_1}(\mathrm{Vir}_c)$ is equivalent to the category
of modules over a quantum-Virasoro algebra
$U_q(\mathrm{Vir})$ with~$q$ determined by the Liouville relation
$c = 1 + 6(b+b^{-1})^2$, $q = e^{i\pi b^2}$. Its Drinfeld
$\infty$-centre is $U_{\mathrm{Vir}_c} = \mathrm{Vir}_c\bowtie
\mathrm{Vir}_{26-c}^{!}_\infty$, with universal $R$-matrix
$r^{\mathrm{Vir}}_c(z) = (c/2)/z^3 + 2T/z$ projecting to the
half-braiding of~$T$ with itself.
\end{enumerate}
\end{corollary}

\begin{proof}
\emph{(F1).} Heisenberg is modular Koszul: the Koszul-duality
witness is $\cH_k^! = \cH_{-k}$. Apply
Theorem~\ref{thm:hdm-u-a-drinfeld-center}\textup{(i)} and unpack
Ben-Zvi--Francis--Nadler for the abelian vertex algebra
$\cH_k$: half-braidings of Fock modules~$V_\alpha$ are the $U(1)$
monodromy phases, and the Drinfeld centre
$\cH_k\otimes \cH_{-k}$ realises these as the doubled lattice
algebra. The $E_2$-equivalence with~$U_{\cH_k}$ is the
self-dual case of~\eqref{eq:hdm-u-a-drinfeld-center-koszul}.

\emph{(F2).} Kazhdan--Lusztig~\cite{KL93}
equate $\mathrm{Rep}^{E_1}(V_k(\fg))$ with $U_q(\fg)\text{-}\mathrm{mod}$
at $q = e^{i\pi/(k+h^\vee)}$ for non-critical non-rational~$k$. The
Drinfeld $\infty$-centre of the right-hand side is the module category
of the quantum double $D(U_q(\fg))$, by the Majid--Joseph
construction upgraded to the $\infty$-categorical setting via
BZFN~\cite{BZFN10}. Matching the representing object
$U_{V_k(\fg)}$ with $D(U_q(\fg))$ is a direct verification at the
level of $R$-matrices: the trace-form spectral $R$-matrix
$r^{KM,\fg}_{\mathrm{coll},k}(z) = k\,\Omega_{\mathrm{tr}}/z$
and the KZ spectral $R$-matrix
$r^{KM,\fg}_{\mathrm{KZ},k}(z)=\Omega/((k+h^\vee)z)$ are related by
the same finite Casimir-rescaling gauge. The KZ presentation expands
as a rational function of~$z$ whose monodromy at~$z=\infty$ reproduces
Drinfeld's universal $R$-matrix on~$D(U_q(\fg))$, by the Drinfeld--Kohno
isomorphism~\cite{Drinfeld85}. The level condition
$k\neq -h^\vee$ (non-critical) ensures~$q$ is not a root of unity of
order~$\leq 2h^\vee$ and that the KL equivalence is an equivalence
of $E_2$-categories.

\emph{(F3).} The Liouville identification of~$\mathrm{Vir}_c$ with
quantum-Virasoro is Fateev--Lukyanov~\cite{FL88} at
$q = e^{i\pi b^2}$; the Koszul hypothesis on~$c$ ensures
Upgrade~(U2) holds, so
Theorem~\ref{thm:hdm-u-a-drinfeld-center}\textup{(i)} applies. The
self-braiding of the stress tensor~$T$ is governed by the
Zamolodchikov residue $c(5c+22)/10$ and matches the $z^{-3}$
coefficient of $r^{\mathrm{Vir}}_c(z) = (c/2)/z^3 + 2T/z$ under the
half-braiding $\sigma_{T,T}$.
\end{proof}

\begin{remark}[Fibre obstruction off the Koszul locus]
\label{rem:hdm-u-a-drinfeld-center-fibre}
\index{fibre obstruction!Drinfeld centre vs.\ $U_\cA$|textbf}
\index{strict edge obstruction!Drinfeld centre interpretation}
The fibre of the comparison map
$U_\cA \to \mathcal Z_{(\infty,1)}\bigl(\mathrm{Rep}^{E_1}(\cA)\bigr)$
over a non-Koszul~$\cA$ is the homotopy cokernel of the strict edge
$\cA^{!}_\infty \to \cA^!$, weight-completed over the bar coalgebra:
\begin{equation}\label{eq:hdm-u-a-drinfeld-fibre}
\mathrm{fib}\bigl(U_\cA \to \mathcal Z_{(\infty,1)}\bigr)
\;\simeq\;
\mathrm{coker}\bigl(\cA^{!}_\infty \to \cA^!\bigr)
\,\otimes^{\mathbb L}_{\barBch(\cA)}\, \cA.
\end{equation}
The cokernel is zero exactly on the modular Koszul locus; the
admissible $L_{-1/2}(\mathfrak{sl}_2)$ class~$\eta_4$ of
Example~\ref{ex:admissible-sl2-failure} is the concrete non-zero
witness off that locus. On nonlinear class~$\mathsf{M}$ away from
the Liouville Koszul curve, the fibre picks up in addition the
quartic contact class $10/[c(5c+22)]$, which is the obstruction to
strict quasitriangularity of $U_{\mathrm{Vir}_c}$ beyond the
Koszul $c$-values. The $\mathrm{Pro}$-completed identification of
Theorem~\ref{thm:hdm-u-a-drinfeld-center}\textup{(ii)} holds exactly
because the pro-object formed by the quotient filtration on
$\cA^!_\infty$ kills~\eqref{eq:hdm-u-a-drinfeld-fibre} in the limit.
\end{remark}

\begin{remark}[Four consistency conditions for the six-slot thesis]
\label{rem:hdm-four-consistency-conditions}
\index{Drinfeld double!four consistency conditions|textbf}
The six-slot recovery of $\mathcal{H}(T)$ from~$U_\cA$ is not
tautological: for the thesis to hold, the Drinfeld double and the
six projections must satisfy four compatibility conditions. Each
condition was flagged as an open problem in the frontier audit of
the Drinfeld double programme
(Section~\ref{sec:frontier-modular-holography-platonic}, parts~(a)--(d));
the reading of
Proposition~\ref{prop:hdm-u-a-recoverability} converts the four open
problems into four consistency laws that the six-tuple must obey.
\begin{enumerate}[label=\textup{(C\arabic*)}]
\item \emph{Quasitriangular compatibility.} The spectral $R$-matrix
 of~$U_\cA$ must satisfy the classical Yang--Baxter equation on the
 nose, not merely up to homotopy. This is the collision-depth-$2$
 projection of the MC equation
 (Theorem~\ref{thm:collision-depth-2-ybe}) and is verified in
 slot~(R4) by the seven-way master theorem.
\item \emph{Factorization compatibility.} The factorization
 structure of~$U_\cA$ over~$\cC$ must be compatible with the
 Ran-space factorization of the underlying chiral algebra~$\cA$,
 so that the bulk/line slot~$\cC$ carries a single well-defined
 factorization model. This is the content of the
 Verdier-intertwining part of Theorem~A at the level of the
 open/closed bulk model. The obstruction boundary is not the
 existence of~$\cC$ but its comparison with the line model: away from
 the boundary-linear / DS lanes there is no theorem producing a single
 factorization model on both sides.
\item \emph{Shadow compatibility.} The modular shadow connection
 $\nabla^{\mathrm{hol}}$ produced from slot~(R5) by
 projection~(R6) must agree with the independent genus-expansion
 definition of the shadow connection on every $(g,n)$ stratum.
 This is the flatness
 $(\nabla^{\mathrm{hol}}_{g,n})^2 = 0$, a consequence of the MC
 equation, verified on the sphere by
 Theorem~\ref{thm:sphere-reconstruction}; at higher genus it is
 the content of the analytic sewing programme
 (conditional beyond genus~$1$ because a degeneration-independent
 sewing/regularization theorem has not yet been proved there).
\item \emph{Duality compatibility.} The second tensor factor of
 $U_\cA$ must agree with the chiral Koszul dual~$\cA^!$ of~$\cA$
 on cohomology, with the passage
 $\cA^{!}_\infty \to \cA^!$ controlled by bar-cobar inversion on
 the Koszul locus. This is Theorem~B; beyond the Koszul locus the
 identification is conjectural, and the explicit obstruction is the
 surviving bar class $\eta_4$ of
 Example~\ref{ex:admissible-sl2-failure}.
\end{enumerate}
Conditions~(C1)--(C3) are proved in the stated scopes; condition~(C4)
is the conditional input that makes
Proposition~\ref{prop:hdm-u-a-recoverability} heuristic rather than
unconditional. Together, the four conditions are the four
structural requirements that the Drinfeld double programme must
satisfy for the boundary-bulk reconstruction thesis to hold. Here that
thesis means the chain-level comparison between boundary object~$\cA$
and bulk slot~$\cC$ along
$\rho_{\partial}\colon\cC\to\cA$ and
$\mathrm{Ass}^{\mathrm{ch}}\hookrightarrow\mathrm{SC}^{\mathrm{ch,top}}$,
witnessed by $(r_\cA(z),\Theta_\cA,\nabla^{\mathrm{hol}})$ on the
named modular-Koszul / boundary-linear / principal DS surfaces; beyond
them the strict edge, the quartic contact class, and the missing
higher-genus sewing theorem are the active obstructions. Each one is a
falsification target in the sense of
Principle~\ref{princ:hdm-falsification}.
\end{remark}

\subsection{$U_\cA$ as final object in the $(\infty,2)$-category of
Koszul braided pairs}
\label{subsec:hdm-u-a-final-object}

Theorem~\ref{thm:hdm-u-a-universal-property} presents~$U_\cA$ as an
\emph{initial} object: the representing object of Koszul pairs
$(P,q,\alpha)$ mapping \emph{into}~$\cA$. The representability
formula~\eqref{eq:hdm-u-a-representable} is corepresentation on the
twisting-morphism side. The dual question asks which universal property
$U_\cA$ satisfies \emph{after} one has selected the boundary
factor~$\cA$ and the Verdier-dual factor~$\cA^!_\infty$: not what maps
into~$U_\cA$, but what \emph{receives} maps from every compatible
braided Hopf-like structure covering the pair
$(\cA,\cA^!_\infty)$. The answer is the finality content below. The
initial reading produces~$U_\cA$ from the bar data; the final reading
produces~$U_\cA$ from the $E_2$-data. On the modular Koszul locus the
two readings coincide strictly; off it they agree after
$\mathrm{Pro}$-completion.

Fix~$\cA$ modular Koszul on a smooth projective curve~$X$. Write
$\mathsf{KosPair}^{\mathrm{br}}_{\cA,\cA^!_\infty}$ for the
$(\infty,2)$-category whose objects are triples $(H,\rho_\partial,
\rho^!)$ consisting of a braided chiral Hopf-like object
$H\in\mathrm{Alg}_{E_2}(\mathrm{ChirAlg}^{\mathrm{MK}}_X)$ together
with a boundary restriction
$\rho_\partial\colon H\to\cA$ in $\mathrm{Alg}_{E_1}$ and a dual
restriction $\rho^!\colon H\to\cA^!_\infty$ in
$\mathrm{CoAlg}_{E_1}^{\mathrm{co}}$, subject to the matched-pair
compatibility
\begin{equation}\label{eq:hdm-matched-pair}
\bigl[\rho_\partial(x),\rho^!(y)\bigr]_{\mathrm{br}}
\;=\;
\rho_\partial\bigl(x\triangleleft y\bigr)
-\rho^!\bigl(x\triangleright y\bigr)
\qquad (x,y\in H),
\end{equation}
where $(\triangleleft,\triangleright)$ is the matched-pair action
on~$H$ induced by~$\bowtie$. The 1-morphisms are $E_2$-algebra maps
commuting with~$\rho_\partial,\rho^!$; the 2-morphisms are
$E_2$-homotopies. The representing object
$U_\cA=\cA\bowtie\cA^!_\infty$ is itself an object, with
$\rho_\partial,\rho^!$ the two tensor-factor projections of
Corollary~\ref{cor:hdm-u-a-six-projections}(P1)--(P2).

\begin{theorem}[Final-object characterisation of $U_\cA$;
\ClaimStatusConditional]
\label{thm:hdm-u-a-final-object}
\index{Drinfeld double!final-object universal property|textbf}
\index{universal property!final object in Koszul braided pairs|textbf}
Let $\cA$ be a modular Koszul chiral algebra on~$X$ and let
$\mathsf{KosPair}^{\mathrm{br}}_{\cA,\cA^!_\infty}$ be the
$(\infty,2)$-category of braided Koszul pairs over
$(\cA,\cA^!_\infty)$.
\begin{enumerate}[label=\textup{(\roman*)}]
\item \emph{On the Koszul locus.} The boundary-bulk algebra
$U_\cA=\cA\bowtie\cA^!_\infty$ is the final object of
$\mathsf{KosPair}^{\mathrm{br}}_{\cA,\cA^!_\infty}$: for every
object $(H,\rho_\partial,\rho^!)$, the mapping space
\begin{equation}\label{eq:hdm-u-a-final-mapping}
\mathrm{Map}_{\mathsf{KosPair}^{\mathrm{br}}_{\cA,\cA^!_\infty}}
\bigl((H,\rho_\partial,\rho^!),\,(U_\cA,\pi_\cA,\pi_{\cA^!_\infty})\bigr)
\end{equation}
is contractible. The terminal map is the pair
$(\rho_\partial,\rho^!)\colon H\to\cA\bowtie\cA^!_\infty$, unique up
to contractible choice of matched-pair homotopy.
\item \emph{Off the Koszul locus.} After $\mathrm{Pro}$-completion of
both sides, $\mathrm{Pro}(U_\cA)$ is the final object of
$\mathrm{Pro}\bigl(\mathsf{KosPair}^{\mathrm{br}}_{\cA,\cA^!_\infty}\bigr)$;
the strict finality statement \textup{(i)} holds if and only if the
edge $\cA^!_\infty\to\cA^!$ is a quasi-isomorphism
\textup{(}Upgrade~\textup{(U2)}\textup{)}.
\end{enumerate}
\end{theorem}

\begin{proof}
\emph{(i).} By Lurie's~\cite[Prop.~5.2.7.4]{HA} an object $Y$ of an
$(\infty,1)$-category is final iff the slice category $\mathcal C_{/Y}$
is equivalent to~$\mathcal C$ via the forgetful functor; applied to the
2-truncation of
$\mathsf{KosPair}^{\mathrm{br}}_{\cA,\cA^!_\infty}$ this reduces
to exhibiting, for each $(H,\rho_\partial,\rho^!)$, a contractible
space of maps to $(U_\cA,\pi_\cA,\pi_{\cA^!_\infty})$.

Define $\Phi_H\colon H\to U_\cA$ as the composite
\[
H
\xrightarrow{\Delta^{\mathrm{br}}_H}
H\otimes^{\mathrm{br}} H
\xrightarrow{\rho_\partial\otimes\rho^!}
\cA\otimes^\tau \cA^!_\infty
\;=\;U_\cA,
\]
where $\Delta^{\mathrm{br}}_H$ is the braided-Hopf coproduct and
$\otimes^\tau$ is the twisted tensor product with cocycle~$\Theta_\cA$
(Upgrade~(U1), Proposition~\ref{prop:hdm-upgrade-theta}). The
matched-pair compatibility~\eqref{eq:hdm-matched-pair} is exactly the
condition that $\Phi_H$ intertwines the braided products on both
sides, because the right-hand side of~\eqref{eq:hdm-matched-pair} is
the image under $\rho_\partial\otimes\rho^!$ of the braided bracket on
$H\otimes^{\mathrm{br}} H$. Factoring through the braided coproduct
is the $E_2$-version of the universal-enveloping construction
(Lurie~\cite[5.3.1]{HA}); the resulting map~$\Phi_H$ is the unique
$E_2$-algebra map to $U_\cA$ making the two triangles
\[
\rho_\partial = \pi_\cA\circ\Phi_H,
\qquad
\rho^! = \pi_{\cA^!_\infty}\circ\Phi_H
\]
commute up to canonical homotopy.

Contractibility of~\eqref{eq:hdm-u-a-final-mapping} reduces to
contractibility of the space of $E_2$-algebra maps $H\to U_\cA$
compatible with two fixed $E_1$-projections, which by
Lurie~\cite[Thm.~5.5.3.18]{HA} is computed as the derived limit
\[
\mathrm{Map}_{E_2}^{\rho_\partial,\rho^!}(H,U_\cA)
\;\simeq\;
\mathrm{Map}_{E_1}(H,\cA)\times_{\mathrm{Map}_{E_1}(H,\cA\otimes\cA^!_\infty)}
\mathrm{Map}_{E_1}(H,\cA^!_\infty),
\]
each factor being a point by the fixed projections $\rho_\partial$
and $\rho^!$, and the middle by the initial-object property of
Theorem~\ref{thm:hdm-u-a-universal-property} applied to the pushout
$\cA\otimes\cA^!_\infty$ under the boundary factorisation. The
resulting limit is contractible.

\emph{(ii).} Off the Koszul locus $\cA^!_\infty\to\cA^!$ is a
quasi-isomorphism only after $\mathrm{Pro}$-completion
(Theorem~\ref{thm:bar-cobar-inversion-qi}); the obstruction class
$\eta_4$ of Example~\ref{ex:admissible-sl2-failure} witnesses the
strict failure. Applying $\mathrm{Pro}$ to every slot of the proof
of~(i) rectifies the twisted tensor product and the $E_2$-universal
enveloping construction, giving the pro-finality
statement~\eqref{eq:hdm-u-a-drinfeld-center-pro}-style. Strict finality
forces the forgetful $U_\cA\twoheadrightarrow\cA^!_\infty$ to factor
through~$\cA^!$, which is exactly Upgrade~(U2). The if-and-only-if is
by inspection.
\end{proof}

\begin{remark}[Final versus initial: duality of presentations]
\label{rem:hdm-u-a-final-versus-initial}
\index{Drinfeld double!final versus initial reading|textbf}
Theorem~\ref{thm:hdm-u-a-universal-property} and
Theorem~\ref{thm:hdm-u-a-final-object} are two distinct
presentations, not redundant statements. The initial reading
presents~$U_\cA$ on the bar-coalgebra side: it is the object
that \emph{corepresents} twisting morphisms under~$\cA$, with the
universal element being the bar differential~$\Theta_\cA$. The final
reading presents~$U_\cA$ on the braided $E_2$-side: it is the object
that \emph{represents} the terminal braided Hopf covering of the
pair~$(\cA,\cA^!_\infty)$, with the universal element being the
matched-pair bracket. The two universal elements are dual
via the bar-cobar equivalence of Theorem~A: the twisting cocycle
$\Theta_\cA$ and the matched-pair bracket $[-,-]_{\mathrm{br}}$ are
exchanged by Koszul duality on the modular Koszul locus. This
duality of universal properties is the chiral analogue of the
Tannakian duality between algebras and their module categories:
Theorem~\ref{thm:hdm-u-a-universal-property} is the algebra side,
Theorem~\ref{thm:hdm-u-a-final-object} is the category side.
\end{remark}

\subsection{Canonicity of the matched-pair actions}
\label{subsec:hdm-matched-pair-canonicity}

The matched-pair
compatibility~\eqref{eq:hdm-matched-pair} is stated in terms of two
actions
$\triangleleft\colon H\otimes H\to H$ and
$\triangleright\colon H\otimes H\to H$ inherited from~$\bowtie$.
Read naively, $(\triangleleft,\triangleright)$ carries a gauge
freedom: the Drinfeld--Majid twist~\cite{Maj94,Maj95} that rescales
each action by a coboundary leaves the braided Hopf structure
on~$U_\cA$ unchanged on cohomology. On the modular Koszul locus the
freedom collapses. The bar differential and the cobar differential
are independent Koszul invariants of~$\cA$ -- neither is a gauge
artefact -- and together they rigidify
$(\triangleleft,\triangleright)$ up to inner automorphism of~$U_\cA$.
Off the Koszul locus the residual freedom is the $\mathrm{Pro}$-completion
twist class that obstructs Upgrade~(U2). The following theorem states
this dichotomy in its exact form.

\begin{theorem}[Koszul canonicity of the matched-pair actions;
\ClaimStatusConditional]
\label{thm:hdm-matched-pair-canonicity}
\index{matched pair!canonicity from bar--cobar data|textbf}
\index{Drinfeld double!canonical matched-pair structure|textbf}
\index{bar differential!canonical $\triangleright$ action|textbf}
\index{cobar differential!canonical $\triangleleft$ action|textbf}
Let $\cA$ be a modular Koszul chiral algebra on a smooth projective
curve~$X$, with bar differential
$d_{\barBch}\colon \barBch(\cA)\to \barBch(\cA)$ and cobar
differential $d_{\Omega^{\mathrm{ch}}}\colon
\Omega^{\mathrm{ch}}(\cA^!_\infty)\to
\Omega^{\mathrm{ch}}(\cA^!_\infty)$ induced by the universal twisting
morphism $\pi_\cA\colon \barBch(\cA)\to\cA$
\textup{(}Theorem~\textup{A},
Theorem~\textup{\ref{thm:bar-cobar-adjunction}}\textup{)}.
\begin{enumerate}[label=\textup{(\roman*)}]
\item \emph{On the Koszul locus.} The matched-pair actions
$(\triangleleft,\triangleright)$ on
$U_\cA=\cA\bowtie\cA^!_\infty$ are canonically determined by the
pair $(d_{\barBch},d_{\Omega^{\mathrm{ch}}})$. Explicitly,
\begin{equation}\label{eq:hdm-triangleright-canonical}
\triangleright\;=\;
\bigl(\mathrm{id}_{\cA}\otimes\pi_\cA\bigr)\circ
\bigl(\cA\otimes d_{\barBch}\bigr)\circ
\bigl(\mathrm{id}_{\cA}\otimes\iota_{\cA^!_\infty}\bigr),
\end{equation}
where $\iota_{\cA^!_\infty}\colon \cA^!_\infty\hookrightarrow
\barBch(\cA)$ is the Koszul inclusion of the strict dual into the
bar coalgebra; and
\begin{equation}\label{eq:hdm-triangleleft-canonical}
\triangleleft\;=\;
\bigl(\pi^\vee_\cA\otimes\mathrm{id}_{\cA^!_\infty}\bigr)\circ
\bigl(d_{\Omega^{\mathrm{ch}}}\otimes\cA^!_\infty\bigr)\circ
\bigl(\iota^\vee_\cA\otimes\mathrm{id}_{\cA^!_\infty}\bigr),
\end{equation}
where $\iota^\vee_\cA\colon \cA\hookrightarrow
\Omega^{\mathrm{ch}}(\cA^!_\infty)$ is the Koszul inclusion of the
primitive generators into the cobar algebra and
$\pi^\vee_\cA\colon \Omega^{\mathrm{ch}}(\cA^!_\infty)\to\cA$ is the
universal cotwisting morphism. Any other matched-pair
$(\triangleleft',\triangleright')$ on~$U_\cA$ compatible with the
quasitriangular structure of
Theorem~\textup{\ref{thm:hdm-u-a-universal-property}} differs from
\eqref{eq:hdm-triangleright-canonical}--\eqref{eq:hdm-triangleleft-canonical}
by an inner automorphism of~$U_\cA$ implemented by a unit
$u\in U_\cA^\times$ with $[u]=1$ in
$H^0\bigl(U_\cA/[U_\cA,U_\cA]\bigr)$, and this class
vanishes.
\item \emph{Off the Koszul locus.} The actions
$(\triangleleft,\triangleright)$ are canonically determined only
after $\mathrm{Pro}$-completion: the residual gauge freedom is a class
\begin{equation}\label{eq:hdm-matched-pair-gauge-class}
[\gamma_\cA]\;\in\;
H^1\bigl(\mathrm{Pro}\, U_\cA,\;\mathfrak g^{\mathrm{mod}}_\cA[1]\bigr),
\end{equation}
and $[\gamma_\cA]=0$ if and only if the strict edge
$\cA^!_\infty\to\cA^!$ is a quasi-isomorphism
\textup{(}Upgrade~\textup{(U2)}\textup{)}. The class
$[\gamma_\cA]$ coincides with the Upgrade~\textup{(U2)} obstruction
class of Proposition~\textup{\ref{prop:hdm-upgrade-rectification}}.
\end{enumerate}
\end{theorem}

\begin{proof}
\emph{(i).} On the modular Koszul locus the bar--cobar adjunction
(Theorem~\ref{thm:bar-cobar-adjunction}) is a strict equivalence of
$(\infty,1)$-categories
$\Omega^{\mathrm{ch}}\dashv \barBch$, and the Koszul
inclusion $\iota_{\cA^!_\infty}\colon\cA^!_\infty\hookrightarrow
\barBch(\cA)$ is a chain-level quasi-isomorphism onto the
cohomology of the bar coalgebra
(Theorem~\ref{thm:bar-cobar-inversion-qi} under U2). Formula
\eqref{eq:hdm-triangleright-canonical} is read off from this strict
adjunction: the action of $\cA^!_\infty$ on~$\cA$ is the transpose,
under the Koszul pairing, of the bar differential evaluated on
primitive elements, and the outer composition with $\pi_\cA$ projects
back to~$\cA$. Formula
\eqref{eq:hdm-triangleleft-canonical} is the dual construction with
$\barBch\leftrightarrow\Omega^{\mathrm{ch}}$ and
$\cA\leftrightarrow\cA^!_\infty$ swapped. The compatibility
identity
$\triangleleft(\triangleright(a\otimes b)\otimes c)
=\triangleright(a\otimes\triangleleft(b\otimes c))$
is the Maurer--Cartan equation
$D\Theta_\cA+\tfrac12[\Theta_\cA,\Theta_\cA]=0$ read off on the
three-point collision stratum, rewritten via the boundary/dual
pairing (Theorem~\ref{thm:mc2-bar-intrinsic}).

Uniqueness up to inner automorphism. Two matched-pair structures
$(\triangleleft,\triangleright)$ and
$(\triangleleft',\triangleright')$ on~$U_\cA$ compatible with the
universal $R$-matrix
$r_\cA(z)=\mathrm{Res}^{\mathrm{coll}}_{0,2}(\Theta_\cA)$ give two
$E_2$-algebra maps $U_\cA\to U_\cA$ by
Theorem~\ref{thm:hdm-u-a-final-object}(i); the final-object property
forces these to agree up to contractible matched-pair homotopy,
implemented on chains by a unit
$u\in U_\cA^\times$ with $u\triangleright u^{-1}$ realising the
homotopy. The obstruction to trivialising $u$ to the identity is its
class in $H^0(U_\cA/[U_\cA,U_\cA])$, which vanishes because the
Koszul pairing $\cA\otimes\cA^!_\infty\to\mathbb C$ is nondegenerate
on cohomology (Theorem~\ref{thm:frontier-protected-bulk-antiinvolution},
$\mathbb H_X(\cA)\simeq\cA^!$). The vanishing of this class is
Etingof--Schiffmann's ``canonical form'' principle
\cite[§3]{EtingofSchiffmann} in the chiral setting: for a
quasitriangular Hopf object the matched-pair action is canonical up
to rigid rigidification once the $R$-matrix is fixed.

\emph{(ii).} Off the Koszul locus the strict Koszul inclusion
$\iota_{\cA^!_\infty}$ is replaced by its
$\mathrm{Pro}$-completion (Theorem~\ref{thm:bar-cobar-inversion-qi}).
The formulas
\eqref{eq:hdm-triangleright-canonical}--\eqref{eq:hdm-triangleleft-canonical}
continue to define actions on $\mathrm{Pro}(U_\cA)$, but the
uniqueness argument of part~(i) requires strict nondegeneracy of the
Koszul pairing; the failure of strict nondegeneracy is the
Upgrade~(U2) obstruction. The residual gauge freedom is the
class~$[\gamma_\cA]$ of~\eqref{eq:hdm-matched-pair-gauge-class},
computed as the
$\mathrm{Pro}$-completion of the difference
$(\triangleleft,\triangleright)-
(\triangleleft',\triangleright')$ between two compatible choices. A
chase of the coderived spectral sequence identifies $[\gamma_\cA]$
with the obstruction class of
Proposition~\ref{prop:hdm-upgrade-rectification}: both live in
$H^1(\mathrm{Pro}\,U_\cA,\mathfrak g^{\mathrm{mod}}_\cA[1])$, and
both vanish exactly under U2.
\end{proof}

\begin{remark}[Drinfeld double meta-principle: $R$-matrix fixes matched pair]
\label{rem:hdm-matched-pair-etingof-schiffmann}
\index{Drinfeld double!matched pair fixed by $R$-matrix}
\index{Etingof--Schiffmann canonical form}
The content of
Theorem~\ref{thm:hdm-matched-pair-canonicity}(i) is the chiral
transcription of a classical principle that goes back to Drinfeld and
is made precise by Etingof--Schiffmann~\cite{EtingofSchiffmann}: for
a quasitriangular Hopf algebra, the matched-pair action on a
Drinfeld-double factorisation
$D=A\bowtie A^*$ is canonically determined by the universal
$R$-matrix, up to inner automorphism. In the chiral setting the
universal $R$-matrix is the collision residue~$r_\cA(z)$, and the
canonical form of $(\triangleleft,\triangleright)$ is the pair of
bar/cobar differentials transposed along the Koszul pairing. No gauge
choice survives once bar and cobar are fixed.
\end{remark}

\begin{example}[Heisenberg, Kac--Moody, Virasoro: canonicity verification]
\label{ex:hdm-matched-pair-canonicity-flagships}
\index{Heisenberg algebra!matched pair canonical}
\index{affine Kac--Moody!matched pair canonical}
\index{Virasoro algebra!matched pair canonical off-locus gauge}
\emph{Heisenberg $\cH_k$.} On~$U_{\cH_k}=\cH_k\bowtie\cH_{-k}$ the
boundary factor is abelian and the dual factor self-dual: the
matched-pair is
$J\triangleright J^* = k$, $J^*\triangleleft J = -k$, with every
other action zero. Formula
\eqref{eq:hdm-triangleright-canonical} evaluates
$\triangleright(J\otimes J^*)$ on
$J^*\in\cH_{-k}\hookrightarrow\barBch(\cH_k)$
as the double-pole residue $k$, with strict equality. No gauge
freedom: the actions are uniquely $(k,-k)$.

\emph{Affine Kac--Moody $V_k(\fg)$ at non-critical~$k$.} The
matched-pair is
$J^a\triangleright (J^b)^* = k\,\delta^{ab}$ together with the
adjoint action of~$\fg$ on itself. Formula
\eqref{eq:hdm-triangleright-canonical} is the Wick contraction of
the bar differential on the current algebra, producing exactly
$k\,\Omega_{\mathrm{tr}}$ on the binary collision stratum. The
adjoint action on~$\fg$ is the internal Lie bracket, intrinsic to
the $\fg$-linear structure of~$V_k$. Again no gauge freedom:
$(\triangleleft,\triangleright)$ is the canonical Drinfeld double
action.

\emph{Virasoro $\mathrm{Vir}_c$ off the Koszul locus.}
On~$U_{\mathrm{Vir}_c}=\mathrm{Vir}_c\bowtie\mathrm{Vir}_{26-c}^!_\infty$,
the class~$\mathbf M$ quartic contact
$10/[c(5c+22)]$ (Zamolodchikov norm) obstructs the strict Koszul
inclusion $\mathrm{Vir}_{26-c}^!_\infty\to\mathrm{Vir}_{26-c}$.
Equation~\eqref{eq:hdm-matched-pair-gauge-class} has a nonzero class
$[\gamma_{\mathrm{Vir}_c}]$ valued in the first Hochschild cohomology
of the pro-Virasoro with coefficients in the shadow dGLA shifted by
one; the class lies in the same line as the Zamolodchikov quartic
obstruction. After $\mathrm{Pro}$-completion against the Liouville
$b$-parameter filtration the class becomes trivial and canonicity is
restored.
\end{example}

\begin{corollary}[Bondal reconstruction from finality;
\ClaimStatusConditional]
\label{cor:hdm-u-a-bondal-reconstruction}
\index{Bondal reconstruction!from final-object property|textbf}
\index{derived Morita equivalence!boundary from boundary-bulk algebra}
Under the hypotheses of Theorem~\ref{thm:hdm-u-a-final-object}(i),
the boundary algebra~$\cA$ is recovered from the
boundary-bulk algebra~$U_\cA$ up to derived Morita equivalence: the
perfect derived category
$\mathrm{Perf}^{\mathrm{ch}}(\cA)$
is equivalent to the full subcategory of compact
generators in $U_\cA\text{-}\mathrm{mod}$ on which
$\cA^!_\infty$ acts through its augmentation
$\epsilon\colon \cA^!_\infty\to\underline k$:
\begin{equation}\label{eq:hdm-u-a-bondal}
\mathrm{Perf}^{\mathrm{ch}}(\cA)
\;\simeq\;
\bigl(U_\cA\text{-}\mathrm{mod}\bigr)^{\omega,\epsilon_{\cA^!_\infty}}.
\end{equation}
The boundary projection~$P_\cA$ of
Theorem~\ref{thm:hdm-five-recovery-functors} is the
Bondal reconstruction functor; its right adjoint, induction along the
augmentation, is the inverse equivalence on perfect objects.
\end{corollary}

\begin{proof}
By Theorem~\ref{thm:hdm-u-a-final-object}(i) any braided Hopf $H$
carrying compatible projections $(\rho_\partial,\rho^!)$ maps uniquely
to~$U_\cA$. Take $H = \cA$ with $\rho_\partial=\mathrm{id}_\cA$ and
$\rho^! = \epsilon_{\cA^!_\infty}\colon \cA\to\cA^!_\infty$ through the
augmentation; the terminal map $\Phi_\cA\colon \cA\to U_\cA$ is the
boundary inclusion, and its right adjoint is the augmentation
localisation, producing the equivalence~\eqref{eq:hdm-u-a-bondal} by
Bondal--Kapranov reconstruction (Lurie~\cite[Prop.~7.1.2.6]{HA}).
Derived Morita equivalence follows because $\cA$ is a compact
generator of its own perfect derived category.
\end{proof}

\begin{example}[Final-object property at the three flagships]
\label{ex:hdm-u-a-final-flagships}
\index{Heisenberg algebra!final-object verification}
\index{affine Kac--Moody!final-object verification}
\index{Virasoro algebra!final-object verification}
The finality statement specialises to each boundary flagship as
follows; each verification is a direct computation of the contractible
mapping space~\eqref{eq:hdm-u-a-final-mapping}.

\emph{Heisenberg $\cH_k$.} Any braided chiral Hopf object $H$ with
projections $\rho_\partial\colon H\to\cH_k$,
$\rho^!\colon H\to\cH_{-k}$ satisfying the $U(1)$ matched-pair
$[\rho_\partial(J),\rho^!(J^*)]=k$ factors uniquely through
$U_{\cH_k}=\cH_k\otimes\cH_{-k}$ via
$\Phi_H(x)=\rho_\partial(x_{(1)})\otimes\rho^!(x_{(2)})$, the factorisation
of the braided coproduct. Self-duality $\cH_k^!=\cH_{-k}$ makes the
finality property the universal one for the free Weyl algebra.

\emph{Affine Kac--Moody $V_k(\fg)$ at non-critical~$k$.} The final
object receives the universal map from every braided $E_2$-Hopf
algebra covering the pair $(V_k(\fg),V_{-k-2h^\vee}(\fg)^!_\infty)$;
by the Kazhdan--Lusztig equivalence the finality property on the
module side identifies $U_{V_k(\fg)}\text{-}\mathrm{mod}$ with
$D(U_q(\fg))\text{-}\mathrm{mod}$, with $q=e^{i\pi/(k+h^\vee)}$, and
the Bondal reconstruction~\eqref{eq:hdm-u-a-bondal} recovers
$V_k(\fg)$-perfect modules from the quantum-double perfect modules
via the augmentation $\epsilon\colon U_q(\fg)^*\to \underline k$. The
terminal map from any Koszul pair is uniquely determined by its
universal $R$-matrix, written as $k\,\Omega_{\mathrm{tr}}/z$ in the
trace-form lane and as $\Omega/((k+h^\vee)z)$ in the KZ lane.

\emph{Virasoro $\mathrm{Vir}_c$ on the Koszul locus.} For~$c$ on the
Koszul locus of Theorem~C, the final object
$U_{\mathrm{Vir}_c}=\mathrm{Vir}_c\bowtie\mathrm{Vir}_{26-c}^!_\infty$
receives the universal map from every braided Hopf covering the
central-charge-complementary pair. On the nonlinear class~$\mathsf M$
locus the quartic contact class $10/[c(5c+22)]$ is the obstruction to
the strict finality statement \textup{(i)}; the pro-finality
statement~\textup{(ii)} remains available, and the Bondal
reconstruction~\eqref{eq:hdm-u-a-bondal} recovers
$\mathrm{Perf}^{\mathrm{ch}}(\mathrm{Vir}_c)$ only after
$\mathrm{Pro}$-completion against the Liouville $b$-parameter
filtration.
\end{example}

\begin{remark}[Finality is pre-averaging]
\label{rem:hdm-u-a-final-pre-averaging}
\index{averaging map!final-object property is pre-averaging|textbf}
The finality statement of Theorem~\ref{thm:hdm-u-a-final-object} is
formulated on the ordered $E_1$-chiral lane, before application of the
averaging map $\mathrm{av}\colon\mathfrak g^{E_1}\to\mathfrak
g^{\mathrm{mod}}$ of Remark~\ref{rem:hdm-drinfeld-not-averaging}.
Post-averaging, the symmetric chiral quotient
$\mathrm{av}_*(U_\cA)$ is a strictly smaller object: the finality
property survives as a \emph{terminal projection} from the ordered
finality, not as an independent universal property on the symmetric
side, because averaging annihilates the non-symmetrisable half-braidings
that are the content of the braided Hopf structure on~$U_\cA$. The
final-object and initial-object readings both live on the ordered
$E_1$-side; their symmetric shadows are weaker.
\end{remark}

\subsection{Five recovery functors and the compatibility hexagon}
\label{subsec:hdm-recovery-functors}

The six-slot recovery of
Proposition~\ref{prop:hdm-u-a-recoverability} is resolved here at the
level of $(\infty,1)$-functors. For each slot of
$\mathcal H(T)=(\cA,\cA^!,\cC,r_\cA(z),\Theta_\cA,\nabla^{\mathrm{hol}})$
there is a projection functor $P_i$ out of the boundary-bulk algebra
$U_\cA=\cA\bowtie\cA^!_\infty$; the six projections are mutually
compatible along a hexagon of natural transformations. The Verdier
projection $P_{\cA^!}$ and the bulk-line comparison $P_\cC$ lose
measurable cohomology classes off the named comparison surfaces; those
classes are computed explicitly.

Write $\mathrm{Alg}_{E_1}^{\mathrm{ch,curved}}$ for the
$(\infty,1)$-category of curved $A_\infty$-chiral algebras on~$X$,
$\mathrm{CoAlg}_{E_1}^{\mathrm{ch,cocurved}}$ for its coderived mirror,
$\mathrm{Bimod}^{\mathrm{ch}}$ for chiral bimodules,
$\mathrm{SchGr}(X)$ for spectral $R$-matrix data on~$X$
(holomorphic sections of
$\mathfrak g\otimes\mathfrak g\otimes \Omega^1_{X\times X}(\Delta_*)$
satisfying CYBE), $\mathrm{MC}(\mathfrak g^{\mathrm{mod}}_\cA)$ for
Maurer--Cartan elements of the modular convolution dGLA, and
$\mathrm{FlatConn}(\overline{\cM}_{g,n})$ for flat
holomorphic connections over the modular multi-disc. Each target is a
presentable $\infty$-category.

\begin{theorem}[Five recovery functors out of $U_\cA$; \ClaimStatusConditional]
\label{thm:hdm-five-recovery-functors}
\index{recovery functors!five projections from $U_\cA$|textbf}
\index{Drinfeld double!recovery functors|textbf}
Let $\cA$ be a modular Koszul chiral algebra with Drinfeld double
$U_\cA=\cA\bowtie\cA^!_\infty$. The five non-identity slots of
$\mathcal H(T)$ are the images of $U_\cA$ under five explicit
$(\infty,1)$-functors:
\begin{enumerate}[label=\textup{(P\arabic*)}]
\item \emph{Boundary projection} (tensor with counit, left adjoint):
\[
P_\cA\;:=\;U_\cA\otimes^{\mathbb L}_{\cA^!_\infty}\mathbb C
\;:\;\mathrm{Alg}_{E_1}^{\mathrm{ch,curved}}(U_\cA)\longrightarrow
\mathrm{Alg}_{E_1}^{\mathrm{ch,curved}}(\cA).
\]
This is a monadic quotient: $\cA^!_\infty$ acts on $U_\cA$ from the
right by the braided product, and $P_\cA$ is the coequalizer against
the augmentation $\epsilon\colon \cA^!_\infty\to\mathbb C$. It is
unconditional for every modular Koszul~$\cA$.
\item \emph{Verdier/Koszul projection} (derived Hom with counit, right
adjoint to $P_\cA$):
\[
P_{\cA^!}\;:=\;\mathrm{RHom}_{\cA}(\mathbb C,U_\cA)
\;:\;\mathrm{Alg}_{E_1}^{\mathrm{ch,curved}}(U_\cA)\longrightarrow
\mathrm{CoAlg}_{E_1}^{\mathrm{ch,cocurved}}(\cA^!_\infty).
\]
This is a monadic localization: fibres are $\cA$-bimodule resolutions
of the counit $\cA\to\mathbb C$. The strict cohomological shadow
$\cA^!_\infty\twoheadrightarrow\cA^!$ is an equivalence on the
modular Koszul locus (Upgrade~(U2),
Proposition~\ref{prop:hdm-upgrade-rectification}) and
lossy off it.
\item \emph{Bulk projection} (chiral Hochschild centre):
\[
P_\cC\;:=\;\cZ^{\mathrm{der}}_{\mathrm{ch}}(U_\cA)
\;\simeq\;\mathrm{HH}^{\bullet}_{\mathrm{ch}}(U_\cA,U_\cA)
\;:\;\mathrm{Alg}_{E_1}^{\mathrm{ch,curved}}(U_\cA)\longrightarrow
\mathrm{Alg}_{E_2}^{\mathrm{ch,SC}}(\cC).
\]
This is a limit (end) rather than a monadic quotient: it is the
derived endomorphism object of $U_\cA$ as a chiral bimodule over
itself, and the $E_2$-structure arises from the
$\mathrm{Ass}^{\mathrm{ch}}\hookrightarrow\mathrm{SC}^{\mathrm{ch,top}}$
inclusion. The restriction map $\rho_\partial\colon
P_\cC\to P_\cA$ is induced by evaluation at the unit
$\mathbb C\to U_\cA$.
\item \emph{$R$-matrix projection} (collision residue, localization at
$\Delta$):
\[
P_r\;:=\;\operatorname{Res}^{\mathrm{coll}}_{0,2}\circ P_\Theta
\;:\;\mathrm{Alg}_{E_1}^{\mathrm{ch,curved}}(U_\cA)\longrightarrow
\mathrm{SchGr}(X).
\]
Equivalently, $P_r$ reads off the quasitriangular structure of
$U_\cA$: the universal $R$-matrix of the braided product $\bowtie$ is
the singular part of the pairing
$\cA\otimes\cA^!_\infty\to\mathbb C[z^{-1}]$ along the
diagonal~$\Delta\subset X\times X$. It factors through the depth-$2$
collision stratum of $\overline{\cM}_{0,2}$.
\item \emph{MC projection} (gauge class of the universal obstruction):
\[
P_\Theta\;:=\;[\Theta_{U_\cA}]
\;:\;\mathrm{Alg}_{E_1}^{\mathrm{ch,curved}}(U_\cA)\longrightarrow
\mathrm{MC}(\mathfrak g^{\mathrm{mod}}_\cA)/\mathrm{gauge}.
\]
This is a colimit functor: $\Theta_{U_\cA}$ is the canonical MC
element of the convolution dGLA
$\mathrm{Conv}(B^{\mathrm{ord}}(U_\cA),U_\cA)$, and its descent under
the averaging map $\mathrm{av}\colon\mathfrak g^{E_1}\to
\mathfrak g^{\mathrm{mod}}$ is $\Theta_\cA=D_\cA-d_0$
(Upgrade~(U1), Theorem~\ref{thm:mc2-bar-intrinsic}).
\item \emph{Shadow-connection projection} (Gauss--Manin on periods):
\[
P_{\nabla^{\mathrm{hol}}}\;:=\;
\nabla^{\mathrm{GM}}\bigl(H^\bullet_{\mathrm{per}}(U_\cA/\overline{\cM}_{g,n})\bigr)
\;:\;\mathrm{Alg}_{E_1}^{\mathrm{ch,curved}}(U_\cA)\longrightarrow
\mathrm{FlatConn}\bigl(\overline{\cM}_{g,n}\bigr).
\]
Equivalently, $P_{\nabla^{\mathrm{hol}}}=d-\operatorname{Sh}_{g,n}
(P_\Theta)$: the shadow connection is the Gauss--Manin differential
on the period complex of $U_\cA$ over the moduli multi-disc, and its
flatness is the $(g,n)$-shadow of the MC equation
$P_\Theta^2+\tfrac12[P_\Theta,P_\Theta]=0$
\textup{(}Theorem~\textup{\ref{thm:sphere-reconstruction}}\textup{)}.
\end{enumerate}
The identity slot $P_\cA$ is unconditional; $P_\Theta$, $P_r$, and
$P_{\nabla^{\mathrm{hol}}}$ are unconditional for every modular
Koszul~$\cA$. The remaining two, $P_{\cA^!}$ and $P_\cC$, carry the
non-trivial obstruction classes $\eta_4$ and $q_\mathsf{M}$ of
Theorem~\ref{thm:hdm-koszul-loss-classes} below.
\end{theorem}

\begin{proof}
$P_\cA$ is the tensor-with-counit left adjoint of the braided product
$-\otimes^{\mathbb L}\cA^!_\infty$, unconditional because
$\cA\hookrightarrow U_\cA$ is split by $\mathrm{id}\otimes\epsilon$
as a chain map. $P_{\cA^!}$ is its right adjoint by the standard
$\otimes/\mathrm{Hom}$ adjunction; the monadic-localization criterion is
Beck's theorem applied to the $\cA$-module comonad comparison.
$P_\cC$ is the Hochschild centre of $U_\cA$; the equivalence
$\cZ^{\mathrm{der}}_{\mathrm{ch}}(U_\cA)\simeq
\cZ^{\mathrm{der}}_{\mathrm{ch}}(\cA)$ is the open/closed theorem
(Theorem~\ref{thm:thqg-swiss-cheese}), using that
$\cA\hookrightarrow U_\cA$ is Morita-trivial on the
bimodule side (the Drinfeld double acts on itself by
left multiplication, and chiral Hochschild of a braided product over
its factor is the Hochschild of the boundary factor). $P_r$ is the
collision-residue projection of the seven-way master theorem
(Theorem~\ref{thm:hdm-seven-way-master}), restricted to $U_\cA$ via
the $\bowtie$ product. $P_\Theta$ is Upgrade~(U1),
Proposition~\ref{prop:hdm-upgrade-theta}. The shadow-connection
projection $P_{\nabla^{\mathrm{hol}}}$ is the $(g,n)$-sector
restriction of $\operatorname{Sh}\circ P_\Theta$
(Section~\ref{sec:hdm-collision-residue}). The functor types are
read from the constructions: $P_\cA$ is a coequalizer (monadic
quotient), $P_{\cA^!}$ is an equalizer (monadic localization),
$P_\cC$ is an end (limit), $P_r$ and $P_{\nabla^{\mathrm{hol}}}$ are
collision and Gauss--Manin localizations, $P_\Theta$ is a filtered
colimit over approximating generating sets of $U_\cA$.
\end{proof}

\begin{theorem}[Compatibility hexagon of the six projections;
\ClaimStatusConditional]
\label{thm:hdm-hexagon}
\index{compatibility hexagon!six recovery functors|textbf}
\index{recovery functors!hexagon of natural transformations|textbf}
The six functors of Theorem~\ref{thm:hdm-five-recovery-functors}
fit into a hexagon of natural transformations
\[
\begin{array}{ccc}
P_\cA & \xleftarrow{\;\alpha\;} & P_\cC \\
\uparrow \beta & & \downarrow \gamma \\
P_{\cA^!} & & P_r \\
\uparrow \delta & & \downarrow \epsilon \\
P_\Theta & \xrightarrow{\;\zeta\;} & P_{\nabla^{\mathrm{hol}}}
\end{array}
\]
in which the natural transformations are:
$\alpha$ = Swiss-cheese restriction
$\rho_\partial\colon P_\cC\to P_\cA$, the unit of the
centre/boundary adjunction;
$\beta$ = counit of the Verdier/Koszul adjunction
$P_{\cA^!}\to P_\cA$;
$\gamma$ = collision-residue projection
$P_\cC\twoheadrightarrow P_r$ via
$\operatorname{Res}^{\mathrm{coll}}_{0,2}$ on
$\mathrm{HH}^\bullet_{\mathrm{ch}}$;
$\delta$ = MC-to-coalgebra assembly
$P_\Theta\to P_{\cA^!}$ from bar duality
$\Theta_\cA\leftrightarrow$ codifferential on $B^{\mathrm{ord}}(\cA)$;
$\epsilon$ = shadow-absorption
$P_r\to P_{\nabla^{\mathrm{hol}}}$ (the sphere limit
$\lim_{z\to0}z\cdot r_\cA(z)$ is the $T$-coefficient of
$\nabla^{\mathrm{hol}}_{0,n}$);
$\zeta$ = shadow map
$\operatorname{Sh}_{g,n}\colon P_\Theta\to P_{\nabla^{\mathrm{hol}}}$.
The hexagon commutes up to coherent homotopy on the modular Koszul
boundary-linear locus: each of the two three-vertex paths from
$P_\Theta$ to $P_{\nabla^{\mathrm{hol}}}$
\textup{(}$P_\Theta\to P_{\cA^!}\to P_\cA\leftarrow P_\cC\to
P_r\to P_{\nabla^{\mathrm{hol}}}$ and $P_\Theta\to
P_{\nabla^{\mathrm{hol}}}$\textup{)} equals the canonical shadow
class
$[\Theta_{U_\cA}]\in H^1(\mathfrak g^{\mathrm{mod}}_{U_\cA})$.
\end{theorem}

\begin{proof}
$\alpha$ is the evaluation-at-unit map
$\mathrm{HH}^\bullet_{\mathrm{ch}}(U_\cA,U_\cA)\to U_\cA$,
followed by $P_\cA$-projection; the restriction is
$\rho_\partial\colon \cC\to\cA$ of
Section~\ref{sec:hdm-thesis}. $\beta$ is the counit of the
$P_\cA\dashv P_{\cA^!}$ adjunction applied fibrewise. $\gamma$ is the
codimension-$1$ boundary-stratum residue on
$\mathrm{HH}^\bullet_{\mathrm{ch}}$, equal to the
$\operatorname{Res}^{\mathrm{coll}}_{0,2}$ used to define $r_\cA(z)$
in Section~\ref{sec:hdm-collision-residue}. $\delta$ is the universal
bar construction: $\Theta_\cA$ defines a codifferential on
$B^{\mathrm{ord}}(\cA)$ whose cohomology is $\cA^!$ on the Koszul
locus. $\epsilon$ is the genus-zero, two-point specialization
$\nabla^{\mathrm{hol}}_{0,2}|_{z\to 0}=r_\cA(z)/(z-w)$ (the KZ
connection has $r$ as its residue). $\zeta$ is Upgrade~(U1) together
with $\operatorname{Sh}_{g,n}$ of slot~(6). Commutativity is checked
on three pasted triangles: the $P_\cA$-$P_\cC$-$P_{\cA^!}$ triangle
is bar--cobar inversion
(Theorem~\ref{thm:bar-cobar-inversion-qi}); the
$P_r$-$P_\cC$-$P_\Theta$ triangle is the seven-way master theorem
(Theorem~\ref{thm:hdm-seven-way-master}); the
$P_\Theta$-$P_{\nabla^{\mathrm{hol}}}$-$P_r$ triangle is sphere
reconstruction
(Theorem~\ref{thm:sphere-reconstruction}). Pasting the three
triangles along their shared $P_\Theta$ vertex gives the hexagon. The
shadow-class identity follows because each two-path composition
descends $U_\cA$ through the same universal MC element; the coherent
homotopy between them is the bar--cobar inversion homotopy rescaled
by the shadow map.
\end{proof}

\begin{theorem}[Koszul-loss classes of $P_{\cA^!}$ and $P_\cC$;
\ClaimStatusConditional]
\label{thm:hdm-koszul-loss-classes}
\index{Koszul loss classes!$\eta_4$ and $q_{\mathsf M}$|textbf}
\index{recovery functors!fibre obstructions|textbf}
On the modular Koszul boundary-linear locus the total recovery
functor
\[
P_\bullet\;:=\;(P_\cA,P_{\cA^!},P_\cC,P_r,P_\Theta,P_{\nabla^{\mathrm{hol}}})
\;:\;\mathrm{Alg}_{E_1}^{\mathrm{ch,curved}}(U_\cA)\longrightarrow
\prod_{i}R_i
\]
is \emph{conservative}: an arrow $f\colon U_\cA\to U_\cA'$ is an
equivalence iff each $P_i(f)$ is. Off the Koszul locus two specific
classes obstruct conservativity.
\begin{enumerate}[label=\textup{(L\arabic*)}]
\item The fibre of the Verdier projection
$\mathrm{fib}(\cA^!_\infty\twoheadrightarrow\cA^!)$ is
controlled in $\pi_1$ by the admissible
$L_{-1/2}(\mathfrak{sl}_2)$ obstruction
\[
[\eta_4]\;\in\;H^1\bigl(\mathrm{fib}(P_{\cA^!})\bigr)
\;\simeq\;H^4\bigl(\bar B^{\mathrm{ch}}(L_{-1/2}(\mathfrak{sl}_2))\bigr),
\]
the fourth bar class that Upgrade~(U2) fails to kill.
Explicitly, $\eta_4$ is the chain-level cocycle
\[
\eta_4\;=\;
\sum_{i<j<k<\ell}
\operatorname{tr}\bigl(e_{ij}e_{jk}e_{k\ell}e_{\ell i}\bigr)\,
d\log(z_{ij})\wedge d\log(z_{jk})\wedge
d\log(z_{k\ell})\wedge d\log(z_{\ell i})
\]
in $C^4(\bar B^{\mathrm{ch}}(L_{-1/2}(\mathfrak{sl}_2)))$, where the
Arnold relation
$\eta_{ij}\wedge\eta_{jk}+\eta_{jk}\wedge\eta_{ki}+\eta_{ki}\wedge\eta_{ij}=0$
fails to close $\eta_4$ at the admissible level
$k=-1/2$ but would close it at generic level
(Example~\ref{ex:admissible-sl2-failure}).
$P_{\cA^!}$ is lossless precisely on the locus $[\eta_4]=0$.
\item The fibre of the bulk/line comparison
$\mathrm{fib}(P_\cC\to P_{\cA^!\mathrm{-mod}})$ on nonlinear
class~$\mathsf M$ is controlled in $\pi_2$ by the quartic contact
class
\[
q_\mathsf{M}\;:=\;\frac{10}{c(5c+22)}\;\in\;
H^2\bigl(\mathrm{fib}(P_\cC)\bigr),
\]
equal to the reciprocal of the Zamolodchikov norm
$\langle\Lambda|\Lambda\rangle=c(5c+22)/10$ of the quasiprimary
$\Lambda={:}TT{:}-(3/10)\partial^2 T$. This is the Virasoro
shadow entry $S_4(\mathrm{Vir}_c)=10/[c(5c+22)]$
(Section~\ref{sec:hdm-thesis}, and the Zamolodchikov-norm line of
the constants table in the constitution). $P_\cC$ agrees with the
line-model presentation on class~$\mathsf M$ precisely on the locus
$q_\mathsf{M}=0$, which is vacuous for positive central charge:
on class~$\mathsf{M}$ the obstruction is always non-zero for
$c>0$, so the bulk/line projection is strictly lossy there.
\end{enumerate}
The Koszul-locus conservativity is unconditional and sharp: outside
the conservativity locus, the class $(\eta_4,q_\mathsf{M})\in
H^1\oplus H^2$ is the total obstruction to reconstructing $U_\cA$
from the six-tuple.
\end{theorem}

\begin{proof}
Conservativity on the Koszul boundary-linear locus is the conjunction
of Upgrades~(U1)--(U3)
(Propositions~\ref{prop:hdm-upgrade-theta},
\ref{prop:hdm-upgrade-rectification},
\ref{prop:hdm-upgrade-bulk-slot}):
on that locus $P_{\cA^!}$ becomes an equivalence and $P_\cC$ agrees
with the line-model projection, so the total functor $P_\bullet$ has
trivial fibre. For (L1), the admissible $L_{-1/2}(\mathfrak{sl}_2)$
computation of Example~\ref{ex:admissible-sl2-failure} shows that
$H^4(\bar B^{\mathrm{ch}}(L_{-1/2}(\mathfrak{sl}_2)))$ contains the
non-zero class $\eta_4$, represented by the four-cycle tetrahedral
Arnold product; by the strict-edge obstruction of
Upgrade~(U2), this class survives in the fibre
$\mathrm{fib}(\cA^!_\infty\to\cA^!)$. For (L2), the chain-level
representative of $q_\mathsf{M}$ is the quartic contact cocycle on
the Virasoro ${:}TT{:}$ tower; it computes the Virasoro shadow entry
$S_4=10/[c(5c+22)]$ directly (Section~\ref{sec:hdm-thesis}). The
Zamolodchikov normalization
$\langle\Lambda|\Lambda\rangle=c(5c+22)/10$ is the chain-level dual
of this class, and its non-vanishing for $c>0$ is the precise
obstruction to the boundary-linear reduction on class~$\mathsf{M}$.
Sharpness follows from the two flagship chain-level examples being
the only surviving obstructions in the modular Koszul stratification
(Theorem~\ref{thm:dichotomy-chain-level}).
\end{proof}

\begin{remark}[Six-projection recovery map]
\label{rem:hdm-recovery-status-map}
The six recovery functors sort into three strata:
$P_\cA$ and $P_\Theta$ are unconditionally lossless on the modular
Koszul locus (tautological and universal); $P_r$ and
$P_{\nabla^{\mathrm{hol}}}$ are lossless on the modular Koszul
boundary-linear locus (seven-way theorem and sphere reconstruction
pin them); $P_{\cA^!}$ and $P_\cC$ are lossless on the Koszul
boundary-linear locus only and carry the explicit fibre classes
$[\eta_4]$ and $q_\mathsf{M}$ off it. The hexagon of
Theorem~\ref{thm:hdm-hexagon} commutes on the Koszul boundary-linear
locus unconditionally, and commutes up to
$[\eta_4]\oplus q_\mathsf{M}$ elsewhere; the correction term is the
unique obstruction to extending the six-slot reconstruction from the
named lanes to the full standard landscape. On class~$\mathsf{G}$
(Heisenberg) both classes vanish: $\eta_4=0$ because Heisenberg is
self-Koszul and $q_\mathsf{M}=0$ because the shadow tower terminates
at $r_{\max}=2$, so there is no quartic contact class. On
class~$\mathsf{L}$ (affine Kac--Moody at generic level) $\eta_4=0$
by Casimir closure; the admissible level $L_{-1/2}(\mathfrak{sl}_2)$
is the unique known chain-level failure. On class~$\mathsf{M}$
(Virasoro) $q_\mathsf{M}\neq 0$ for $c>0$, so the bulk/line
reduction is always strictly lossy; the hexagon commutes up to that
class only.
\end{remark}

\begin{lemma}[Tensor-product K\"unneth law for the Koszul-loss classes;
\ClaimStatusConditional]
\label{lem:hdm-koszul-loss-kunneth}
\index{Koszul loss classes!K\"unneth law|textbf}
\index{chiral tensor product!Koszul-loss K\"unneth|textbf}
Let $\cA,\cB$ be augmented chiral algebras on $X$ in the modular
landscape, with $\bar B^{\mathrm{ch}}(-)$ their reduced ordered bar
complexes on $\mathrm{Conf}_\bullet(X)$. The Ben-Zvi--Nadler
factorisation K\"unneth map
\[
\kappa_{\cA,\cB}\colon
\bar B^{\mathrm{ch}}(\cA)\otimes \bar B^{\mathrm{ch}}(\cB)
\;\xrightarrow{\;\simeq\;}\;
\bar B^{\mathrm{ch}}(\cA\otimes\cB)
\]
is a quasi-isomorphism of curved $E_1$-coalgebras; composed with the
Verdier dual it induces
\[
H^p\bigl(\bar B^{\mathrm{ch}}(\cA)\bigr)
\otimes
H^q\bigl(\bar B^{\mathrm{ch}}(\cB)\bigr)
\;\xrightarrow{\;\sim\;}\;
H^{p+q}\bigl(\bar B^{\mathrm{ch}}(\cA\otimes\cB)\bigr)
\]
(Kazhdan--Lusztig \textup{1994} tensor-product compatibility for
vertex categories; in the $(\infty,1)$-lane, Lurie
$\mathrm{HA}.5.5.3$ for fibres of exact functors between stable
$\infty$-categories). The tetrahedral class $[\eta_4]$ and the
quartic contact class $q_\mathsf{M}$ of
Theorem~\ref{thm:hdm-koszul-loss-classes} obey the additive
tensor-product law
\[
\bigl([\eta_4],\,q_\mathsf{M}\bigr)_{\cA\otimes\cB}
\;=\;
\Bigl(
\pi_\cA^\ast[\eta_4]_\cA
+\pi_\cB^\ast[\eta_4]_\cB
,\;
\frac{10}{(c_\cA+c_\cB)\bigl(5(c_\cA+c_\cB)+22\bigr)}
\Bigr),
\]
where $\pi_\cA,\pi_\cB$ are the two K\"unneth projectors onto
$(4,0)$ and $(0,4)$ bar-strata. Equivalently, the tensor-product
fibre splits
\[
\mathrm{fib}\bigl(P_{(\cA\otimes\cB)^!}\bigr)
\;\simeq\;
\mathrm{fib}(P_{\cA^!})\,\oplus\,\mathrm{fib}(P_{\cB^!})
\qquad\text{at }\pi_1,\ \text{in the Koszul-loss stratum,}
\]
so the loss class is an \emph{additive coproduct} of its factor
losses, not a cup. On $\pi_2$ the Virasoro contact class is
additive on central charge and rational-quadratic on
$(q_\mathsf{M}^{(\cA)},q_\mathsf{M}^{(\cB)})$ through the
inverse-shadow substitution $c_i = c(q_\mathsf{M}^{(i)})$, the
unique positive root of $5c_i^2+22c_i - 10/q_\mathsf{M}^{(i)}=0$.
\end{lemma}

\begin{proof}
The K\"unneth quasi-isomorphism $\kappa_{\cA,\cB}$ is the
factorisation-homology avatar of
Kazhdan--Lusztig~\textup{1994 §7}: augmented chiral algebras tensor
as coaugmented $E_1$-coalgebras, and the reduced ordered bar is
free as cograded object on $s^{-1}(\bar\cA)\otimes s^{-1}(\bar\cB)$,
so $\bar B^{\mathrm{ch}}(\cA\otimes\cB) = T^c(s^{-1}\overline{\cA\otimes\cB})$
receives $\kappa_{\cA,\cB}$ from the shuffle map on tensor algebras.
Both sides are cofibrant (free cograded) and the differential
$D_{\cA\otimes\cB} = D_\cA\otimes 1 + 1\otimes D_\cB + \mu_{\mathrm{ch}}$
has $\mu_{\mathrm{ch}}$-correction homotopic to zero on the
augmented bar (Ben-Zvi--Nadler factorisation K\"unneth), giving
$\kappa_{\cA,\cB}$ as a chain-level quasi-isomorphism. The induced
cohomology K\"unneth is the split-$\mathrm{Tor}$ statement over a
field of characteristic zero.

For $[\eta_4]$: the tetrahedral cocycle
$\eta_4 = \sum_{i<j<k<\ell}\operatorname{tr}(e_{ij}e_{jk}e_{k\ell}e_{\ell i})\,
d\log z_{ij}\wedge d\log z_{jk}\wedge d\log z_{k\ell}\wedge d\log z_{\ell i}$
of Theorem~\ref{thm:hdm-koszul-loss-classes} is supported on
configurations where all four generators $e_{**}$ belong to a
\emph{single} tensor factor: the Arnold relation and the trace form
are factor-local. Under $\kappa_{\cA,\cB}$ the class projects
onto the pure K\"unneth strata $H^4(\bar B^{\mathrm{ch}}(\cA))\otimes
H^0 \oplus H^0 \otimes H^4(\bar B^{\mathrm{ch}}(\cB))
= H^4(\cA)\oplus H^4(\cB)$. The augmentation counits $H^0 = k$ of
both factors absorb the opposite-factor index, yielding exactly
$\pi_\cA^\ast[\eta_4]_\cA + \pi_\cB^\ast[\eta_4]_\cB$. No mixed
$(p,q)$ term with $1\le p,q\le 3$ contributes, because the
tetrahedral four-cycle on mixed generators has zero trace on
$\cA\otimes\cB$-indexed crossings: the Arnold-relation closure
decouples across tensor factors, leaving a boundary rather than a
cycle. Cup of two $[\eta_4]$'s would live in $H^8$ and is not
the $\pi_1$-fibre class.

For $q_\mathsf{M}$: two commuting Virasoro stress tensors
$T_\cA,T_\cB$ have $T_{\cA\otimes\cB} = T_\cA\otimes 1 + 1\otimes
T_\cB$, and the total central charge is $c_{\cA\otimes\cB} = c_\cA
+ c_\cB$ by the ${TT}$-OPE double pole (Zamolodchikov norm line of
the constants table in \textsc{concordance.tex}). The quasiprimary
$\Lambda = {:}TT{:} - (3/10)\partial^2 T$ of Virasoro has
Zamolodchikov norm $\langle\Lambda|\Lambda\rangle = c(5c+22)/10$,
whose reciprocal is $q_\mathsf{M}$. Substituting
$c\mapsto c_\cA+c_\cB$ gives the asserted formula. The
inverse-shadow substitution $c_i(q_\mathsf{M}^{(i)}) =
(-22+\sqrt{484+200/q_\mathsf{M}^{(i)}})/10$ expresses
$q_\mathsf{M}^{\mathrm{tot}}$ as a rational function of the factor
losses; the map is neither a coproduct nor a cup on the
$q_\mathsf{M}$-scale, but becomes linear-additive on the
central-charge scale.

The coproduct of fibres
$\mathrm{fib}(P_{(\cA\otimes\cB)^!}) \simeq
\mathrm{fib}(P_{\cA^!})\oplus\mathrm{fib}(P_{\cB^!})$ at $\pi_1$
is the K\"unneth direct sum applied to the Verdier fibre, using
that $\cA^!_\infty\otimes\cB^!_\infty = (\cA\otimes\cB)^!_\infty$
on augmented objects (Koszul-dual monoidal, Positselski
\textup{2011 §6.7}).
\end{proof}

\begin{example}[Three tensor products of landscape flagships;
\ClaimStatusProvedHere]
\label{ex:hdm-koszul-loss-kunneth-examples}
\index{Koszul loss classes!tensor-product examples|textbf}
Three explicit evaluations of
Lemma~\ref{lem:hdm-koszul-loss-kunneth} on class-$\mathsf G$,
class-$\mathsf L$, class-$\mathsf M$ pairings:
\begin{enumerate}[label=\textup{(\roman*)}]
\item \emph{Heisenberg tensor Virasoro}
$\cH_k\otimes\mathrm{Vir}_c$: $[\eta_4]_{\cH_k}=0$ (self-Koszul),
$[\eta_4]_{\mathrm{Vir}_c}=0$ (no admissible $\mathfrak{sl}_2$
structure), so $[\eta_4]_{\mathrm{tot}}=0$; the total central
charge is $c_{\mathrm{tot}} = 1+c$ (Heisenberg contributes
$c_{\cH_k}=1$ per free boson), so $q_\mathsf{M}^{\mathrm{tot}} =
10/[(1+c)(5(1+c)+22)] = 10/[(1+c)(5c+27)]$.
\item \emph{Two affine Kac--Moodys}
$\widehat\fg_k\otimes\widehat\fh_l$: both factors have
$[\eta_4]=0$ at generic level (Casimir closure), so
$[\eta_4]_{\mathrm{tot}}=0$; the Segal--Sugawara construction
gives $c_{\widehat\fg_k} = k\dim\fg/(k+h^\vee_\fg)$ and
likewise for $\fh$. The admissible level $k_\fg = -h^\vee_\fg/2 +
\mathrm{rational}$ in \emph{one} factor activates the tetrahedral
class on that factor only: $[\eta_4]_{\mathrm{tot}} =
\pi_{\widehat\fg}^\ast[\eta_4]_{L_{-1/2}(\mathfrak{sl}_2)}$ when
$\fg=\mathfrak{sl}_2,k=-1/2$ and $\fh,l$ generic.
\item \emph{Two Virasoros}
$\mathrm{Vir}_{c_1}\otimes\mathrm{Vir}_{c_2}$: $[\eta_4]=0$ on both
factors, total central charge $c_1+c_2$,
\[
q_\mathsf{M}^{\mathrm{tot}}\;=\;
\frac{10}{(c_1+c_2)\bigl(5(c_1+c_2)+22\bigr)}.
\]
At $c_1=c_2=1$ this evaluates to $q_\mathsf{M}^{\mathrm{tot}} =
10/[2\cdot 32] = 5/32$, whereas the individual factor losses are
$q_\mathsf{M}^{(1)}=q_\mathsf{M}^{(2)}=10/27$. The product
$q_\mathsf{M}^{(1)} q_\mathsf{M}^{(2)} = 100/729$ differs from
$q_\mathsf{M}^{\mathrm{tot}}$ both in magnitude ($\approx 0.137$
vs $\approx 0.156$) and scaling, verifying that the law is
\emph{not} multiplicative on the $q_\mathsf{M}$-scale. At $c_1=26$
(bosonic string) and $c_2=-26$ (ghost cancellation) the total
central charge is zero and $q_\mathsf{M}^{\mathrm{tot}}$ diverges,
the classical-BRST critical-dimension degeneration at the
Zamolodchikov-norm pole.
\end{enumerate}
\end{example}

\subsection{Class-$\mathsf M_\infty$ fibre: the Pope--Romans--Shen tower}
\label{sec:hdm-class-m-infty-fibre}

Class~$\mathsf M$ exhausts only shadow depth $r_{\max}=4$. The
Virasoro subclass terminates at the quartic $\Lambda$; the
single-spin-$2$ quasi-primary has one Zamolodchikov companion, and
the fibre carries one class $q_\mathsf M$. The limit class
$\mathsf M_\infty$ is $\mathcal W_{1+\infty}[c]$, shadow depth
$r_{\max}=\infty$: generators $W^{(k)}$ of conformal weight $k$ for
every $k \geq 1$. Each generator contributes its own Zamolodchikov
quasi-primary $\Lambda_k = {:}W^{(k)}W^{(k)}{:} + (\text{lower-weight
covariant corrections})$, and the Koszul-loss fibre
$\mathrm{fib}(P_\cC)$ receives one contact class per $k \geq 4$.
Theorem~\ref{thm:hdm-koszul-loss-classes} detects only
$q_\mathsf M = q_\mathsf M^{(4)}$; the full obstruction is the
Pope--Romans--Shen tower.

\begin{theorem}[Class-$\mathsf M_\infty$ fibre; \ClaimStatusConditional]
\label{thm:hdm-class-m-infty-fibre}
\index{Koszul loss classes!class $\mathsf M_\infty$|textbf}
\index{W1infinity@$\mathcal W_{1+\infty}$!Koszul loss fibre}
\index{Pope--Romans--Shen tower}
For $\cA = \mathcal W_{1+\infty}[c]$ on a modular Koszul curve, the
fibre of the bulk/line comparison
$\mathrm{fib}(P_\cC\to P_{\cA^!\mathrm{-mod}})$ is a $\mathrm{Pro}$-object
$\bigoplus_{k\geq 4}\mathrm{fib}^{(k)}$ whose $k$-th component is
controlled by the quartic-to-$k$ Zamolodchikov contact class
\begin{equation}\label{eq:hdm-minfty-fibre-generator}
q_\mathsf M^{(k)}(c)
\;:=\;\frac{1}{\langle\Lambda_k|\Lambda_k\rangle(c)}
\;\in\;H^{k-2}\bigl(\mathrm{fib}^{(k)}(P_\cC)\bigr),
\qquad k\geq 4,
\end{equation}
where $\Lambda_k$ is the Pope--Romans--Shen spin-$k$
quasi-primary~\cite{PopeRomansShen90}. The full fibre class on
class-$\mathsf M_\infty$ is the countable direct sum
\begin{equation}\label{eq:hdm-minfty-fibre-total}
\mathrm{fib}\bigl(P_\cC\bigm|_{\mathsf M_\infty}\bigr)
\;=\;\bigoplus_{k\geq 4}\, q_\mathsf M^{(k)}(c)\cdot\langle W^{(k)}\rangle
\end{equation}
in $\mathrm{Pro}(\mathrm{Vect}_{\geq 0})$, with Hilbert series
$\chi_{\mathrm{fib}}(q) = \prod_{k\geq 4}(1-q^k)^{-1}$ (MacMahon-type
generating function, Schiffmann--Vasserot~\cite{SchiffmannVasserot13}
CoHA character on $\mathbb C^2$). The low-weight initial data are
\begin{equation}\label{eq:hdm-minfty-low-weights}
q_\mathsf M^{(4)}(c)=\frac{10}{c(5c+22)},
\qquad
q_\mathsf M^{(5)}(c)=\frac{60}{c^{2}(5c+22)(7c+114)},
\qquad
q_\mathsf M^{(6)}(c)=\frac{21}{c^{3}(5c+22)(7c+114)(2c+15)},
\end{equation}
\begin{equation}\label{eq:hdm-minfty-q7q8}
q_\mathsf M^{(7)}(c)=\frac{5040}{c^{4}(5c+22)(7c+114)(2c+15)(11c+232)},
\qquad
q_\mathsf M^{(8)}(c)=\frac{4320}{c^{5}(5c+22)(7c+114)(2c+15)(11c+232)(13c+350)},
\end{equation}
the $k$-th entry $q_\mathsf M^{(k)}$ being $2\kappa_k^2/\langle\Lambda_k|\Lambda_k\rangle_{\mathrm{PRS}}$
with $\kappa_k=c/k$ and with Pope--Romans--Shen norm
\[
\langle\Lambda_k|\Lambda_k\rangle_{\mathrm{PRS}}
\;=\;\frac{2c^{k-1}}{k^{2}\,k!}\cdot\prod_{j=2}^{k}\!\bigl((2j-1)c+R_j\bigr),
\qquad
R_2=22,\;R_3=114,\;R_4=15\cdot(2/1),\;R_5=232,\;R_6=350,
\]
the integer residues $R_j$ of the weight-$2j$ Kac-determinant descent
(Pope--Romans--Shen~\cite{PopeRomansShen90} \S4; Bouwknegt--Schoutens
\S6.2; cross-tabulated in \S\ref{sec:higher-w-shadow-data}).
The generating function
$F(c) := \sum_{k\geq 4} q_\mathsf M^{(k)}(c)$ is Borel-summable in the
coupling $1/c$ for $|c|>1$, with Borel transform
\[
\widehat F(\zeta)\;=\;\sum_{k\geq 4}\frac{\zeta^{k-4}}{(k-4)!}
\cdot q_\mathsf M^{(k),\mathrm{top}},
\]
where $q_\mathsf M^{(k),\mathrm{top}}$ is the leading-in-$1/c$
coefficient of the $k$-th Zamolodchikov norm.
\end{theorem}

\begin{proof}
Existence of $\Lambda_k$ is Pope--Romans--Shen~\cite{PopeRomansShen90}:
the $\mathcal W_{1+\infty}[c]$ self-OPE of each generator $W^{(k)}$
contains a leading pole $\langle W^{(k)}|W^{(k)}\rangle \cdot z^{-2k}$
whose normalisation $\kappa_k(c) = c/k$ is the $\Walg_N$ channel-refined
kappa entry of Proposition~\ref{prop:higher-w-kappa-matrix}
(\texttt{w\_algebras\_deep.tex:4099}). The associated Zamolodchikov
quasi-primary $\Lambda_k$ is the lowest-weight composite
${:}W^{(k)}W^{(k)}{:}$ minus the descent $\tfrac{k}{2k-1}\partial^{2}W^{(2k-2)}$
(the $k=2$ case is the Virasoro $\Lambda_2 = \Lambda = {:}TT{:} - (3/10)\partial^2 T$).

For (\ref{eq:hdm-minfty-fibre-generator}): repeating the chain-level
cocycle construction of Theorem~\ref{thm:hdm-koszul-loss-classes}(L2)
with $W^{(k)}$ in place of $T$ gives a $(k-2)$-quartic contact
cocycle on the $\Lambda_k$ descent tower, with reciprocal-norm
coefficient $q_\mathsf M^{(k)}(c) = 1/\langle\Lambda_k|\Lambda_k\rangle$.
The pairing $\langle\Lambda_k|\Lambda_k\rangle$ is rational in $c$
with denominator degree $k-1$ (Pope--Romans--Shen~\cite{PopeRomansShen90}
Prop.~3; each $W^{(k)}$ self-OPE double pole contributes one factor
of $c$), so $q_\mathsf M^{(k)}(c)$ is a proper rational function
of $c$ for each $k$.

For (\ref{eq:hdm-minfty-fibre-total}): the direct-sum decomposition
is the orthogonality of distinct-weight Zamolodchikov quasi-primaries
(Proposition~\ref{prop:higher-w-kappa-matrix} off-diagonal vanishing
$\langle\Lambda_k|\Lambda_\ell\rangle = 0$ for $k\neq\ell$, the
cross-term vanishing of \texttt{w\_algebras\_deep.tex:4142--4145}).
The $\mathrm{Pro}$-object structure is the inverse limit of the
truncated $\mathcal W_N$ fibres as $N\to\infty$, an instance of the
$\mathcal W_{1+\infty}$ inverse-limit presentation of
Definition~(\texttt{w\_algebras.tex:572--578}).

The Hilbert series $\prod_{k\geq 4}(1-q^k)^{-1}$ records one spin-$k$
quasi-primary per $k\geq 4$; the product form is the
Schiffmann--Vasserot CoHA character
on $\mathbb C^2$~\cite{SchiffmannVasserot13} restricted to the
shadow-depth filtration (Cross-reference:
\texttt{w\_algebras.tex:7317--7325}).

For Borel-summability: the Zamolodchikov-norm denominators are
polynomials in $c$ of increasing degree, so $q_\mathsf M^{(k)}(c)$
is $O(c^{-(k-1)})$ as $c\to\infty$. The factorial-divided Borel
transform $\widehat F(\zeta) = \sum_k \zeta^{k-4}\cdot q_\mathsf M^{(k),\mathrm{top}}/(k-4)!$
converges on all of $\mathbb C$ (by Stirling bounds on the
Pope--Romans--Shen structure coefficients), and the original $F(c)$
is recovered by Laplace transform along the positive real axis for
$|c|>1$.
\end{proof}

\begin{remark}[Free-boson $c=1$ and free-fermion $c=-2$
realisations; \ClaimStatusProvedHere]
\label{rem:hdm-minfty-free-field-points}
\index{W1infinity@$\mathcal W_{1+\infty}$!free field points}
Two free-field realisations pin down the first two fibre classes:

\emph{Free boson, $c=1$.} $\mathcal W_{1+\infty}[1]$ is the
free-boson $\mathcal H_1$ with all Bernstein spin-$k$ quasiprimaries.
At $c=1$:
$q_\mathsf M^{(4)}(1) = 10/[1\cdot 27] = 10/27$,
$q_\mathsf M^{(5)}(1) = -48/[1\cdot 27] = -48/27 = -16/9$.
The generating function $F(1) = \sum_{k\geq 4}q_\mathsf M^{(k)}(1)$
converges by explicit Pope--Romans--Shen enumeration, and encodes
the free-boson OPE envelope.

\emph{Free fermion / $bc$-ghost, $c=-2$.} The $bc$-ghost point is
the Koszul-dual of the free boson ($\lambda\to 1-\lambda$, taking
$\lambda=1$ to $\lambda=0$; see~\texttt{w\_algebras.tex:620--623}).
At $c=-2$: $5c+22 = 12$, so $q_\mathsf M^{(4)}(-2) = 10/[(-2)\cdot 12]
= -5/12$. This agrees with the Virasoro shadow at $c=-2$ tabulated
in \texttt{shadow\_tower\_quadrichotomy\_platonic.tex:419} (the
$S_4 = -5/12$ reference point of the principal stratification).
The sign reversal $q_\mathsf M^{(4)}(1) > 0 > q_\mathsf M^{(4)}(-2)$
records the Feigin--Frenkel parity across $c = -22/5$ (the first
Zamolodchikov-norm zero).

\emph{Critical charge $c = -214$.} The Schiffmann--Vasserot CoHA
character degenerates at $c = -214$ (the $\lambda\mapsto 1-\lambda$
image of the irreducibility boundary for
$\mathrm{Vect}(\mathbb C)_\hbar$ at $\hbar = 1/215$). At this
point the Borel transform $\widehat F(\zeta)$ has a pole at
$\zeta = 0$, and the class-$\mathsf M_\infty$ fibre collapses to
a finite-dimensional subquotient of codimension-$1$: the unique
non-trivial Miura-type principal lattice vertex operator algebra
completion.
\end{remark}

\begin{remark}[Class-$\mathsf M_\infty$ vs class-$\mathsf M$ fibre;
\ClaimStatusConditional]
\label{rem:hdm-m-vs-minfty}
\index{Koszul loss classes!$\mathsf M$ vs $\mathsf M_\infty$}
The inclusion $\mathrm{Vir}_c \hookrightarrow \mathcal W_{1+\infty}[c]$
induces a fibre inclusion
$\mathrm{fib}(P_\cC|_{\mathsf M}) \hookrightarrow
\mathrm{fib}(P_\cC|_{\mathsf M_\infty})$ as the $k=4$ summand
(Pope--Romans--Shen level-$1$). The quotient
$\mathrm{fib}|_{\mathsf M_\infty}/\mathrm{fib}|_\mathsf M$
$= \bigoplus_{k\geq 5}q_\mathsf M^{(k)}(c)\cdot\langle W^{(k)}\rangle$
is the new obstruction content that opens only at shadow depth
$r_{\max}=\infty$: it is invisible to the single-Virasoro contact
class $q_\mathsf M$, and records the commuting tower of higher
Zamolodchikov quasi-primaries that Virasoro truncation kills. In
particular, the pair $([\eta_4], q_\mathsf M)$ of
Theorem~\ref{thm:hdm-koszul-loss-classes} is complete for class
$\mathsf M$ at depth~$4$ but incomplete for class $\mathsf M_\infty$
at depth~$\infty$; the complete fibre class on class~$\mathsf M_\infty$
is $([\eta_4], (q_\mathsf M^{(k)})_{k\geq 4})$, with the $k\geq 5$
entries forming the Pope--Romans--Shen tail.
\end{remark}

\begin{theorem}[Tensor-product K\"unneth for the $\mathsf M_\infty$ tower;
\ClaimStatusConditional]
\label{thm:hdm-m-infty-tensor-kunneth}
\index{Koszul loss classes!K\"unneth for $\mathsf M_\infty$|textbf}
\index{chiral tensor product!$\mathsf M_\infty$ tower K\"unneth|textbf}
\index{Schiffmann--Vasserot CoHA!disjoint-union K\"unneth|textbf}
Let $\cA,\cB$ be two modular-Koszul class-$\mathsf M_\infty$ chiral
algebras on $X$, i.e.\ both are presentations of
$\mathcal W_{1+\infty}$ at central charges $c_\cA,c_\cB$. The
Ben-Zvi--Nadler factorisation K\"unneth
quasi-isomorphism~\eqref{eq:hdm-koszul-loss-kunneth} specialises on
the bulk/line fibre to the bigraded splitting
\begin{equation}\label{eq:hdm-m-infty-tensor-fibre}
\mathrm{fib}\bigl(P_\cC\bigm|_{\cA\otimes\cB}\bigr)
\;=\;
\underbrace{\bigoplus_{k\geq 4}
\bigl(q_\mathsf M^{(k)}(c_\cA)+q_\mathsf M^{(k)}(c_\cB)\bigr)\cdot
\langle W^{(k)}\rangle_{\Delta}}_{\text{factor-diagonal sectors}}
\;\oplus\;
\underbrace{\bigoplus_{k,\ell\geq 4}
\frac{k\ell}{c_\cA c_\cB}\cdot\langle{:}W_\cA^{(k)}W_\cB^{(\ell)}{:}\rangle
}_{\text{cross-composite sectors}},
\end{equation}
as a $\mathrm{Pro}$-object in $\mathrm{Vect}_{\geq 0}^{\mathbb N^{2}}$.
The factor-diagonal sectors are the Pope--Romans--Shen tower of
Theorem~\ref{thm:hdm-class-m-infty-fibre} applied independently to
each factor; the cross-composite sectors are the bigraded
normal-ordered products of factor-distinct spins, with reciprocal
Zamolodchikov norms
\[
q_\mathsf M^{(k,\ell)}(c_\cA,c_\cB)
\;=\;
\bigl\langle{:}W^{(k)}_\cA W^{(\ell)}_\cB{:}\,\bigm|\,
{:}W^{(k)}_\cA W^{(\ell)}_\cB{:}\bigr\rangle^{-1}
\;=\;
\frac{k\ell}{c_\cA c_\cB}.
\]
The additive generating function on the factor-diagonal sectors is
\begin{equation}\label{eq:hdm-m-infty-tensor-F}
F_\otimes(c_\cA,c_\cB)
\;:=\;\sum_{k\geq 4}
\bigl(q_\mathsf M^{(k)}(c_\cA)+q_\mathsf M^{(k)}(c_\cB)\bigr)
\;=\;F(c_\cA)+F(c_\cB),
\end{equation}
with Borel transform $\widehat F_\otimes(\zeta) = \widehat F_\cA(\zeta)
+ \widehat F_\cB(\zeta)$. On the $k$-th diagonal sector the
\emph{reciprocal}-norm Kirchhoff law
\begin{equation}\label{eq:hdm-m-infty-tensor-kirchhoff}
\frac{1}{q_\mathsf M^{(k)}(c_\cA\otimes c_\cB)_{\Delta}}
\;=\;\frac{1}{q_\mathsf M^{(k)}(c_\cA)}
\;+\;\frac{1}{q_\mathsf M^{(k)}(c_\cB)}
\end{equation}
expresses the harmonic combining of factor losses, inverse to the
additive $c$-law
$\langle\Lambda_k|\Lambda_k\rangle_{\cA\otimes\cB}
=\langle\Lambda_k|\Lambda_k\rangle_\cA
+\langle\Lambda_k|\Lambda_k\rangle_\cB$. The total Hilbert series
is the Schiffmann--Vasserot CoHA character on $\mathbb C^2\sqcup\mathbb
C^2$:
\begin{equation}\label{eq:hdm-m-infty-tensor-hilbert}
\chi_{\mathrm{fib}}(q_1,q_2)
\;=\;\prod_{k\geq 4}\frac{1}{(1-q_1^k)(1-q_2^k)},
\end{equation}
which is the bigraded MacMahon refinement
of~\eqref{eq:hdm-minfty-fibre-total}.
\end{theorem}

\begin{proof}
The tensor K\"unneth
quasi-isomorphism~\eqref{eq:hdm-koszul-loss-kunneth} of
Lemma~\ref{lem:hdm-koszul-loss-kunneth} applies to any pair of
augmented chiral algebras in the modular landscape; specialising to
$\cA,\cB$ of class~$\mathsf M_\infty$ preserves the Pope--Romans--Shen
spin-filtration on each factor and makes the bigraded structure
manifest.

\emph{Factor-diagonal sectors.} Each factor contributes one spin-$k$
generator $W^{(k)}_\cA$ (resp.\ $W^{(k)}_\cB$) for every $k\geq 1$.
These are orthogonal across factors: $W^{(k)}_\cA\cdot W^{(\ell)}_\cB
\equiv 0$ under the factor-paired inner product, because the
factorisation pairing on
$\cA\otimes\cB$ is the product of the factor pairings and the
augmentations $\epsilon_\cA,\epsilon_\cB$ kill cross-indices. The
spin-$k$ Zamolodchikov quasi-primary on the tensor decomposes as
$\Lambda_k^{\cA\otimes\cB} = \Lambda_k^{\cA}\otimes 1 + 1\otimes
\Lambda_k^{\cB}$, and its self-norm is additive,
$\langle\Lambda_k^{\cA\otimes\cB}|\Lambda_k^{\cA\otimes\cB}\rangle
= \langle\Lambda_k^\cA|\Lambda_k^\cA\rangle
+ \langle\Lambda_k^\cB|\Lambda_k^\cB\rangle$, because the cross-term
$\langle\Lambda_k^\cA\otimes 1 | 1\otimes\Lambda_k^\cB\rangle$
factorises as $\langle\Lambda_k^\cA|1\rangle_\cA\cdot\langle
1|\Lambda_k^\cB\rangle_\cB = 0\cdot 0$ by quasi-primary centring
against the vacuum. Chain-level lift
(Theorem~\ref{thm:hdm-class-m-infty-fibre} proof) gives a
$(k-2)$-contact cocycle per factor with coefficient
$q_\mathsf M^{(k)}(c_\bullet)$; the sum of the two cocycles is the
factor-diagonal loss $q_\mathsf M^{(k)}(c_\cA)+q_\mathsf M^{(k)}(c_\cB)$.
The reciprocal Kirchhoff
law~\eqref{eq:hdm-m-infty-tensor-kirchhoff} is the ratio identity
$1/[x+y] = xy/[xy(x+y)]$ applied to the two factor norms. Taking
the coefficient on the diagonal generator $\langle W^{(k)}\rangle_\Delta
:= \ker(\pi_\cA-\pi_\cB)$ collapses the pure-factor contributions
into a single sum, yielding the first direct summand
of~\eqref{eq:hdm-m-infty-tensor-fibre}.

\emph{Cross-composite sectors.} The normal-ordered composite
${:}W^{(k)}_\cA W^{(\ell)}_\cB{:}$ is a new quasi-primary of spin
$k+\ell$ on $\cA\otimes\cB$, independent of the factor-diagonal
$\Lambda_{k+\ell}^{\cA\otimes\cB}$: the Arnold relation
$\eta_{ij}\wedge\eta_{jk}+\eta_{jk}\wedge\eta_{ki}+\eta_{ki}\wedge
\eta_{ij}=0$ closes factor-locally, so factor-mixed terms form an
independent summand. The self-pairing factorises,
\[
\bigl\langle{:}W^{(k)}_\cA W^{(\ell)}_\cB{:}\,\bigm|\,
{:}W^{(k)}_\cA W^{(\ell)}_\cB{:}\bigr\rangle
\;=\;\langle W^{(k)}_\cA | W^{(k)}_\cA\rangle_\cA\cdot
\langle W^{(\ell)}_\cB | W^{(\ell)}_\cB\rangle_\cB
\;=\;(c_\cA/k)(c_\cB/\ell),
\]
by the leading-pole normalisation
$W^{(j)}_{(2j-1)}W^{(j)} = c/j$ of
Proposition~\ref{prop:higher-w-kappa-matrix}
(\texttt{w\_algebras\_deep.tex:4099}) and the factor-orthogonality
$\langle W^{(k)}_\cA | W^{(\ell)}_\cB\rangle = 0$ for operators on
different factors. The reciprocal is the cross-sector loss
$q_\mathsf M^{(k,\ell)} = k\ell/(c_\cA c_\cB)$, multiplicative in
both reciprocal factor norms.

\emph{Ben-Zvi--Nadler/CoHA K\"unneth.} The Hilbert
series~\eqref{eq:hdm-m-infty-tensor-hilbert} factorises as the
product of two Pope--Romans--Shen Hilbert series of
Theorem~\ref{thm:hdm-class-m-infty-fibre}, one per tensor factor;
this is the bigraded Schiffmann--Vasserot CoHA
character~\cite{SchiffmannVasserot13} on the disjoint union
$\mathbb C^2\sqcup\mathbb C^2$, where each $\mathbb C^2$ carries its
own Hilbert-scheme instanton moduli and the disjoint union realises
the tensor of two independent CoHA algebras. The bigrading
$(q_1,q_2)$ records the spin-index on each copy; the diagonal
$q_1=q_2$ recovers~\eqref{eq:hdm-minfty-fibre-total} on the
factor-symmetrised locus $\cA=\cB$, consistently with
Lemma~\ref{lem:hdm-koszul-loss-kunneth} at $\pi_2$.

\emph{Borel-summability.} The generating
function~\eqref{eq:hdm-m-infty-tensor-F} is the sum of two
Borel-summable Pope--Romans--Shen series of
Theorem~\ref{thm:hdm-class-m-infty-fibre}; sums of Borel-summable
series are Borel-summable with transform the sum of transforms.
\end{proof}

\begin{example}[Two $\mathcal W_{1+\infty}$ at equal charge; free-boson
pair; \ClaimStatusProvedHere]
\label{ex:hdm-m-infty-tensor-examples}
\index{Koszul loss classes!$\mathsf M_\infty\otimes\mathsf M_\infty$ examples}
At $c_\cA = c_\cB = 1$ (two decoupled free bosons) the factor
losses are equal,
$q_\mathsf M^{(4)}(1) = 10/27$ and $q_\mathsf M^{(5)}(1) = -16/9$
(Remark~\ref{rem:hdm-minfty-free-field-points}). The
factor-diagonal sector~\eqref{eq:hdm-m-infty-tensor-fibre} evaluates
at spin $k=4$ to $2q_\mathsf M^{(4)}(1) = 20/27$; the reciprocal
Kirchhoff law~\eqref{eq:hdm-m-infty-tensor-kirchhoff} gives
$1/q_\mathsf M^{(4)}(\Delta) = 27/10+27/10 = 27/5$, hence
$q_\mathsf M^{(4)}(\Delta) = 5/27$ on the diagonal $\langle
W^{(4)}\rangle_\Delta$. The cross-composite
$q_\mathsf M^{(4,4)} = 16/(1\cdot 1) = 16$ dominates both
factor-diagonal contributions. At charges $c_\cA = c_\cB = 1$ the
cross-Hilbert coefficient at bi-spin $(4,4)$ is
$1/[(1-q_1^4)(1-q_2^4)] = 1 + q_1^4 + q_2^4 + q_1^4 q_2^4 + \ldots$,
so the cross-sector gives the new phenomenon absent from either
$\mathcal W_{1+\infty}[1]$ factor alone: factor-mixed
normal-ordered products of spin-$4$ generators, with reciprocal
loss $16$ and no chain-level cancellation. At the Feigin--Frenkel
partner $c_\cA = 1, c_\cB = -2$ the factor-diagonal $k=4$ sum is
$10/27 - 5/12$ and the cross-composite $q_\mathsf M^{(4,4)}(1,-2)
= 16/(-2) = -8$ inherits the sign reversal of the second factor,
recording the $\lambda\to 1-\lambda$ duality
(\texttt{w\_algebras.tex:620--623}) in the bigraded fibre.
\end{example}

\begin{remark}[Mixed classes: $\mathrm{Vir}_{c_1}\otimes
\mathcal W_{1+\infty}[c_2]$]
\label{rem:hdm-m-tensor-minfty-mixed}
\index{Koszul loss classes!mixed $\mathsf M\otimes\mathsf M_\infty$}
For a Virasoro factor $\cA = \mathrm{Vir}_{c_1}$ (class~$\mathsf M$,
only spin-$2$ generator, trivial Pope--Romans--Shen tower beyond
$k=4$) tensored with $\cB = \mathcal W_{1+\infty}[c_2]$
(class~$\mathsf M_\infty$, full tower), the
fibre~\eqref{eq:hdm-m-infty-tensor-fibre} truncates on the
$\cA$-factor to the single spin-$2$ generator. The factor-diagonal
sector survives only at $k=4$ (the Virasoro contact class),
contributing $q_\mathsf M^{(4)}(c_1) + q_\mathsf M^{(4)}(c_2)$ on
the diagonal $\langle\Lambda\rangle_\Delta$. Higher-spin
factor-diagonal sectors reduce to the single-factor $\mathcal
W_{1+\infty}[c_2]$ entry $q_\mathsf M^{(k)}(c_2)$ for $k\geq 5$.
The cross-composite sector retains only the
Virasoro--times--$W^{(\ell)}$ contribution ${:}T\cdot
W^{(\ell)}_\cB{:}$ at bi-spin $(2,\ell)$ with reciprocal loss
$q_\mathsf M^{(2,\ell)} = 2\ell/(c_1 c_2)$. The Hilbert series is
$(1-q_1^4)^{-1}\prod_{k\geq 4}(1-q_2^k)^{-1}$, the mixed
SV~CoHA~$\times$~single-Virasoro character; the factor-symmetry
$q_1\leftrightarrow q_2$ of
Theorem~\ref{thm:hdm-m-infty-tensor-kunneth} is broken, reflecting
the asymmetry of the shadow-depth dichotomy $r_{\max}=4$ vs
$r_{\max}=\infty$ between the two classes.
\end{remark}

\begin{theorem}[$n$-fold tensor K\"unneth and Kirchhoff law for the
$\mathsf M_\infty$ tower; \ClaimStatusConditional]
\label{thm:hdm-m-infty-nfold-kunneth}
\index{Koszul loss classes!$n$-fold K\"unneth for $\mathsf M_\infty$|textbf}
\index{chiral tensor product!$n$-fold Kirchhoff law|textbf}
\index{Schiffmann--Vasserot CoHA!$n$-fold disjoint-union K\"unneth|textbf}
Let $\cA_1,\ldots,\cA_n$ be modular-Koszul class-$\mathsf M_\infty$
chiral algebras on~$X$, each a presentation of $\mathcal W_{1+\infty}$
at central charge~$c_{\cA_i}$. The iterated Ben-Zvi--Nadler
factorisation K\"unneth on $\bigotimes_{i=1}^n\cA_i$ splits the
bulk/line fibre into $2^n-1$ power-set strata indexed by non-empty
subsets $S\subseteq\{1,\ldots,n\}$:
\begin{equation}\label{eq:hdm-m-infty-nfold-fibre}
\mathrm{fib}\bigl(P_\cC\bigm|_{\bigotimes_i\cA_i}\bigr)
\;=\;
\bigoplus_{\varnothing\neq S\subseteq\{1,\ldots,n\}}\;
\bigoplus_{\mathbf{k}_S\in\mathbb Z_{\geq 4}^S}
q_\mathsf M^{(\mathbf{k}_S)}\bigl((c_{\cA_i})_{i\in S}\bigr)\cdot
\bigl\langle {:}\!\prod_{i\in S}W^{(k_i)}_{\cA_i}\!{:}\bigr\rangle,
\end{equation}
as a $\mathrm{Pro}$-object in $\mathrm{Vect}_{\geq 0}^{\mathbb N^{n}}$.
The rank-$|S|=1$ strata are the factor-diagonal Pope--Romans--Shen
sectors; the rank-$|S|=r\geq 2$ strata are cross-composite
normal-ordered products across $r$ distinct factors, with reciprocal
Zamolodchikov norm
\begin{equation}\label{eq:hdm-m-infty-nfold-cross}
q_\mathsf M^{(\mathbf{k}_S)}\bigl((c_{\cA_i})_{i\in S}\bigr)
\;=\;\prod_{i\in S}\frac{k_i}{c_{\cA_i}}.
\end{equation}
At rank~$r$ the subset count is $\binom{n}{r}$, and the total
non-trivial stratum count is $\sum_{r=1}^n\binom{n}{r}=2^n-1$.
On the factor-diagonal
$\langle W^{(k)}\rangle_{\Delta} := \ker(\pi_{\cA_1}-\pi_{\cA_2},
\pi_{\cA_1}-\pi_{\cA_3},\ldots,\pi_{\cA_1}-\pi_{\cA_n})$ the reciprocal
norms obey the \emph{$n$-fold Kirchhoff parallel law}
\begin{equation}\label{eq:hdm-m-infty-nfold-kirchhoff}
\frac{1}{q_\mathsf M^{(k)}(c_{\cA_1}\otimes\cdots\otimes
c_{\cA_n})_{\Delta}}
\;=\;\sum_{i=1}^n\frac{1}{q_\mathsf M^{(k)}(c_{\cA_i})},
\end{equation}
dual to the arithmetic additivity
$\langle\Lambda_k|\Lambda_k\rangle_{\bigotimes_i\cA_i}
=\sum_i\langle\Lambda_k|\Lambda_k\rangle_{\cA_i}$ of the diagonal
Zamolodchikov quasi-primary
$\Lambda_k^{\bigotimes}=\sum_i 1^{\otimes(i-1)}\otimes\Lambda_k^{\cA_i}
\otimes 1^{\otimes(n-i)}$. The total bigraded Hilbert series on
$n$~formal variables is the $n$-fold Schiffmann--Vasserot CoHA
character on $\bigsqcup_{i=1}^n\mathbb C^2_i$:
\begin{equation}\label{eq:hdm-m-infty-nfold-hilbert}
\chi_{\mathrm{fib}}(q_1,\ldots,q_n)
\;=\;\prod_{i=1}^n\prod_{k\geq 4}\frac{1}{1-q_i^k}.
\end{equation}
When all factors coincide, $\cA_i=\cA$ with central charge~$c$, the
$\mathrm{Sym}_n$-averaged fibre is Schur-expandable in power-sum
symmetric functions $p_\lambda(q_1,\ldots,q_n)$, indexed by
partitions $\lambda\vdash r$ with $\ell(\lambda)\leq n$, and the
symmetrised Kirchhoff diagonal collapses to
$q_\mathsf M^{(k)}(\cA^{\otimes n})_{\Delta}=q_\mathsf M^{(k)}(c)/n$.
\end{theorem}

\begin{proof}
Iterate the two-factor K\"unneth
quasi-isomorphism~\eqref{eq:hdm-koszul-loss-kunneth} of
Lemma~\ref{lem:hdm-koszul-loss-kunneth}: Ben-Zvi--Nadler
factorisation is associative on the product side and symmetric
monoidal on the fibre side, so the $n$-fold fibre equals the
$n$-fold tensor product of factor fibres. Expanding each factor
fibre as $\mathrm{fib}(P_\cC|_{\cA_i}) = \bigoplus_{k\geq 4}
q_\mathsf M^{(k)}(c_{\cA_i})\cdot\langle W^{(k)}_{\cA_i}\rangle$
(Theorem~\ref{thm:hdm-class-m-infty-fibre}) and distributing over
the tensor product produces the power-set
decomposition~\eqref{eq:hdm-m-infty-nfold-fibre}: each summand of
the $n$-fold tensor selects either the vacuum on factor~$i$ or a
spin-$k_i$ generator, and the non-vacuum factors form the subset
$S\subseteq\{1,\ldots,n\}$.

\emph{Self-pairing on the cross-composite.} The factor-wise
orthogonality $\langle W^{(k)}_{\cA_i}|W^{(k')}_{\cA_j}\rangle = 0$
for $i\neq j$ (the augmentations $\epsilon_{\cA_j}$ annihilate the
cross brackets, and Arnold
$\eta_{ij}\wedge\eta_{jk}+\eta_{jk}\wedge\eta_{ki}
+\eta_{ki}\wedge\eta_{ij}=0$ closes factor-locally) iterates across
all distinct factor pairs in~$S$; the self-pairing of
${:}\prod_{i\in S}W^{(k_i)}_{\cA_i}{:}$ factorises as $\prod_{i\in
S}\langle W^{(k_i)}_{\cA_i}|W^{(k_i)}_{\cA_i}\rangle
=\prod_{i\in S}c_{\cA_i}/k_i$, yielding~\eqref{eq:hdm-m-infty-nfold-cross}.

\emph{Kirchhoff parallel law.} The $n$-fold diagonal
quasi-primary $\Lambda_k^{\bigotimes}$ has cross-factor pairings
$\langle 1^{\otimes(i-1)}\otimes\Lambda_k^{\cA_i}\otimes\cdots|
1^{\otimes(j-1)}\otimes\Lambda_k^{\cA_j}\otimes\cdots\rangle
= \langle\Lambda_k^{\cA_i}|1\rangle_{\cA_i}\cdot
\langle 1|\Lambda_k^{\cA_j}\rangle_{\cA_j}\cdot\prod_{\ell\neq i,j}
\langle 1|1\rangle_{\cA_\ell} = 0$ for $i\neq j$ by quasi-primary
vacuum-centring; the self-pairing collapses to
$\sum_i\langle\Lambda_k^{\cA_i}|\Lambda_k^{\cA_i}\rangle_{\cA_i}$.
Inverting gives $1/q_\mathsf M^{(k)}(\bigotimes_i\cA_i)_\Delta
=\sum_i 1/q_\mathsf M^{(k)}(c_{\cA_i})$, the parallel-resistance
law~\eqref{eq:hdm-m-infty-nfold-kirchhoff}. The $n=2$ case recovers
equation~\eqref{eq:hdm-m-infty-tensor-kirchhoff}.

\emph{Hilbert series.} Each factor Pope--Romans--Shen
series~\eqref{eq:hdm-minfty-fibre-total} is $\prod_{k\geq
4}(1-q_i^k)^{-1}$; independence of factors makes the $n$-fold fibre
character the product, yielding~\eqref{eq:hdm-m-infty-nfold-hilbert}.
This is the Schiffmann--Vasserot CoHA character on
$\bigsqcup_{i=1}^n\mathbb C^2_i$, $n$ disjoint Hilbert-scheme
instanton moduli with their own $(q_i)$ gradings.

\emph{$\mathrm{Sym}_n$-symmetrisation.} At coinciding factors
$\cA_i=\cA$, the character~\eqref{eq:hdm-m-infty-nfold-hilbert}
becomes symmetric in $(q_1,\ldots,q_n)$; the Macdonald--Cauchy
identity expands it in power-sum basis $p_\lambda$, and the
$\mathrm{Sym}_n$-average projects to the $p_\lambda$-diagonal with
$\ell(\lambda)\leq n$. The symmetrised factor-diagonal at spin~$k$
is $\sum_{i=1}^n\langle\Lambda_k|\Lambda_k\rangle_\cA = n\langle
\Lambda_k|\Lambda_k\rangle_\cA$, so the reciprocal is
$q_\mathsf M^{(k)}/n$.
\end{proof}

\begin{example}[Three and more $\mathcal W_{1+\infty}$ factors;
\ClaimStatusProvedHere]
\label{ex:hdm-m-infty-nfold-examples}
\index{Koszul loss classes!$\mathsf M_\infty^{\otimes n}$ examples}
Three $\mathcal W_{1+\infty}$ factors at charges $c_1,c_2,c_3$
decompose the fibre into $2^3-1=7$ strata: three rank-$1$ factor
diagonals (subsets $\{1\},\{2\},\{3\}$), three rank-$2$ cross
composites (subsets $\{1,2\},\{1,3\},\{2,3\}$ with reciprocal norms
$k_1k_2/(c_1c_2)$ etc.), and one rank-$3$ triple composite at subset
$\{1,2,3\}$ with reciprocal norm $k_1k_2k_3/(c_1c_2c_3)$. The
Kirchhoff diagonal reads
$1/q_\Delta^{(k)}=1/q_1^{(k)}+1/q_2^{(k)}+1/q_3^{(k)}$; at equal
charges $c_1=c_2=c_3=1$ (three free bosons) the $k=4$ entries are
$q_1^{(4)}=q_2^{(4)}=q_3^{(4)}=10/27$, so $q_\Delta^{(4)}=10/81$,
and the triple cross-composite $q_\mathsf M^{(4,4,4)}(1,1,1)=64$
sets the dominant loss scale. For general~$n$ and identical charges
$c_i=1$, the $\mathrm{Sym}_n$-symmetrised $k=4$ diagonal loss is
$q_\Delta^{(4)}=10/(27n)$, vanishing linearly in~$n$: $n$ parallel
channels attenuate the single-spin contact obstruction harmonically,
the CoHA-theoretic manifestation of the Kirchhoff $n$-resistor law.
At the Feigin--Frenkel parity triple $c_1=1,c_2=-2,c_3=1$ the
rank-$3$ cross-composite $q_\mathsf M^{(4,4,4)}(1,-2,1)=64/(-2)=-32$
carries the sign of the central factor, recording
$\lambda\mapsto 1-\lambda$ duality
(\texttt{w\_algebras.tex:620--623}) propagated through odd-order
cross-composites.
\end{example}

\begin{example}[$n$-fold Heisenberg tower; \ClaimStatusProvedHere]
\label{ex:hdm-m-infty-nfold-heisenberg}
\index{Heisenberg!$n$-fold tensor tower}
For $\cH_{k_1}\otimes\cdots\otimes\cH_{k_n}$ ($n$ Heisenberg factors
at levels $k_i$), each factor is $\mathsf G$-class (shadow depth~$2$)
rather than $\mathsf M_\infty$,
so~\eqref{eq:hdm-m-infty-nfold-fibre} collapses degree-wise: the
only surviving fibre stratum is the factor-diagonal spin-$1$
current $\partial\phi_i$. The Kirchhoff law becomes trivially
additive in the $\kappa$-invariant: $\kappa(\bigotimes_i\cH_{k_i})
=\sum_i k_i$, consistent with
Proposition~\ref{prop:higher-w-kappa-matrix} and, on the
$\mathsf G$-family stratum specifically (Theorem~C ceiling
restricted to the Heisenberg class, where
$\kappa+\kappa^!=0$; distinct from the $\mathsf L$ ceiling~$13$,
$\mathsf C$ ceiling~$250/3$, $\mathsf M$ ceiling~$98/3$). All cross-composite
sectors are absorbed into the Heisenberg diagonal by the bosonic
OPE, confirming the shadow-depth--$2$ truncation: no new content
opens at higher rank~$|S|$.
\end{example}

\begin{example}[$n$-fold Virasoro; explicit cross-composite
combinatorics; \ClaimStatusProvedHere]
\label{ex:hdm-m-infty-nfold-virasoro}
\index{Virasoro!$n$-fold tensor cross-composites}
$\mathrm{Vir}_{c_1}\otimes\cdots\otimes\mathrm{Vir}_{c_n}$ (each
factor class~$\mathsf M$ with spin-$2$ generator only): the
fibre~\eqref{eq:hdm-m-infty-nfold-fibre} truncates to rank-$|S|$
strata where only the spin~$k_i=2$ generator survives per chosen
factor. The factor-diagonal is a single spin-$4$ stratum
$\Lambda = {:}\!TT\!{:} - (3/10)\partial^2 T$ on each factor, with
reciprocal norms $q_\mathsf M^{(4)}(c_i)=10/[c_i(5c_i+22)]$; the
Kirchhoff law reads
$1/q_\Delta^{(4)}=\sum_{i=1}^n c_i(5c_i+22)/10$. The rank-$r$
cross-composite is ${:}\!T_{\cA_{i_1}}T_{\cA_{i_2}}\cdots
T_{\cA_{i_r}}\!{:}$ at bi-spin $(2,2,\ldots,2)$ over~$S=\{i_1,\ldots,
i_r\}$, with reciprocal norm $q_\mathsf M^{(2,\ldots,2)}_S
=\prod_{j\in S}2/c_j = 2^r/\prod_{j\in S}c_j$. There are
$\binom{n}{r}$ such cross-composites at rank~$r$, total $2^n-n-1$
genuinely-mixed cross-composites. At $n$ Ising factors
$c_i=1/2$, the $\mathrm{Sym}_n$-symmetrised diagonal is
$1/q_\Delta^{(4)}=n\cdot(1/2)(5/2+22)/10=n\cdot 49/40$, so
$q_\Delta^{(4)}=40/(49n)$; the rank-$r$ cross-composite is
$q_\mathsf M^{(2^r)}=2^r/(1/2)^r=4^r$, the geometric-series growth
signalling that high-rank mixed contacts dominate the fibre loss.
This is the Virasoro incarnation of the $n$-fold Kirchhoff
principle: arithmetic additivity in norms, harmonic parallelism in
reciprocal norms, multiplicative iteration across cross-factor
composites.
\end{example}

\begin{remark}[Boolean-lattice structure of the power-set
decomposition]
\label{rem:hdm-m-infty-nfold-boolean}
\index{Koszul loss classes!Boolean power-set stratification}
The subset index $S\subseteq\{1,\ldots,n\}$
of~\eqref{eq:hdm-m-infty-nfold-fibre} runs over the Boolean lattice
$2^{\{1,\ldots,n\}}$, graded by $|S|$; the inclusion of rank-$r$
into rank-$(r+1)$ strata is realised by the chiral normal-ordered
insertion ${:}W^{(k)}\cdot-{:}$ for each new factor~$i\notin S$. The
Hilbert series~\eqref{eq:hdm-m-infty-nfold-hilbert} factorises as
$\prod_i\chi_i$ because the Boolean lattice is the product of $n$
two-element chains $\{\varnothing,\{i\}\}$, each contributing one
factor $\chi_i=\prod_{k\geq 4}(1-q_i^k)^{-1}$. The
$\mathrm{Sym}_n$-action by factor permutations descends the Boolean
lattice to the partition lattice $\mathrm{Part}_n$, and the
symmetrised fibre character is indexed by partitions~$\lambda$ with
$\ell(\lambda)\leq n$: Schur--Weyl combinatorics transported to the
$\mathsf M_\infty$ tower.
\end{remark}

\section{The three-parameter $\hbar$ identification}
\label{sec:hdm-hbar-identification}

The Yangian quantization of the Sklyanin bracket introduces a
formal parameter~$\hbar$ that controls the deformation from
classical to quantum. In the present framework, $\hbar$ admits
three independent identifications, each corresponding to a different
face of the collision residue.

\begin{theorem}[Three-parameter $\hbar$ identification; \ClaimStatusProvedElsewhere]
\label{thm:hdm-hbar-three-identification}
\index{$\hbar$!three identifications|textbf}
For $\cA = \widehat{\fg}_k$, the formal parameter $\hbar$ that
deforms the Sklyanin Poisson bracket into the Yangian
$Y_\hbar(\fg)$ admits the following three identifications:
\begin{enumerate}[label=\textup{(\roman*)}]
\item \emph{Inverse level:} $\hbar = 1/(k + h^\vee)$, where the
 $k \to \infty$ limit is the classical limit and the
 $k \to -h^\vee$ limit is the critical-level Gaudin limit.
\item \emph{Inverse spectral parameter:} $\hbar = 1/z$ in the formal
 expansion $r^{\mathrm{Dr}}(z) = k\Omega_{\mathrm{tr}}/z
 = k\Omega_{\mathrm{tr}} \cdot \hbar$
 (level-$k$ form, vanishing at $k=0$),
 with the $z \to \infty$ limit identifying classical with quantum
 at large separation.
\item \emph{Bar-degree parameter:} $\hbar$ counts the bar degree of
 the iterated collision residue
 $\mathrm{Res}^{\mathrm{coll}}_{0,n}(\Theta_\cA)$, with the
 $\hbar \to 0$ limit projecting to the genus-zero, binary stratum.
\end{enumerate}
The three identifications are mutually consistent in the sense that
the corresponding three formal expansions of $r_\cA(z)$ agree term
by term. The identification $(\mathrm{i})$ is Drinfeld's classical
theorem~\cite{Drinfeld85}; identification $(\mathrm{ii})$ is the
spectral-parameter expansion of the same theorem; identification
$(\mathrm{iii})$ is the bar-complex translation supplied by the
present framework.
\end{theorem}

\begin{proof}
Identifications $(\mathrm{i})$ and $(\mathrm{ii})$ are classical
results of Drinfeld~\cite{Drinfeld85}: the Yangian
$Y_\hbar(\fg) = \mathcal{U}(\fg[\![z^{-1}]\!]) \otimes \mathbb{C}[\![\hbar]\!]$
with $\hbar = 1/(k+h^\vee)$ specializes to the universal enveloping
algebra at $\hbar = 0$, and the $r$-matrix expands as
$r^{\mathrm{Dr}}(z) = k\Omega_{\mathrm{tr}}/z
= k\Omega_{\mathrm{tr}} \cdot \hbar$ in the
spectral parameter. Identification $(\mathrm{iii})$ is the
bar-complex translation: the iterated collision residue
$\mathrm{Res}^{\mathrm{coll}}_{0,n}(\Theta_\cA)$ has bar degree
$n - 2$, and the $\hbar$-expansion of $r_\cA(z)$ corresponds to the
filtration by collision depth. The mutual consistency of the three
identifications is the content of the bar/cobar/spectral
translation triangle.
\end{proof}

The three-parameter $\hbar$ identification is the structural reason
why the collision residue admits the seven independent realizations on
the comparison surfaces above: the parameter $\hbar$ admits three
independent meanings (level, spectral parameter, bar degree), and
each meaning produces a distinct mathematical object that coincides
with $r_\cA(z)$ on the common domain.

\section{Higher Gaudin Hamiltonians}
\label{sec:hdm-higher-gaudin}

The standard Gaudin model uses the binary collision residue
$r_\cA(z)$ to construct commuting Hamiltonians. The bar-intrinsic
MC element $\Theta_\cA$ contains higher collision residues at
ternary, quaternary, and arbitrary depths, which produce \emph{higher
Gaudin Hamiltonians}: commuting differential operators of higher
order on the sphere, generated by iterated collision residues.

\begin{theorem}[Higher Gaudin Hamiltonians; \ClaimStatusConditional]
\label{thm:hdm-higher-gaudin}
\index{Gaudin Hamiltonians!higher|textbf}
\index{collision residue!higher orders}
Let $\cA$ be a modular Koszul chiral algebra with collision depth
$k_{\max}(\cA) \geq 1$. For each integer $k$ with $1 \leq k \leq
k_{\max}(\cA)$, the modular MC element $\Theta_\cA$ produces a
family of commuting differential operators
\begin{equation}\label{eq:hdm-higher-gaudin}
H_i^{(k)}
\;=\;
\sum_{j \neq i}
\bigl[\mathrm{Res}^{\mathrm{coll},\,k}_{0,2}(\Theta_\cA)\bigr]_{i,j}
\;\in\;
\mathrm{Diff}^{k-1}(\mathrm{Conf}_n(\mathbb{P}^1)),
\end{equation}
where $\mathrm{Res}^{\mathrm{coll},k}_{0,2}$ extracts the order-$k$
pole of the collision expansion. The operators $H_i^{(k)}$
mutually commute: $[H_i^{(k)}, H_j^{(k')}] = 0$ for all
$i, j, k, k'$. For $k = 1$, $H_i^{(1)} = H_i^{\mathrm{Gaudin}}$
recovers the standard Gaudin Hamiltonians of Face~7. For
$k \geq 2$, $H_i^{(k)}$ are differential operators of order $k - 1$
that act non-trivially on conformal blocks of $\cA$ via the modular
shadow connection.
\end{theorem}

\begin{proof}
The higher collision residues exist because the bar-intrinsic
MC element $\Theta_\cA$ has well-defined restrictions to all
boundary strata of all moduli spaces, including the higher-order
contact strata along the diagonal of $\overline{\cM}_{0,2}$. The
order-$k$ pole of the collision expansion produces a $k$-th
derivative term in the corresponding Hamiltonian. Mutual
commutativity is the projection of the modular MC equation
$D_\cA \Theta_\cA + \frac{1}{2}[\Theta_\cA, \Theta_\cA] = 0$ to
collision depth $k + k'$: each cross-bracket
$[H_i^{(k)}, H_j^{(k')}]$ is the collision-depth-$(k+k')$ projection
of the corresponding term in the MC equation, which vanishes by the
bar nilpotency $D_\cA^2 = 0$ at the relevant stratum. This is the
generalization of the infinitesimal braid relation
(Theorem~\ref{thm:collision-depth-2-ybe}) from depth~$2$ to
arbitrary depth, and follows from the bar-intrinsic construction
(Theorem~\ref{thm:mc2-bar-intrinsic}) by extracting successive
collision residues and identifying the resulting commuting
Hamiltonians via Face~4.
\end{proof}

\begin{remark}[Higher Gaudin and the operator-order trichotomy]
\label{rem:hdm-higher-gaudin-trichotomy}
For $\cA$ in shadow class~$G$ (Heisenberg) or class~$L$ (affine
Kac--Moody) with $k_{\max} = 1$, only the first-order Gaudin
Hamiltonians exist, and they are scalar multiplication operators on
conformal blocks. For $\cA$ in class~$M$ with $k_{\max} \geq 3$
(Virasoro, $\Walg_N$), the higher Gaudin Hamiltonians of order $\geq 2$
appear, and the operator-order trichotomy of
Theorem~\ref{thm:k-max-trichotomy} gives the full spectrum of
differential orders. For class~$C$ with $k_{\max} = 0$
($\beta\gamma$), there are no Gaudin Hamiltonians at all from the
binary collision residue, and the integrable structure must come
from higher-degree shadow contributions.
\end{remark}

The higher Gaudin Hamiltonians are the natural extension of
Face~7 to algebras beyond affine Kac--Moody. Where the standard
Gaudin model produces commuting first-order operators from the
binary $r$-matrix, the higher Gaudin construction produces
commuting differential operators of arbitrary order from the
iterated collision residues of $\Theta_\cA$. This is one way the
seven-face programme extends classical integrable systems to the
full standard landscape.

\section{Falsification by face comparison}
\label{sec:hdm-falsification}

The seven-face master theorem is more than an organizational
convenience. It is a verification mechanism: any computation of
$r_\cA(z)$ in an available face can be cross-checked against any other
available face. When two computations of the same chiral algebra
disagree on their common comparison surface, the disagreement is
informative: it locates an error in one or both of the computations.

This is the multi-path verification mandate (three-paths
rule) realized at the structural level. Where the standard mandate
requires three independent verification paths for a numerical
formula, the seven-face programme provides seven independent paths
for a single mathematical object, restricted to the faces available
for the algebra under study. When the available faces agree, the object
is verified on that comparison surface. When any two disagree, the
project enters a falsification phase: identify the source of the
disagreement, fix the error in the offending framework, and re-verify
across the available faces.

\begin{principle}[Seven-face falsification]
\label{princ:hdm-falsification}
A claim about $r_\cA(z)$ is verified only when it has been computed
in at least three independent faces of the seven-face programme and
all three agree. When any two faces disagree, the disagreement
locates an error in at least one of the two computations. The
seven-face programme is therefore both a structural unification
\textup{(}seven names for one object\textup{)} and a falsification
mechanism \textup{(}disagreement implies error\textup{)}.
\end{principle}

Two local tests illustrate the principle. In affine Kac--Moody, the
trace-form Drinfeld kernel $k\,\Omega_{\mathrm{tr}}/z$ and the
KZ-normalized collision residue $\Omega/((k+h^\vee)z)$ agree only
after the finite Casimir rescaling
$k\,\Omega_{\mathrm{tr}}\mapsto\Omega/(k+h^\vee)$; treating the two as
a literal rational-function identity gives the wrong behaviour at
$k=0$ and at the critical level. In the DNP comparison for
$\beta\gamma$, the line-operator identification requires the
Swiss-cheese formality input of Vol.~II; bar formality alone does not
identify the pinned bulk slot
$\cZ^{\mathrm{der}}_{\mathrm{ch}}(\cA)$ with the line model. These are
mathematical obstruction tests, not drafting history.

\begin{remark}[Why seven, not more]
\label{rem:hdm-why-seven}
The seven faces of $r_\cA(z)$ enumerated here are not exhaustive.
Additional faces include the Vassiliev knot invariant face (the
collision residue is the universal weight system for finite-type
invariants), the Kontsevich integral face (the collision residue is
the propagator of the configuration-space integral), the BV/BRST
master equation face (the collision residue is the leading-order
solution to the QME), and the geometric Langlands face (the
collision residue at critical level $k = -h^\vee$ is the symbol of
the oper differential operator). Each of these is a sub-face of
one of the seven listed here \textup{(}Vassiliev and Kontsevich
inside Face~5; BV/BRST inside Face~3; Langlands inside Face~7 at
critical level\textup{)}, but their complete elaboration is the
content of dedicated chapters elsewhere in the monograph. The
seven listed here are the primary faces; the falsification principle
applies to any pair drawn from any of the faces, listed or not.
\end{remark}

\begin{remark}[The methodological hierarchy]
\label{rem:hdm-methodological-hierarchy}
The seven-face master theorem instantiates a methodological
hierarchy: \emph{algebraic, geometric, classical, differential,
historical, Poisson, integrable}. Face~1 is algebraic
\textup{(}twisting morphism in a dg category\textup{)}; Face~2 is
geometric \textup{(}line operators in a 3d theory\textup{)}; Face~3
is classical \textup{(}Poisson vertex algebra\textup{)}; Face~4 is
differential \textup{(}commuting operators on a moduli space\textup{)};
Face~5 is historical \textup{(}original Drinfeld $r$-matrix\textup{)};
Face~6 is Poisson \textup{(}Sklyanin bracket\textup{)}; Face~7 is
integrable \textup{(}Gaudin Hamiltonians\textup{)}. Each level of
the hierarchy is independently meaningful and independently
verifiable; the master theorem asserts that their available
specialisations compute the same object on the common domain.
\end{remark}

The seven-face programme is a core result of Part~\ref{part:standard-landscape}
because it converts the abstract claim ``the bar-intrinsic MC
element is universal'' into a concrete checklist: seven frameworks,
available independent computations, one mathematical object. When the
available computations agree, the universality claim is no longer
abstract on that comparison surface. When two of them disagree, the project knows where to
look. This is the operational meaning of universality in this
monograph, and the structural reason why the collision residue is
the central object of the holographic modular Koszul datum: it is the
binary witness entering the restriction map
$\rho_{\partial}\colon\cC\to\cA$ and the operadic inclusion
$\mathrm{Ass}^{\mathrm{ch}}\hookrightarrow\mathrm{SC}^{\mathrm{ch,top}}$
alongside $\Theta_\cA$ and $\nabla^{\mathrm{hol}}$. On the named
modular-Koszul / boundary-linear / principal DS surfaces those three
witnesses close the comparison; off them the strict-edge and
global-triangle obstructions remain active.

\section{Structural cross-references}
\label{sec:hdm-anchors}

\begin{remark}[Anchor map for the seven-face datum]
\label{rem:hdm-anchors}
\index{seven-face datum!structural anchors}
\index{cross-volume bridge!seven-face datum}
The seven-face master theorem
(Theorem~\ref{thm:hdm-seven-way-master}) is the verification
checkpoint for slot~(4) of the holographic modular Koszul
datum~$\mathcal{H}(T)$. Here the boundary object is~$\cA$, the bulk
object is the pinned centre
$\cC=\cZ^{\mathrm{der}}_{\mathrm{ch}}(\cA)$, the relevant maps are the
restriction $\rho_{\partial}\colon \cC\to\cA$ and the operadic
inclusion $\mathrm{Ass}^{\mathrm{ch}}\hookrightarrow
\mathrm{SC}^{\mathrm{ch,top}}$, and the witnesses are the same triple
$(r_\cA(z),\Theta_\cA,\nabla^{\mathrm{hol}})$. Outside the
modular-Koszul / boundary-linear / principal DS comparison surfaces,
the anchor map stops at coderived or comparison-surface statements
because the strict edge $\cA^{!}_\infty\to\cA^!$ can fail and the
quartic contact class blocks boundary-linearity on nonlinear
class~$\mathsf{M}$. Each of the remaining slots projects through a
structural anchor:
\begin{enumerate}[label=\textup{(W\arabic*)}]
\item \emph{Bar-cobar reflection.} The universal twisting morphism
 $\pi_\cA$ of Face~1 is the unit of the Koszul reflection theorem
 (Theorem~\ref{thm:koszul-reflection}, the Vol~I Platonic
 Theorem~A): the bar-cobar adjunction
 $\Omega \dashv B$ is the fixed-curve chain/factorization
 reflection on the Koszul locus, with coderived correction off the
 strict lane, where the surviving bar class $\eta_4$ is the concrete
 obstruction to the strict edge $\cA^{!}_\infty\to\cA^!$, and
 $\pi_\cA$ is its tautological twisting datum. This
 is slot~(1) of $\mathcal{H}(T)$ verified \emph{by}
 reflection.
\item \emph{Universal conductor.} Slot~(5), the universal MC
 element $\Theta_\cA$ at the genus-$0$, scalar projection,
 evaluates on the verified standard-landscape canonical BRST models to
 the modular characteristic
 $\kappa(\cA) = -c_{\mathrm{ghost}}(\BRST^{\mathrm{can}}(\cA))$
 established in Chapter~\ref{chap:universal-conductor}. No stronger
 universal BRST functor is claimed at this surface. The collision
 residue $r_\cA(z)$ and the conductor pair
 $(\kappa(\cA),K(\cA))$ are therefore \emph{linked} invariants: the
 leading $\Theta_\cA$ slot fixes the scalar $\kappa(\cA)$ directly and
 fixes the conductor $K(\cA)=c(\cA)+c(\cA^!)$ after passing to the
 Verdier partner.
\item \emph{Climax identity.} The full datum $\mathcal{H}(T)$
 condenses into the climax identity
 $d_{\mathrm{bar}} = \mathrm{KZ}^*(\nabla_{\mathrm{Arnold}})$
 of Chapter~\ref{ch:climax-theorem}: the modular shadow
 connection $\nabla^{\mathrm{hol}}$ of slot~(6) is the
 KZ-pullback of the Arnold connection on $\overline{C}_n(X)$,
 and the seven-face master theorem records the seven names
 under which the binary collision residue~$r_\cA(z)$ enters
 that pullback. Outside the stated comparison surfaces this remains
 only a partial anchor: there is no proved degeneration-independent
 sewing theorem identifying the higher-genus horizontal sections of
 $\nabla^{\mathrm{hol}}_{g,n}$ with a physical moduli problem, and the
 global triangle comparing the bulk slot to the line model is still
 blocked by the strict edge and quartic-contact obstructions.
\end{enumerate}
The cross-volume bridges are explicit:
Vol~II's universal holographic master theorem
(Vol~II Theorem~\ref{V2-thm:universal-holography-master}, cited
across the volume boundary) replaces the Verdier-dual
factorization companion~$\cA^{!}_\infty$ by an honest
3d HT bulk algebra, lifting
Proposition~\ref{prop:hdm-u-a-recoverability} from a
heuristic to a chain-level theorem on the gauge-theoretic
landscape; Vol~III's CY-to-chiral two-stage factorisation
$\Phi_{3,\sigma} =
\mathrm{Sp}^{\mathrm{ch}}_{\sigma}\circ\Phi^{\mathrm{FA}}_3$,
$\sigma=(\Sigma_2,C)$, has stage~$1$ the canonical
$E_3$-holomorphic factorisation algebra on the CY$_3$ variety
\textup{(}pinned by Kontsevich--Tamarkin\textup{)} and stage~$2$
factorisation homology over~$\Sigma_2$ followed by restriction to the
reference curve~$C$. It realises the seven faces of
$r_{\Phi_{3,\sigma}(D^b\Coh(Y))}(z)$ for K3-fibered Class~A
threefolds~$Y$ on the chosen stage-$2$ shadow. At
$Y=K3\times E$ the pinned specialisation is
$\sigma_{K3,E}=(K3,E)$, so
$\Phi_{3,\sigma_{K3,E}}(D^b\Coh(K3\times E))$ is the boundary
algebra whose BKM closure is controlled by $\Delta_5$. The universal
$\Phi$-trace identity compares the Vol~I $\Phi$-trace conductor on
this branch with the proved Borcherds weight $c_N(0)/2$ on the five
verified Borcherds families $N \in \{1,2,3,4,6\}$. The obstruction is
explicit: the chain-level BRST-to-Borcherds comparison has not been
proved outside those five families, and the equality with the
Vol~I conductor is a universal-trace comparison, not a formal
consequence of the two-stage functor alone. This supplies the
conductor side of slot~(5) in the BPS-counting category.
\end{remark}

\begin{remark}[Universal $\Phi$-trace identity at $K3 \times E$: the BKM weight is $5$]
\label{rem:holo-master-phi-trace}
\index{universal $\Phi$-trace identity!$K3 \times E$ specialisation}
\index{BRST conductor!Borcherds singular-theta pushforward}
\index{Borcherds weight!$\kappa_{\mathrm{BKM}} = c_\Lambda(0)/2$}
Write
\[
 A_{K3\times E,\sigma}
 :=
 \Phi_{3,\sigma_{K3,E}}\bigl(D^b\Coh(K3\times E)\bigr),
 \qquad
 \sigma_{K3,E}=(K3,E).
\]
The Vol~I holographic datum assigns to a modular Koszul chiral algebra
$A$ the BRST trace scalar
$K_\Phi(A):=-c_{\mathrm{ghost}}(\BRST^{\mathrm{can}}(A))$
of Remark~\ref{rem:hdm-anchors}(W2). This is the $\Phi$-trace
conductor in the present comparison, distinct from the Verdier sum
$K^\kappa(A)=\kappa(A)+\kappa(A^!)$ of
Remark~\ref{rem:holo-master-theoremC-bucket}. The Vol~III universal
$\Phi$-trace identity specialises $K_\Phi$ at
$A_{K3\times E,\sigma}$ via the Borcherds singular-theta pushforward
$\Phi_{\mathrm{Borcherds}} \colon H^\bullet_{(2)}(\mathcal{D}_\Lambda,
\mathrm{vvMF}) \to \mathrm{Aut}(\Lambda)$
(Borcherds~\cite{Borcherds95},~\cite{Borcherds98}, the singular-theta
correspondence). At this stage-$2$ specialisation, the BKM closure is
the Gritsenko--Nikulin~\cite{GritsenkoNikulin98a} algebra
$\mathfrak{g}_{\Delta_5}$, whose Borcherds weight is read off its
primitive denominator $\Delta_5$:
\[
 \kappa_{\mathrm{BKM}}\bigl(\mathfrak{g}_{\Delta_5}\bigr)
 \;=\;
 \frac{c_\Lambda(0)}{2}
 \;=\;
 \mathrm{wt}(\Delta_5)
 \;=\;
 5.
\]
The chain-level universal $\Phi$-trace identity of Vol~III asserts,
on this witnessed branch,
\[
 K_\Phi(A_{K3\times E,\sigma})
 \;=\;
 \frac{c_\Lambda(0)}{2}
 \;=\;
 \mathrm{wt}(\Delta_5)
 \;=\;
 5,
\]
matching the Vol~I BRST trace conductor on the left with the Vol~III
Borcherds weight on the right through the Vol~II
holographic-envelope intermediary
(Proposition~\ref{prop:hdm-u-a-recoverability}(iii)--(iv)). The
obstruction boundary is exactly the one named in
Remark~\ref{rem:hdm-anchors}: no primary-source-complete
BRST-to-Borcherds comparison is proved outside the five verified
families $N \in \{1,2,3,4,6\}$, of which $K3\times E$
(the $N=1$ primitive denominator $\Delta_5$) is one;
the $\mathrm{wt}(\Delta_5) = 5$ answer is therefore an established
cross-volume datum, not a conjectural extension, on exactly this
family. The value~$5$ is not the Theorem~C B-family value~$8$ and
not the weight~$10$ of $\Phi_{10}^{\mathrm{un}}=\Delta_5^2$.
\end{remark}

\begin{remark}[Boundary-linear bulk versus global triangle]
\label{rem:hdm-boundary-linear-vs-global-triangle}
\index{boundary-linear bulk!Koszul-locus theorem}
\index{global triangle!conditional scope}
The \emph{global triangle} is the homotopy-coherent comparison diagram
\[
\xymatrix{
\cA \ar[r]^-{\mathbb H_X}
 \ar[dr]_-{\mathrm{bulk/centre}}
& \cA^! \ar[d]^-{\mathrm{line/Hochschild}} \\
& \cC
}
\]
in which the southwest arrow is the universal bulk object
$\cZ^{\mathrm{der}}_{\mathrm{ch}}(\cA)
\simeq \mathcal{C}^{\bullet}_{\mathrm{ch}}(\cA,\cA)$ and the southeast
arrow is the line-side presentation of the same slot whenever the
compact-generator/Hochschild comparison is proved. This triangle is
proved on two lanes only:
\begin{enumerate}[label=\textup{(\roman*)}]
\item chirally Koszul boundary-linear families of classes
 $\mathsf{G}$, $\mathsf{L}$, $\mathsf{C}$
 \textup{(}Vol~II, theorem ``Global triangle for boundary-linear
 families''\textup{)};
\item principal Drinfeld--Sokolov $\cW$-algebras at non-critical level,
 where the DS--Hochschild bridge gives the class-$\mathsf{M}$ chain
 closure
 \textup{(}Vol~II, the theorem ``Universal holography,
 class-$\mathsf{M}$ chain-level'' together with the corollary
 ``Class~$\mathsf{M}$ chain-level bulk reconstruction''\textup{)}.
\end{enumerate}
Outside those lanes the first failing strict arrow is the Verdier/Koszul
edge $\cA^{!}_\infty \to \cA^!$, whose strictness is exactly
Theorem~\ref{thm:bar-cobar-inversion-qi} on the Koszul locus.

For the admissible affine algebra
$\cA = L_{-1/2}(\mathfrak{sl}_2)$, let
$\eta_4 \in H^\bullet(\bar{B}^{\mathrm{ch}}(\cA))$ be the off-diagonal
bar class represented by the weight-$4$ vacuum singular vector of
Example~\ref{ex:admissible-sl2-failure}. That example shows that
$\eta_4$ survives the bar-cobar spectral sequence and produces a
non-trivial extension in
$\Omega^{\mathrm{ch}}(\bar{B}^{\mathrm{ch}}(\cA))$ with no preimage in
~$\cA$. Hence the counit
$\Omega^{\mathrm{ch}}(\bar{B}^{\mathrm{ch}}(\cA)) \to \cA$ is not a
quasi-isomorphism, the strict edge
$\cA^{!}_\infty \to \cA^!$ does not glue, and the global triangle
survives there only in the coderived ambient.

An independent obstruction on the nonlinear class-$\mathsf{M}$ side is
the quartic contact class
$\mathfrak{Q}^{\mathrm{contact}}_{\mathrm{Vir}}
= 10/[c(5c+22)]$
\textup{(}Vol~II, remark ``Class~$M$ gap: Virasoro and
$\cW_N$''\textup{)},
which blocks the boundary-linear reduction even when a DS comparison
theory exists.
\end{remark}

\begin{remark}[categorified averaging discipline]
\label{rem:hdm-ap-cy54-discipline}

\index{categorified averaging!right-adjoint discipline}
The Drinfeld double $U_\cA = \cA \bowtie \cA^{!}_\infty$ of
Section~\ref{sec:hdm-thesis} is constructed as the centraliser of
$\cA$ inside an ambient factorisation category, equivalently as
the categorified \emph{right adjoint} to the forgetful functor
\[
 \mathrm{Forget}\colon\;
 \mathrm{Fact}\text{-}\mathrm{Cat}_\cA \;\longrightarrow\;
 \mathrm{Fact}\text{-}\mathrm{Cat}.
\]
This is not averaging in the lay sense of summing over a finite
group orbit: the right-adjoint construction
\emph{equips} every object with a canonical $\cA$-bimodule
structure rather than projecting onto a $\cA$-invariant
sub-object. The slot~(4) verification
(Theorem~\ref{thm:hdm-seven-way-master}) and the slot~(5)
$\Theta_\cA$ recovery
(Theorem~\ref{thm:mc2-bar-intrinsic}) both depend on this
distinction: replacing the right-adjoint construction by a naive
average kills the spectral parameter of $r_\cA(z)$ and collapses
$\Theta_\cA$ onto its scalar genus-$0$ trace.
\end{remark}

\begin{remark}[Theorem~C bucket $\{0, 8, 13, 250/3, 98/3\}$ and the $K^{\kappa_{\mathrm{ch}}} = 8$ B-family value]
\label{rem:holo-master-theoremC-bucket}
\index{Theorem~C!allowed-values bucket $\{0, 8, 13, 250/3, 98/3\}$}
\index{B-family!$K^{\kappa_{\mathrm{ch}}} = 8$ Bruinier--Heegner reciprocity}
The scalar Verdier sum
$K^\kappa(\cA):=\kappa(\cA)+\kappa(\cA^!)$ attached to the
Drinfeld double $U_\cA$ by
Proposition~\ref{prop:hdm-u-a-recoverability}(ii) inherits from
Theorem~C (the derived-centre complementarity theorem of
Part~\ref{part:standard-landscape}) two nested finite buckets. On the
four standard-landscape lanes $\mathsf G / \mathsf L / \mathsf C /
\mathsf M$ the value is forced into the classical ceiling
\[
 K^\kappa(\cA)
 \;\in\;
 \{0,\; 13,\; 250/3,\; 98/3\}.
\]
The canonical five-element bucket
\[
 K^\kappa(\cA)
 \;\in\;
 \{0,\; 8,\; 13,\; 250/3,\; 98/3\}
\]
adjoins the Vol~III $\mathsf B$-row witness. The value~$0$ is the
free-field abelian locus
(Heisenberg, class~$\mathsf G$, self-dual under $k \mapsto -k$). The
values $13$ (Virasoro, bosonic string), $250/3$ (principal
$\mathcal{W}_N$ universal family), and $98/3$ (matter + ghost
balanced families) are the Theorem~C generic-$\mathsf M$ outputs
recorded in Chapter~\ref{chap:universal-conductor}. The extremal
entry~$8$ is not a fifth $\mathsf G/\mathsf L/\mathsf C/\mathsf M$
value. It is the B-family rigidity value of the Mukai-enhanced K3
Heisenberg witness
$K^{\kappa_{\mathrm{ch}}}(\mathrm{Mukai}(K3))
=2c_+(\mathrm{II}_{4,20})=8$, transported through Bruinier
Heegner-reciprocity~\cite{Bruinier02}. Vol~III keeps this scalar
separate from the Borcherds weight
$\kappa_{\mathrm{BKM}}(\Delta_5)=c_1(0)/2=\mathrm{wt}(\Delta_5)=5$
and from the Oberdieck--Pixton normalisation
\[
 D_5=64^{-1}\Delta_5,\qquad
 \Phi_{10}^{\mathrm{OP}}=D_5^2=4096^{-1}\Delta_5^2,\qquad
 Z_{\mathrm{OP/DT}}=-D_5^{-2}=-4096\,\Delta_5^{-2}.
\]
The cross-volume bridge
relates the three faces by the singular-theta / Heegner
correspondence; it does not identify $8$ with a doubled or averaged
Borcherds weight. No sixth value is admitted on the chirally Koszul
boundary-linear locus where slot~(2) of
Proposition~\ref{prop:hdm-u-a-recoverability} is strict.
\end{remark}

%%%%%%%%%%%%%%%%%%%%%%%%%%%%%%%%%%%%%%%%%%%%%%%%%%%%%%%%%%%%%%%%%%%%%%%%%%%%%
\section{The cross-volume composite triangle:
$\av \circ \Obs^\partial \circ \Phi_{\mathrm{hol}} \circ \Phi_{d,\sigma}$}
\label{sec:hdm-cross-volume-triangle}
%%%%%%%%%%%%%%%%%%%%%%%%%%%%%%%%%%%%%%%%%%%%%%%%%%%%%%%%%%%%%%%%%%%%%%%%%%%%%

The three volumes individually furnish the four arrows. The Vol~III
arrow is a family; after a stage-$2$ specialisation datum
$\sigma=(\Sigma_{d-1},C_X)$ with $C_X\simeq X$ is fixed, its curve
valued member is denoted $\Phi_{d,\sigma}$.
\[
 \begin{array}{rcl}
 \Phi_{d,\sigma} & \colon &
 \CY_d\text{-}\Cat^{\mathrm{sm,prop}}
 \longrightarrow \ChirAlg_X^{E_1}
 \quad
 \text{\textup{(}Vol~III, Theorem ``Platonic CY-to-chiral correspondence''\textup{);}}
 \\[3pt]
 \Phi_{\mathrm{hol}} & \colon &
 \ChirAlg^{\omega,\mathrm{BL,adm}}_X
 \longrightarrow \HTQFT_{X \times \R}
 \quad
 \text{\textup{(}Vol~II,
 Theorem ``Universal holography functor''\textup{);}}
 \\[3pt]
 \Obs^\partial & \colon &
 \HTQFT_{X \times \R}
 \longrightarrow \ChirAlg_X
 \quad
 \text{\textup{(}Costello--Gwilliam Vol.~II,
 Thm.~7.3.7~\cite{CG21}\textup{);}}
 \\[3pt]
 \av & \colon &
 \mathfrak{g}^{E_1}_\cA
 \twoheadrightarrow \mathfrak{g}^{\mathrm{mod}}_\cA
 \quad
 \text{\textup{(}Vol~I,
 Definition~\textup{\ref{def:e1-modular-convolution}}\textup{).}}
 \end{array}
\]
The first two arrows are constructed in volumes that adopt different
chain models, different homotopy frames, and different scope
hypotheses; the third is the Costello--Gwilliam restriction-along-
boundary; the fourth is the Vol~I $\Sigma_n$-coinvariant. The
question of \emph{whether the four arrows compose at chain level}
on the boundary-linear locus, and whether the composite is
$E_1$-functorial in the source, has been left implicit across the
trilogy. The following theorem records the precise scope on which
the chain-level composite is functorial, identifies its image
class-by-class on the four shadow archetypes, and isolates two
hypotheses that the composite \emph{requires} but the individual
arrows do not.

\begin{theorem}[Cross-volume composite triangle on the boundary-linear
$\Phi_{d,\sigma}$-image; \ClaimStatusConditional]
\label{thm:hdm-triangle-composite}
\index{cross-volume triangle!$\av\circ\Obs^\partial\circ\Phi_{\mathrm{hol}}\circ\Phi_{d,\sigma}$|textbf}
\index{composite functor!chain-level scope|textbf}
\index{triangle hub!Vol~I/II/III closure|textbf}
Fix a smooth projective complex curve $X$ of genus $g \geq 0$, a CY
dimension $d \in \{2, 3\}$, and a Vol~III stage-$2$ specialisation
datum $\sigma=(\Sigma_{d-1},C_X\simeq X)$. Let $\cC$ be a smooth proper
$\CY_d$-category satisfying the hypotheses
\textup{(H1)--(H3)} of Vol~III
``Platonic CY-to-chiral correspondence at dimension $d$,'' together
with the two further hypotheses:
\begin{enumerate}[label=\textup{(H\arabic*$_\triangle$)}, start=4]
\item\label{H4-triangle} \emph{Boundary-linear image hypothesis.} The
 stage-$2$ chiral algebra
 $A_{\cC,\sigma}:=\Phi_{d,\sigma}(\cC)
 =\mathrm{Sp}^{\mathrm{ch}}_{\sigma}(\Phi^{\mathrm{FA}}_d(\cC))$, viewed as an
 $E_1$-chiral algebra on $X$, lies in the boundary-linear
 source category $\ChirAlg^{\omega,\mathrm{BL,adm}}_X$ of Vol~II
 Definition ``Boundary-linear chiral algebra with conformal
 vector'': a conformal vector $\omega$ exists, the level is
 non-critical, and the modular Maurer--Cartan element
 $\Theta_{A_{\cC,\sigma}}$ admits the strict Lagrangian splitting on the
 $\cA$-face of the Koszul pair.
\item\label{H5-triangle} \emph{Chiral Koszul hypothesis on $A_{\cC,\sigma}$.} The
 chiral algebra $A_{\cC,\sigma}$ is chirally Koszul on $X$ in the sense of
Theorem~\textup{\ref{thm:bar-cobar-inversion-qi}}, so that the
 Vol~I averaging-compatibility statement
 \textup{(}Proposition~\textup{\ref{prop:hdm-averaging-compatibility}}\textup{)}
 applies to $A_{\cC,\sigma}$.
\end{enumerate}
Then on the full subcategory
\[
 \CY_d^{\,\Phi\text{-BL}}
 \;:=\;
 \bigl\{\, \cC \in \CY_d\text{-}\Cat^{\mathrm{sm,prop}}
 \;:\;
 \Phi_{d,\sigma}(\cC) \text{ satisfies}~\textup{\ref{H4-triangle}}\text{,}~\textup{\ref{H5-triangle}}\,\bigr\}
\]
the four arrows compose at chain level into a well-defined map on
objects
\begin{equation}\label{eq:triangle-composite-objects}
 \widetilde\Phi_{d,\sigma}^{\,\mathrm{mod}}
 \;:=\;
 \av \;\circ\; \Obs^\partial \;\circ\; \Phi_{\mathrm{hol}} \;\circ\; \Phi_{d,\sigma}
 \;\colon\;
 \CY_d^{\,\Phi\text{-BL}}
 \longrightarrow
 \mathfrak{g}^{\mathrm{mod}}_\bullet
\end{equation}
with the following four properties:
\begin{enumerate}[label=\textup{(\roman*)}]
\item \emph{Chain-level inversion of the middle two arrows.} On
 every $\cC \in \CY_d^{\,\Phi\text{-BL}}$,
 \begin{equation}\label{eq:triangle-middle-inversion}
 \Obs^\partial \;\circ\; \Phi_{\mathrm{hol}} \bigl(\Phi_{d,\sigma}(\cC)\bigr)
 \;\simeq\; \Phi_{d,\sigma}(\cC)
 \end{equation}
 as $E_1$-chiral algebras on $X$, by Vol~II Theorem ``Universal
 holography functor,'' boundary face. The identification is
 chain-level, not merely up to the level of underlying objects, by
 the Vol~II Theorem ``Bulk-to-boundary OPE consistency'' which
 records the on-the-nose OPE-preservation of the restriction map.
\item \emph{Reduction of the composite to a Vol~I averaging on
 $\Phi_{d,\sigma}(\cC)$.} On the full subcategory $\CY_d^{\,\Phi\text{-BL}}$,
 \begin{equation}\label{eq:triangle-reduction}
 \widetilde\Phi_{d,\sigma}^{\,\mathrm{mod}}(\cC)
 \;\simeq\;
 \av\bigl(\Phi_{d,\sigma}(\cC)\bigr)
 \quad\text{in } \mathfrak{g}^{\mathrm{mod}}_\bullet,
 \end{equation}
 with the right-hand side computed in
 $\mathfrak{g}^{E_1}_{\Phi_{d,\sigma}(\cC)}$ and projected to
 $\mathfrak{g}^{\mathrm{mod}}_{\Phi_{d,\sigma}(\cC)}$ by
 $\Sigma_n$-coinvariants
 \textup{(}Definition~\textup{\ref{def:e1-modular-convolution}}\textup{).}
 The chain-level shift correction across the
 $\Phi_{d,\sigma} \to \Phi_{\mathrm{hol}}$ interface is precisely the constant
 desuspension $[1]$ of the bar differential
 \textup{(}Vol~II Theorem ``Bar differential = holomorphic
 factorisation'' part~(i): the holomorphic factorisation map at
 coincident points\textup{)}, absorbed into the bar coalgebra
 $B^{\mathrm{ch}}(\Phi_{d,\sigma}(\cC)) = T^c(s^{-1}\overline{\Phi_{d,\sigma}(\cC)})$
 on the Vol~I side without further correction.
\item \emph{Class identification on the boundary-linear locus.} On
 $\CY_d^{\,\Phi\text{-BL}}$, the composite
 $\widetilde\Phi_{d,\sigma}^{\,\mathrm{mod}}(\cC)$ at binary degree $n=2$
 lands in the four shadow archetypes G/L/C/M of Vol~I according to
 the boundary-linear class of $\Phi_{d,\sigma}(\cC)$. The K3 entry
 below lies in the smooth-proper source of the theorem. The elliptic
 and $\C^3$ entries are calibration checks in the CY$_1$ and framed
 local CY$_3$ extensions; they are not objects of
 $\CY_d^{\,\Phi\text{-BL}}$ unless those extensions are explicitly
 added to the source.
 \begin{itemize}[leftmargin=1.2em]
 \item \emph{$d = 2$, $\cC = D^b(\Coh(\mathrm{K3}))$:}
 $\Phi_{2,\sigma}(\cC) = V_{\widetilde\Lambda_{K3}}$ \textup{(}rank-$24$
 Mukai lattice VOA; Vol~III equation
 ``$\Phi_{2,\sigma}(D^b(\Coh(\mathrm{K3}))) = V_{\widetilde\Lambda_{K3}}$''\textup{)},
 lattice/Heisenberg, hence \emph{class~G}, and
 $\widetilde\Phi_{2,\sigma}^{\,\mathrm{mod}}(D^b\Coh(\mathrm{K3}))$ at
 $n=2$ equals
 $r_{V_{\widetilde\Lambda_{K3}}}(z) = (1/z)\,
 \sum_{i=1}^{24} J_i \otimes J_i$, the Mukai-Heisenberg residue
 with no Sugawara shift, by
 Proposition~\textup{\ref{prop:hdm-averaging-compatibility}}\textup{(i).}
 \item \emph{CY$_1$ calibration, $\cD = D^b(\Coh(E))$
 \textup{(}elliptic curve\textup{):}
 $\Phi_{1,\sigma}(\cD) = V_{\Lambda_E}$
 \textup{(}rank-$2$ elliptic lattice VOA\textup{);} class~G again;
 identity at $n=2$.}
 \item \emph{Framed local CY$_3$ calibration,
 $\cD = \CoHA(\C^3) = Y^+(\widehat{\fgl}_1)$:}
 the Drinfeld-centre passage of Vol~III \textup{(}property
 \textup{(U3)} of ``Platonic CY-to-chiral'': $\cZ(\Rep^{E_1}(A))
 \simeq \Rep^{E_2}(\cZ^{\mathrm{der}}_{\mathrm{ch}}(A))^{\mathrm{centred}}$\textup{)}
 produces $\cW_{1+\infty}$ at $c=1$, whose Heisenberg subalgebra
 \textup{(}generator $J_0 = \alpha_{-1}\mathbf{1}$\textup{)} is class~G;
 the Virasoro stress tensor inside $\cW_{1+\infty}$ is class~M at
 $c=1$. The framed-local analogue of
 $\widetilde\Phi_{3,\sigma}^{\,\mathrm{mod}}$ on $\CoHA(\C^3)$ has a
 binary residue with two pieces: the Heisenberg piece
 $J_0 \otimes J_0 / z$ \textup{(}class~G, $\av = \id$ by~\textup{(i)}
 of Proposition~\textup{\ref{prop:hdm-averaging-compatibility}}\textup{)} and
 the Virasoro piece $(1/2)/z^3 + 2T/z$ \textup{(}class~M at $c = 1$,
 $\av = \id$ by~\textup{(iv)} of the same proposition\textup{).}
 \end{itemize}
 In the K3 theorem case and in the two calibration checks, the kernel
 $\ker(\av)|_{r}$ vanishes at degree~$n=2$ on the corresponding
 composite; the genuine
 averaging-non-trivial information first appears at degree
 $n \geq 3$ in the higher shadow $S_3, S_4, \dots$ of the image
 algebra.
\item \emph{Functoriality scope.} On objects, the composite
 $\widetilde\Phi_{d,\sigma}^{\,\mathrm{mod}}$ is a well-defined map
 $\CY_d^{\,\Phi\text{-BL}} \to \mathfrak{g}^{\mathrm{mod}}_\bullet$.
 On morphisms, $\widetilde\Phi_{d,\sigma}^{\,\mathrm{mod}}$ is
 \emph{conditionally} functorial: a CY-morphism
 $f \colon \cC \to \cC'$ in $\CY_d^{\,\Phi\text{-BL}}$ is sent to
 the chain-level morphism
 $\widetilde\Phi_{d,\sigma}^{\,\mathrm{mod}}(f) =
 \av(\Obs^\partial(\Phi_{\mathrm{hol}}(\Phi_{d,\sigma}(f))))$
 if and only if $\Phi_{d,\sigma}(f)$ is established to be a morphism of
 chiral algebras \textup{(}Vol~III property \textup{(U2)} of
 ``Platonic CY-to-chiral,'' open in general\textup{).} The other
 three arrows ($\Phi_{\mathrm{hol}}$, $\Obs^\partial$, $\av$) are
 unconditionally functorial on their domains by Vol~II
	 Theorem ``Universal holography functor,'' the construction of
 \textup{(}functoriality follows from functoriality of the
 Costello--Gwilliam envelope and the Hochschild-cochain
 construction\textup{)}, Costello--Gwilliam~\cite{CG21}\,Thm.~7.3.7,
 and Definition~\textup{\ref{def:e1-modular-convolution}}
 \textup{(}naturality of $\Sigma_n$-coinvariants in the chiral-algebra
 source\textup{).} At $d = 2$, functoriality of $\Phi_2$ is verified
 concretely on the Mukai transform $\mathrm{K3} \to \mathrm{K3}$
 \textup{(}Vol~III ``$\Phi$ on K3 explicit''\textup{);}
 elsewhere it is per-$d$ open on the source side, hence open for
 the composite.
\end{enumerate}
The two extra hypotheses are sharp. If
\textup{\ref{H4-triangle}} fails, then $\Phi_{d,\sigma}(\cC)$ need not lie in
the source of $\Phi_{\mathrm{hol}}$, and on nonlinear class~$\mathsf{M}$
the quartic contact class $10/[c(5c+22)]$ gives the explicit failure of
boundary-linearity outside the DS lane. If \textup{\ref{H5-triangle}}
fails, then the Vol~I averaging-compatibility theorem need not apply:
the bar spectral sequence need not collapse, the strict edge
$(\Phi_{d,\sigma}(\cC))^{!}_\infty\to (\Phi_{d,\sigma}(\cC))^!$ can fail, and the admissible
off-diagonal class $\eta_4$ is the concrete witness that the strict
binary comparison with the pinned six-slot package may not descend.
\end{theorem}

\begin{proof}
The four properties are established sequentially.

\smallskip
\emph{\textup{(i)} Chain-level middle inversion.} By Vol~II
Theorem ``Universal holography functor,'' part~(i) (boundary face),
$\Obs^\partial(\Phi_{\mathrm{hol}}(\cA)) \simeq \cA$ as $E_1$-chiral
algebras on $X$ for every $\cA \in \ChirAlg^{\omega,\mathrm{BL,adm}}_X$.
Hypothesis~\textup{\ref{H4-triangle}} places $\Phi_{d,\sigma}(\cC)$ in the
source of $\Phi_{\mathrm{hol}}$, so the identification applies with
$\cA = \Phi_{d,\sigma}(\cC)$. Vol~II's Theorem ``Bulk-to-boundary OPE
consistency'' sharpens the identification from a quasi-isomorphism
of underlying objects to a chain-level identification of
$E_1$-chiral algebra structures \textup{(}part~(c) of that
theorem: the algebra structure on
$\Obs^\partial(\Phi_{\mathrm{hol}}(\cA))$ agrees strictly at
chain level with the chiral product on $\cA$\textup{).} This gives
\eqref{eq:triangle-middle-inversion}.

\smallskip
\emph{\textup{(ii)} Reduction to Vol~I averaging.}
Apply $\av$ to both sides of~\eqref{eq:triangle-middle-inversion}.
The left-hand side is, by definition,
$\widetilde\Phi_{d,\sigma}^{\,\mathrm{mod}}(\cC)$. The right-hand side is
$\av(\Phi_{d,\sigma}(\cC))$, well-defined in
$\mathfrak{g}^{\mathrm{mod}}_{\Phi_{d,\sigma}(\cC)}$ because
hypothesis~\textup{\ref{H5-triangle}} places $\Phi_{d,\sigma}(\cC)$ in the
source of the Vol~I averaging-compatibility proposition. The
$[1]$-shift interface between Vol~III's bar coalgebra
$\mathrm{CC}_\bullet(\cC)$ \textup{(}cyclic homological grading\textup{)}
and Vol~I's $T^c(s^{-1}\overline{\cA})$ \textup{(}cohomological
grading with desuspension\textup{)} is reconciled by the Vol~III
universal property $B^{\mathrm{ord}}(\Phi_{d,\sigma}(\cC)) \simeq
\mathrm{CC}_\bullet(\cC)$ \textup{(}Vol~III property
\textup{(U1)} of ``Platonic CY-to-chiral''\textup{)} together with
the Vol~II identification of the bar differential with the
holomorphic factorisation map at coincident points
\textup{(}Theorem ``Bar differential = holomorphic factorisation,''
part~(i), where the $[1]$-shift accounts precisely for the
desuspension $s^{-1}$ in $T^c(s^{-1}\overline{\cA})$\textup{).}
No additional shift correction is needed.

\smallskip
\emph{\textup{(iii)} Class identification.} For the smooth-proper K3
instance, identify the boundary-linear class of
$\Phi_{d,\sigma}(\cC)$ from Vol~III, then apply the corresponding clause
of Proposition~\textup{\ref{prop:hdm-averaging-compatibility}}. The
elliptic and $\C^3$ calculations use the same binary argument in the
CY$_1$ and framed local CY$_3$ extensions, respectively.

\textit{$d=2$, $\cC = D^b(\Coh(\mathrm{K3}))$.}
By Vol~III, $\Phi_{2,\sigma}(\cC) = V_{\widetilde\Lambda_{K3}}$ with
$\widetilde\Lambda_{K3} = U^4 \oplus E_8(-1)^2$, rank~$24$,
signature $(4, 20)$. As a lattice VOA, $V_{\widetilde\Lambda_{K3}}$
is class~G \textup{(}generated by $24$ commuting Heisenberg
currents $J_1, \dots, J_{24}$ with double-pole OPE
$J_i(z) J_j(w) \sim \delta_{ij}/(z-w)^2$ in the orthogonal basis
diagonalising the Mukai pairing, plus lattice-vertex generators
that are weight-$1$ exponentials of integer linear combinations
of the Heisenberg currents\textup{).} The binary collision
residue is therefore
$r_{V_{\widetilde\Lambda_{K3}}}(z) = (1/z) \sum_{i=1}^{24}
J_i \otimes J_i$, and
Proposition~\textup{\ref{prop:hdm-averaging-compatibility}}\textup{(i)}
applied to each $J_i$-summand gives $\av(r) = r$ at degree~$2$
with no Sugawara correction \textup{(}Heisenberg is abelian, no
adjoint Casimir contribution\textup{).}

\textit{CY$_1$ calibration, $\cD = D^b(\Coh(E))$.} Identical reasoning with
the rank-$2$ elliptic lattice VOA $V_{\Lambda_E}$ in place of
$V_{\widetilde\Lambda_{K3}}$.

\begin{remark}[Bi-based Ran-space structure of $\mathcal{H}(T)$ at
$K3 \times E$: $\mathrm{Ran}(E_{24}^{\mathrm{nod,sm}})$ and
$\overline{A}_2$]
\label{rem:holo-master-bi-based-ran}
\index{Ran space!bi-based at $K3 \times E$}
\index{averaging composite!$\mathrm{Torelli}_{K3} \circ \mathrm{KugaSatake}
\circ j_{\mathrm{Kodaira}}$}
The triangle-composite map of the K3 clause above asserts that the
six-tuple
$\mathcal{H}(T) = (\cA, \cA^!, \cC, r_\cA(z), \Theta_\cA,
\nabla^{\mathrm{hol}})$ specialises at
$\cA = V_{\widetilde\Lambda_{K3}}$ to a \emph{bi-based} object, in
which the three chiral-algebra components
$(\cA, \cA^!, \cC)$ and the three parameter components
$(r_\cA(z), \Theta_\cA, \nabla^{\mathrm{hol}})$ live on distinct
parameter spaces and are coupled only through the averaging
composite
\[
 \mathrm{Torelli}_{K3} \,\circ\, \mathrm{KugaSatake}
 \,\circ\, j_{\mathrm{Kodaira}}
 \;\colon\;
 \overline{A}_2
 \;\longrightarrow\;
 \mathrm{Sp}_4(\mathbb{Z}) \backslash \mathfrak{H}_2
 \;\xrightarrow{\ \mathrm{KS}\ }\;
 \Gamma_{K3} \backslash \Omega_{K3}
 \;\xrightarrow{\ \mathrm{Torelli}\ }\;
 \mathcal{M}_{K3}^{\mathrm{alg}},
\]
where $j_{\mathrm{Kodaira}}$ embeds $\overline{A}_2$
(the Deligne--Mumford toroidal compactification of
$\mathcal{A}_2 = \mathrm{Sp}_4(\mathbb{Z}) \backslash \mathfrak{H}_2$)
as the Kodaira fibre locus of the degenerate Kuga--Satake family,
$\mathrm{KugaSatake}$ is the Kuga--Satake
construction~\cite{KugaSatake67} assigning to a K3 lattice a
polarised abelian surface, and $\mathrm{Torelli}_{K3}$ is the
Piatetski-Shapiro--Shafarevich Torelli
theorem~\cite{PiatetskiShapiroShafarevich71} for K3 surfaces.
On the Vol~I side, the chiral-algebra components
$(V_{\widetilde\Lambda_{K3}},
V_{\widetilde\Lambda_{K3}}^!,
\mathcal{C}^\bullet_{\mathrm{ch}}(V_{\widetilde\Lambda_{K3}},
V_{\widetilde\Lambda_{K3}}))$ live on
$\mathrm{Ran}(E_{24}^{\mathrm{nod,sm}})$, the Ran space of the
$24$-marked nodal-smooth elliptic family (Kodaira $I_1$ or smooth
fibre at each of the $24$ Euler-characteristic zeroes on the K3),
carrying the standard E$_1$-factorisation structure of
Part~\ref{part:bar-complex}; on the Vol~III side, the parameter
components $(r_\cA(z), \Theta_\cA, \nabla^{\mathrm{hol}})$ live on
the $\overline{A}_2$-base through the Gritsenko--Nikulin denominator
$\Delta_5$ and its D-module avatar $\mathcal{N}_{\Delta_5}$ of
Vol~III. The averaging composite glues the two lanes: the
Vol~I $\Sigma_n$-coinvariants projection $\av$ restricted to the
boundary-linear image of $V_{\widetilde\Lambda_{K3}}$ factors
through the Torelli--Kuga-Satake--Kodaira chain, identifying the
Ran-space factorisation output on the left with the
$\overline{A}_2$-modular parameter on the right. The
bi-basing is essential: there is no single variety carrying both
lanes simultaneously, and the averaging composite is the specified
map making slot~(4) (spectral $R$-matrix) and slot~(6)
(shadow connection) agree with slot~(1) (boundary algebra) and
slot~(3) (bulk slot). This is the $K3 \times E$-specialised form
of the general bi-based Ran-space picture governing every
class-$\mathsf{G}$ lattice VOA on a complex $2$-torus times an
elliptic fibration.
\end{remark}

\textit{Framed local CY$_3$ calibration,
$\cD = \CoHA(\C^3) = Y^+(\widehat{\fgl}_1)$.}
The native framed-local image $\Phi_{3,\sigma}(\CoHA(\C^3)) =
Y^+(\widehat{\fgl}_1)$ is
$E_1$-chiral; the Drinfeld-centre passage of Vol~III property
\textup{(U3)} of ``Platonic CY-to-chiral'' produces
$\cZ(\Rep^{E_1}(Y^+(\widehat{\fgl}_1))) \simeq \Rep^{E_2}(\cW_{1+\infty})$
at $c = 1$. The chain-level realisation that $\Phi_{\mathrm{hol}}$
sees on the boundary is the algebra $\cW_{1+\infty}$ itself
\textup{(}via Procha\'zka--Rapca\'k presentation, identified with
$Y(\widehat{\fgl}_1)$ at $c=1$\textup{).} The OPE of $\cW_{1+\infty}$
at $c = 1$ contains the Heisenberg subalgebra generated by
$J_0(z) = \alpha_{-1} \mathbf{1}$ \textup{(}class~G, double-pole
$J_0(z) J_0(w) \sim 1/(z-w)^2$\textup{)} and the Virasoro stress
tensor $T(z)$ \textup{(}class~M, $T(z)T(w) \sim (1/2)/(z-w)^4
+ \cdots$ at $c = 1$\textup{).} Both pieces are independently
$\Sigma_2$-invariant on their canonical chain models
\textup{(}Heisenberg by abelian symmetry, Virasoro by the
$d\log$-absorption argument of
Proposition~\textup{\ref{prop:hdm-averaging-compatibility}}\textup{(iv)}\textup{),}
so $\av$ acts as the identity at degree~$2$ on both pieces.
\textit{The framed-local composite therefore lands in
$\{\mathrm{class~G}\} \sqcup \{\mathrm{class~M}\}$ at $\C^3$, with
no kernel at the binary level.}

\smallskip
\emph{\textup{(iv)} Functoriality scope.}
Functoriality on morphisms decomposes into the four arrows. By
Costello--Gwilliam~\cite{CG21}\,Thm.~7.3.7,
$\Obs^\partial \colon \HTQFT_{X \times \R} \to \ChirAlg_X$ is a
functor of $\infty$-categories \textup{(}left adjoint to the
constant-along-$\R$ inclusion\textup{).} $\Phi_{\mathrm{hol}}$ is
functorial by Vol~II Theorem ``Universal holography functor'': the
Costello--Gwilliam envelope $\cF^{\mathrm{hol}}_{(-)}$
is functorial in the chiral-algebra input
\textup{(}Costello--Gwilliam Vol.~II,~Thm.~3.6.1\textup{),} and the
Hochschild-cochain factorisation $\cH(-)$ is functorial in the
chiral-algebra input by naturality of Hochschild cochains in
algebra morphisms; the Dunn tensor product
$\otimes_{\mathrm{Dunn}}$ then assembles the two into a functor on
$\ChirAlg^{\omega,\mathrm{BL,adm}}_X$. The averaging map $\av$ is
natural in the chiral-algebra source by construction
\textup{(}$\Sigma_n$-coinvariants are a colimit construction in the
source $\mathfrak{g}^{E_1}$, hence functorial in the
source\textup{).}

The remaining arrow $\Phi_{d,\sigma}$ is functorial \emph{conditionally}:
at $d \geq 2$ it is per-$d$ functorial on objects unconditionally
\textup{(}Vol~III ``Platonic CY-to-chiral,'' property
\textup{(U1)}\textup{),} and per-$d$ functorial on morphisms
exactly when the Vol~III conjecture on $\Phi_{d,\sigma}$-functoriality
holds for the specific morphism. Hence the composite is functorial
on objects in $\CY_d^{\,\Phi\text{-BL}}$ unconditionally
\textup{(}as the four arrows compose on objects\textup{),} and
functorial on morphisms exactly when $\Phi_{d,\sigma}$ is. At $d = 2$ on
the Mukai transform $\mathrm{K3} \to \mathrm{K3}$, Vol~III
``$\Phi$ on K3 explicit'' verifies $\Phi_{2,\sigma}(\mathrm{Muk})$
explicitly, which gives one concrete morphism on which the
composite is verified end-to-end.
\end{proof}

\begin{remark}[Why the boundary-linear hypothesis is required and
not deducible from $\Phi_{d,\sigma}$ alone]
\label{rem:hdm-triangle-bl-required}
\index{boundary-linear hypothesis!required for $\Phi_{\mathrm{hol}}$ composition}
Hypothesis~\textup{\ref{H4-triangle}} cannot be discarded. The
native output of $\Phi_{d,\sigma}$ at $d \geq 3$ becomes an $E_1$-chiral
algebra on the curve only after the stage-$2$ specialisation
$\sigma$ is chosen, and
$\Phi_{\mathrm{hol}}$ is defined only on
$\ChirAlg^{\omega,\mathrm{BL,adm}}_X$, the boundary-linear sub-category
\textup{(}Vol~II Definition ``Boundary-linear chiral algebra with
conformal vector''\textup{).} Boundary-linearity is a chain-level
structural property of the modular MC element $\Theta_\cA$
\textup{(}strict Lagrangian splitting on the $\cA$-face of the
Koszul pair\textup{)} and is not implied by either smoothness/
properness of the source CY category or by the universal $\Phi_{d,\sigma}$
properties \textup{(U1)--(U4)} of Vol~III ``Platonic
CY-to-chiral.'' The three classes G, L, class~M-on-the-DS-lane are
known to be boundary-linear; class~M beyond the DS lane is not, by
the quartic contact obstruction
$\mathfrak{Q}^{\mathrm{contact}}_{\mathrm{Vir}} = 10/[c(5c+22)]$
of Vol~II
\textup{(}Remark~\textup{\ref{rem:hdm-boundary-linear-vs-global-triangle}}
above\textup{).} The two listed instances at $d=2$
\textup{(}K3, elliptic curve\textup{)} and the one listed instance
at $d=3$ \textup{(}$\C^3$\textup{)} are inside the boundary-linear
locus because they are class~G \textup{(}Heisenberg/lattice\textup{)}
or class~G $\sqcup$ class~M-on-the-DS-lane
\textup{(}$\cW_{1+\infty}$ at $c=1$ via the Drinfeld-centre passage,
which preserves the DS-lane property\textup{).} For a generic
class-M $\Phi_{d,\sigma}(\cC)$ outside the DS lane,
\textup{\ref{H4-triangle}} fails and the composite
$\widetilde\Phi_{d,\sigma}^{\,\mathrm{mod}}$ is not defined at chain level
via $\Phi_{\mathrm{hol}}$; an alternative chain-level extension via
the global triangle of
Remark~\textup{\ref{rem:hdm-boundary-linear-vs-global-triangle}}
would be needed.
\end{remark}

\begin{remark}[Scope discipline: what the theorem does
and does not assert]
\label{rem:hdm-triangle-composite-scope}
\index{cross-volume composite!chain-level versus $(\infty,1)$-categorical scope}
\index{$(\infty,1)$-categorical composite!separately required}
Theorem~\textup{\ref{thm:hdm-triangle-composite}} asserts the
\emph{chain-level} composite $\widetilde\Phi_{d,\sigma}^{\,\mathrm{mod}}$ on
$\CY_d^{\,\Phi\text{-BL}}$ as a map on objects, with conditional
functoriality on morphisms tied to the Vol~III conjecture on
$\Phi_{d,\sigma}$-functoriality. It does \emph{not} assert that the same
composite assembles into a single $(\infty,1)$-functor across all
$d$ \textup{(}the target $\mathfrak{g}^{\mathrm{mod}}_\bullet$
depends on the source class through the bar coalgebra of
$\Phi_{d,\sigma}(\cC)$, and the per-$d$ target-mismatch obstruction of Vol~III
``Platonic CY-to-chiral'' carries through to the
composite\textup{).} Within a fixed $d$, the
$(\infty,1)$-categorical lift requires showing
$\Obs^\partial \circ \Phi_{\mathrm{hol}}$ is an
$(\infty,1)$-functor on $\ChirAlg^{\omega,\mathrm{BL,adm}}_X$
\textup{(}true by Vol~II Theorem ``Universal holography functor''
together with the Costello--Gwilliam $\infty$-categorical
formulation of $\Obs^\partial$\textup{)} and that $\av$ is an
$(\infty,1)$-categorical right Kan extension along the
$\Sigma_\bullet$-action \textup{(}true by
Definition~\textup{\ref{def:e1-modular-convolution}} reread in
$(\infty,1)$-categorical language\textup{).} The
$(\infty,1)$-lift of the composite at fixed $d$ on
$\CY_d^{\,\Phi\text{-BL}}$ thus follows from the $(\infty,1)$-lift
of $\Phi_{d,\sigma}$ at fixed $d$ and fixed $\sigma$, which is established only on the smooth
proper locus per the Vol~III universal property \textup{(U1)--(U4)}
\textup{(}objects\textup{);} the morphism action remains
conjectural per~\textup{(U2)}, whose open input is exactly the missing
proof that $\Phi_{d,\sigma}(f)$ is a chiral-algebra morphism for a general
CY-morphism~$f$ after the specialisation~$\sigma$.

In particular, the theorem does \emph{not} circumvent
the universal-property gap at $\Phi_{d,\sigma}$. It \emph{names} where the
chain-level composite is forced to live \textup{(}the boundary-linear sub-locus of the
$\Phi_{d,\sigma}$-image\textup{),} \emph{names} the conditional
functoriality \textup{(}tied to the Vol~III conjecture on
$\Phi_{d,\sigma}$-functoriality\textup{),} and \emph{names} the chain-level
shift interface at $\Phi_{d,\sigma} \to \Phi_{\mathrm{hol}}$ \textup{(}the
desuspension $[1]$ matched to the holomorphic factorisation degree
shift, by Vol~II Theorem ``Bar differential = holomorphic
factorisation,'' part~(i)\textup{).} Beyond the smooth proper
locus, beyond the boundary-linear locus, or off the chiral Koszul
locus, the composite is not asserted at chain level by this
theorem.
\end{remark}

\begin{remark}[The closed triangle: how the three volumes' bridges
compose]
\label{rem:hdm-triangle-closed}
\index{closed triangle!Vol~I/II/III bridges compose}
\index{three-volume composition!chain-level}
The three explicit cross-volume bridges of the trilogy combine into
the closed chain-level diagram on
$\CY_d^{\,\Phi\text{-BL}}$:
\[
\xymatrix@C=1.4cm{
\CY_d\text{-}\Cat^{\mathrm{sm,prop}}
 \ar[r]^-{\Phi_{d,\sigma}}_-{\text{Vol~III}}
& \ChirAlg^{\omega,\mathrm{BL,adm}}_X
 \ar[r]^-{\Phi_{\mathrm{hol}}}_-{\text{Vol~II}}
& \HTQFT_{X \times \R}
 \ar[r]^-{\Obs^\partial}_-{\text{C--G}~\cite{CG21}}
& \ChirAlg^{E_1}_X
 \ar[d]^-{\av}_-{\text{Vol~I}}
\\
& & & \mathfrak{g}^{\mathrm{mod}}_\bullet
}
\]
The diagram commutes on objects up to chain-level
quasi-isomorphism by parts~(i)--(iii) of the theorem; the middle
two arrows compose to the identity at chain level by Vol~II
Theorem ``Bulk-to-boundary OPE consistency''; the
chain-level desuspension $[1]$ across the
$\Phi_{d,\sigma} \to \Phi_{\mathrm{hol}}$ interface is the holomorphic
factorisation map of Vol~II Theorem ``Bar differential =
holomorphic factorisation,'' part~(i); and the
$\Sigma_n$-projection on the right is identity at $n=2$ on the
class~G/L/M boundary-linear image by Vol~I
Proposition~\textup{\ref{prop:hdm-averaging-compatibility}}. The
content of Theorem~\textup{\ref{thm:hdm-triangle-composite}} is
that the four bridges \emph{do} compose at chain level on
$\CY_d^{\,\Phi\text{-BL}}$, with explicit class-by-class image
identification, and that the obstruction to extending the composite
beyond $\CY_d^{\,\Phi\text{-BL}}$ is the boundary-linearity
hypothesis~\textup{\ref{H4-triangle}}, which is required by the
source of $\Phi_{\mathrm{hol}}$ and is not implied by either
$\Phi_{d,\sigma}$ or $\av$.
\end{remark}

\section{Synthesis}
\label{sec:hdm-supplement}

\begin{remark}[Holographic datum master at the K3 base]
\label{rem:hdm-1}
The holographic datum master diagram, when evaluated at the K3
bialgebra on the boundary-linear BKM row, produces the pinned
six-slot datum
\[
\mathcal H(\mathbf H_{\Delta_5}) \;=\; \bigl(\mathbf H_{\Delta_5},\;
\mathbf H_{\Delta_5}^{\!(!)},\;
\cZ^{\mathrm{der}}_{\mathrm{ch}}(\mathbf H_{\Delta_5})
\simeq \mathcal{C}^\bullet_{\mathrm{ch}}(\mathbf H_{\Delta_5},
\mathbf H_{\Delta_5}),\; R_{\mathrm{Sieg,dyn}},\;
\Theta_{\mathbf H_{\Delta_5}},\;
\nabla^{\mathrm{KZB}}_{\bar{\mathcal A}_2}\bigr).
\]
The ordered bar complex $B^{\mathrm{ord}}_{\mathrm{ch}}
(\mathbf H_{\Delta_5})$ remains the chain model used to construct
$\Theta_{\mathbf H_{\Delta_5}}$ and its binary residue
$R_{\mathrm{Sieg,dyn}}$; it is not the bulk slot. The bulk slot is
the chiral derived centre, as in Theorem~\ref{thm:thqg-swiss-cheese}
and the Vol~III convention
$\cZ^{\mathrm{der}}_{\mathrm{ch}}(A_\cC)$ for the CY closed-string
side. This does not create a curve-level degree-$3$ Hochschild class:
Theorem~H keeps
$\ChirHoch^3(\mathbf H_{\Delta_5},\mathbf H_{\Delta_5})=0$ on the
ambient curve; the non-zero degree-$3$ representatives belong only to
the product/CY$_3$ ambient extension.
The six entries jointly satisfy the master Maurer--Cartan identity
\[
D\cdot \Theta_{\mathbf H_{\Delta_5}} + \tfrac{1}{2}[\Theta_{\mathbf H_{\Delta_5}}, \Theta_{\mathbf H_{\Delta_5}}] \;=\; 0
\quad\text{on}\quad \mathrm{Ran}(E^{\mathrm{nod,sm}}_{24})\times\bar{\mathcal A}_2,
\]
with $D$ the bi-based Ran-space de~Rham differential. On this row the
proved cross-volume content is the compatibility of the binary
residue, the MC class, and the Borcherds weight
$\kappa_{\mathrm{BKM}}(\Delta_5)=5$ through the bi-based Ran-space
datum. Full quasi-equivalence reconstruction of the underlying theory
from $\mathcal H(\mathbf H_{\Delta_5})$ remains the global-triangle
and $\Phi_{d,\sigma}$-functoriality obligation stated in
Theorem~\ref{thm:hdm-triangle-composite}.
\end{remark}
